% Appendix G: Part II Physical Dynamics Proofs (Chapters 10--16)
\label{sec:appendix-dynamics}

% --- Proofs from ch10_physical_dynamics.tex ---

\begin{proof}[Proof of Proposition~\ref{prop:temporal-realization-law-identity}]
\label{proof:prop:temporal-realization-law-identity}
Corollary~\ref{cor:temporal-interface-precedes-physical-time} supplies the
derived temporal interface \((D=3,[\Theta])\) and shows that the earlier
structural chapters together with the continuum bridge stop there.
Theorem~\ref{thm:temporal-realization-forces-selected-cut-forcing-package}
identifies the least-motion faithful selected cut with the realized horizon
cut, and Theorem~\ref{thm:canonical-selected-schur-realization} then places
exact visibility, exact Schur return, and the same selected-cut forcing data
on that datum. Definition~\ref{def:selected-visible-effective-law} is exactly
the operator formula
\[
\mathcal L_{\mathrm{phys}}(z)
= P_\ast A_{\mathrm{coh}} P_\ast + z\,\CCSigmaP{P_\ast}{z}
\]
when read on that selected cut. The notation
\[
t_{\mathrm{phys}} := \tau_{P_\ast}(D=3,[\Theta])
\]
then names the realized temporal variable on that same cut. The two displayed
definitions therefore record one law on one selected datum: realized time and
realized visible dynamics.
\end{proof}

\begin{proof}[Proof of Theorem~\ref{thm:exact-visible-selected-cut-package-forces-self-sufficient-realization}]
\label{proof:thm:exact-visible-selected-cut-package-forces-self-sufficient-realization}
Item \emph{(i)} of
Definition~\ref{def:exact-visible-selected-cut-package} is exactly the
exact-visibility clause of
Definition~\ref{def:self-sufficient-realization-principle}. Item \emph{(ii)}
states that the selected visible law is read on that same selected cut by the
internal effective-law calculus of the coherent bundle itself, so the law is
intrinsically closed there rather than supplemented by external datum. Item
\emph{(iii)} says that any apparent constitutive contribution is already part of
that same internally induced hidden-complement correction and therefore does
not require exogenous constitutive completion. Because that correction comes
from elimination of the hidden complement of the same coherent bundle, such an
apparent constitutive term is exactly returned residual from eliminated
sectors of that bundle. Hence all four clauses of
Definition~\ref{def:self-sufficient-realization-principle} hold.
\end{proof}

\begin{proof}[Proof of Proposition~\ref{prop:scoped-temporal-completion-unique}]
\label{proof:prop:scoped-temporal-completion-unique}
Because both data satisfy the temporal realization principle, item
\emph{(i)} of Definition~\ref{def:physical-realization-principle} applies to
each of them. That clause states that the physically realized cut is unique
among faithful, closure-admissible, promotion-stable characteristic observer
data up to horizon-preserving equivalence. Hence the two selected cuts are
horizon-preserving equivalent on the realized temporal class, proving
\emph{(i)}.

The same clause also states that, on the physically realized cut, the
observable dynamics are exactly the corresponding selected visible effective
law. Since the two cuts are equivalent realizations of the same realized
temporal class, the induced observable dynamics agree up to that equivalence.
This proves \emph{(ii)}.

Finally, item \emph{(ii)} of
Definition~\ref{def:physical-realization-principle} identifies every
admissible visible source as exchange with another visible sector or with the
hidden complement of the same coherent bundle, identifies all observable
forcing as returned residual, and identifies the irreversible content of the
visible dynamics as unrecovered residual. Item \emph{(iii)} together with
Definition~\ref{def:self-sufficient-realization-principle} identifies any
apparent constitutive term as internally induced hidden-sector elimination and
forbids exogenous constitutive completion. These readings therefore coincide on
both admissible completions, proving \emph{(iii)}. The final sentence is just
the summary of \emph{(i)}--\emph{(iii)}.
\end{proof}

\begin{proof}[Proof of Theorem~\ref{thm:coherent-law-exact-elimination-forces-exact-visible-selected-cut-package}]
\label{proof:thm:coherent-law-exact-elimination-forces-exact-visible-selected-cut-package}
Item \emph{(i)} of
Definition~\ref{def:coherent-law-exact-elimination-package} is exactly item
\emph{(i)} of Definition~\ref{def:exact-visible-selected-cut-package}. By
Definition~\ref{def:selected-visible-effective-law}, the operator
\(\mathcal L_{\mathrm{vis},h}(z)\) is already the exact effective operator
obtained from the coherent operator \(A_{\mathrm{coh}}\) by eliminating the
hidden complement of the same selected cut. Hence the selected observed law is
read on that cut by the internal effective-law calculus of the same coherent
bundle, which is item \emph{(ii)} of
Definition~\ref{def:exact-visible-selected-cut-package}. Next, item
\emph{(ii)} identifies the selected return field with the residual return field
of the same coherent operator and selected projector, and item \emph{(iii)}
states that the only hidden-side contribution visible on the selected cut is
the exact Schur correction. By
Proposition~\ref{prop:return-residual-is-schur}, that correction is exactly the
returned residual forcing from the eliminated sector. Therefore any apparent
constitutive contribution on the selected cut is already contained in the
internally induced hidden-complement correction and introduces no external
datum. This is exactly item \emph{(iii)} of
Definition~\ref{def:exact-visible-selected-cut-package}.
\end{proof}

\begin{proof}[Proof of Theorem~\ref{thm:selected-visible-effective-law-forces-no-auxiliary-shadow-elimination}]
\label{proof:thm:selected-visible-effective-law-forces-no-auxiliary-shadow-elimination}
Let
\[
F(u):=(I-zA_{\mathrm{coh}})u-g
\]
on the ambient coherent bundle, with \(P:=P_h\) and \(Q:=I-P_h\). By
Definition~\ref{def:nonlinear-pq-split}, the coherent law splits on that same
bundle into the visible equation \(F_P(p,q)=0\) and the hidden-complement
equation \(F_Q(p,q)=0\). This is clause \emph{(i)} of
Definition~\ref{def:no-auxiliary-internal-shadow-elimination-package}.

Because \(\mathcal L_{\mathrm{vis},h}(z)\) is defined by the exact Schur
correction \(\CCSigmaP{P_h}{z}\), the shadow resolvent
\[
(I-zQA_{\mathrm{coh}}Q)^{-1}
\]
exists on the selected cut. Corollary~\ref{cor:linear-shadow-specialization}
therefore furnishes the unique shadow solve
\[
\Psi(p)
=
(I-zQA_{\mathrm{coh}}Q)^{-1}\bigl(Qg+z\,QA_{\mathrm{coh}}P_hp\bigr),
\]
and Theorem~\ref{thm:nonlinear-shadow-elimination} identifies the full
coherent law on the selected cut with the reduced observed equation together
with the reconstruction rule \(q=\Psi(p)\). This is clause \emph{(ii)}.

In the selected-law specialization \(Qg=0\),
Corollary~\ref{cor:linear-shadow-specialization} gives
\[
\bigl(I-z\,\mathcal L_{\mathrm{vis},h}(z)\bigr)p=P_hg.
\]
By Definition~\ref{def:selected-visible-effective-law},
\(\mathcal L_{\mathrm{vis},h}(z)\) is exactly the effective operator obtained
from \(A_{\mathrm{coh}}\) on the selected cut, so the reduced observed
equation is exactly the selected visible effective law.
The same corollary identifies the returned visible forcing through the same
shadow solve, and Proposition~\ref{prop:return-residual-is-schur} identifies
that returned forcing with the residual return / Schur correction. This is
clause \emph{(iii)}.

Finally, the displayed formula for \(\Psi(p)\) depends only on the visible
datum \(p\), the same coherent operator \(A_{\mathrm{coh}}\), and the same
selected projector \(P_h\). It introduces no hidden labels, quotient choice,
gauge slice, or exogenous constitutive datum. This is exactly clause
\emph{(iv)}. Hence the no-auxiliary internal shadow elimination package holds
on the selected cut.
\end{proof}

\begin{proof}[Proof of Theorem~\ref{thm:no-auxiliary-internal-shadow-elimination-forces-coherent-law-exact-elimination}]
\label{proof:thm:no-auxiliary-internal-shadow-elimination-forces-coherent-law-exact-elimination}
By clause \emph{(ii)} and
Theorem~\ref{thm:nonlinear-shadow-elimination}, the full coherent law on the
selected datum is exactly equivalent to the reduced observed equation together
with the reconstruction rule given by the internal shadow solve. Clause
\emph{(iii)} identifies the linear selected-law specialization of that reduced
equation with the selected visible effective law
\(\mathcal L_{\mathrm{vis},h}(z)\), so clause \emph{(i)} of
Definition~\ref{def:coherent-law-exact-elimination-package} holds.

The same clause \emph{(iii)} identifies the returned visible forcing of the
selected shadow solve with the residual return field and its Schur correction,
using Corollary~\ref{cor:linear-shadow-specialization} together with
Proposition~\ref{prop:return-residual-is-schur}. Therefore the selected return
field on that cut is exactly the residual return field of the same coherent
operator and selected projector, which is clause \emph{(ii)} of
Definition~\ref{def:coherent-law-exact-elimination-package}.

Finally, clause \emph{(iv)} is exactly the no-extra-state condition: the shadow
solve introduces no hidden labels, quotient choice, gauge slice, or exogenous
constitutive datum. Since every hidden-side visible contribution is already
read through that exact internal elimination on the same coherent bundle, the
only visible hidden-side contribution is the Schur correction / returned
residual identified above. This is clause \emph{(iii)} of
Definition~\ref{def:coherent-law-exact-elimination-package}. Hence the
coherent-law exact elimination package holds on the selected cut.
\end{proof}

\begin{proof}[Proof of Theorem~\ref{thm:temporal-realization-forces-causal-propagation}]
\label{proof:thm:temporal-realization-forces-causal-propagation}
By Definition~\ref{def:physical-realization-principle}, the selected visible
law persists as an ordered visible dynamics along the chosen horizon family;
this ordered persistence is exactly the physical time read on the realized cut.
The same definition further states that all observable forcing is returned
residual from hidden or interface propagation in the same coherent bundle.
Finally, Definition~\ref{def:self-sufficient-realization-principle} excludes
any exogenous constitutive or source completion on that cut. Therefore visible
influence is admissible only when it is carried by the realized causal order
and re-enters through returned residual on that same bundle. No primitive
acausal visible action remains.
\end{proof}

\begin{proof}[Proof of Theorem~\ref{thm:physical-symmetries-are-temporal-stabilizers}]
\label{proof:thm:physical-symmetries-are-temporal-stabilizers}
By observer covariance, admissible observer data in the same realized temporal
class yield horizon-preserving equivalent visible laws. Therefore any
admissible transformation preserving that realized class together with the
selected visible law and its residual-return bookkeeping acts trivially on the
physical content of the visible dynamics, hence is a physical symmetry. For the
same reason, any physical symmetry must preserve the realized temporal class
and the selected visible law, since otherwise it would change the physically
realized dynamics. Thus the physical symmetry group is exactly the stabilizer
of the temporally realized structure.
\end{proof}

\begin{proof}[Proof of Theorem~\ref{thm:noether-bridge-on-temporally-realized-reversible-sectors}]
\label{proof:thm:noether-bridge-on-temporally-realized-reversible-sectors}
Let \(\phi_\varepsilon\) be the history obtained by acting on the realized
history by the one-parameter temporal stabilizer. Because the stabilizer
preserves the action, one has
\[
\frac{d}{d\varepsilon}\Big|_{\varepsilon=0} S[\phi_\varepsilon] = 0.
\]
Because the realized history is stationary, the bulk first variation already
vanishes on shell. The remaining first-variation term is therefore a pure
carrier-side boundary / divergence term, which defines the associated current
\(J\). Since the full first variation is zero, that boundary / divergence term
must vanish, giving \(\partial J = 0\). Thus continuous temporal stabilizers
produce conserved currents and charges on reversible realized sectors.
\end{proof}

\begin{proof}[Proof of Theorem~\ref{thm:coupled-exchange-balance}]
\label{proof:thm:coupled-exchange-balance}
By Definition~\ref{def:exchange-balanced-coupled-family},
\[
\sum_{\alpha\in A} S_\alpha
=
\sum_{\alpha\in A}\sum_{\beta\in A\setminus\{\alpha\}}\partial J_{\alpha\beta}.
\]
Because \(\partial\) is linear, the double sum groups into unordered pairs
\(\{\alpha,\beta\}\). For each such pair,
\[
\partial J_{\alpha\beta}+\partial J_{\beta\alpha}
=
\partial\bigl(J_{\alpha\beta}+J_{\beta\alpha}\bigr)
=
0
\]
by the antisymmetry relation \(J_{\alpha\beta}=-J_{\beta\alpha}\). Hence every
pairwise exchange contribution cancels, proving
\(\sum_{\alpha\in A}S_\alpha=0\). If \(A=\{X,Y\}\), then
\[
S_X=\partial J_{XY}
\qquad\text{and}\qquad
S_Y=\partial J_{YX}=-\partial J_{XY},
\]
so \(S_X=-S_Y\).
\end{proof}

\begin{proof}[Proof of Theorem~\ref{thm:nonempty-schur-induced-mode-classes-have-singleton-reduced-selection}]
\label{proof:thm:nonempty-schur-induced-mode-classes-have-singleton-reduced-selection}
Because the canonical selected Schur bridge supplies an admissible realized
class \(\mathfrak R_{\mathrm{sel}}(h)\), the selection interface of
Theorem~\ref{thm:selection-interface} applies to every realization in that
class. The subclass
\(\mathfrak R^{\mathrm{Schur}}_{\mathfrak m}(h)\subseteq \mathfrak R_{\mathrm{sel}}(h)\)
therefore inherits the same finite reduced selection mechanism, the same
well-posed admissibility constraint, and the same observer-motion objective on
the selected cut.

By Definition~\ref{def:schur-induced-mode-candidate-class}, every realization
in \(\mathfrak R^{\mathrm{Schur}}_{\mathfrak m}(h)\) carries a visible
subcarrier \(\mathcal Y_R\subseteq \mathcal X_h\) that is faithful,
reconstruction-compatible, closure-admissible, stable under admissible
promotion, and irreducible for the compatibility signature and admissible
source profile of the mode \(\mathfrak m\). Since the class is nonempty by
hypothesis, these transported subcarriers define a nonempty class of
mode-compatible candidates on the selected cut. This proves \emph{(i)}.

For \emph{(ii)}, the canonical selected Schur bridge includes
horizon-preserving equivalence across the whole admissible class
\(\mathfrak R_{\mathrm{sel}}(h)\). Hence any two realizations
\(R,R'\in \mathfrak R^{\mathrm{Schur}}_{\mathfrak m}(h)\) are related by a
horizon-preserving equivalence intertwining the selected observed law, the
admissible reduced operator data, and the normalization package. By
construction of
\(\mathfrak R^{\mathrm{Schur}}_{\mathfrak m}(h)\), the defining properties of
the transported subcarriers \(\mathcal Y_R\) and \(\mathcal Y_{R'}\) are stated
entirely in terms of bridge-invariant data: faithfulness, admissibility,
promotion stability, preservation of the mode signature, and irreducibility.
Therefore the reduced selection data carried by \(R\) and \(R'\) represent the
same admissible reduced datum modulo horizon-preserving equivalence. The
reduced selection data on
\(\mathfrak R^{\mathrm{Schur}}_{\mathfrak m}(h)\) are thus singleton modulo
horizon-preserving equivalence.

Item \emph{(iii)} now follows from
Theorem~\ref{thm:mode-compatible-selection-forces-least-motion-realization},
applied to the mode-compatible candidate class produced in \emph{(i)} together
with the singleton-collapse statement of \emph{(ii)}.
\end{proof}

\begin{proof}[Proof of Theorem~\ref{thm:intrinsic-selected-channel-package-forces-three-schur-mode-classes}]
\label{proof:thm:intrinsic-selected-channel-package-forces-three-schur-mode-classes}
Let \(R_\ast\in\mathfrak R_{\mathrm{sel}}(h)\) be a realization whose selected
datum is the least-motion characteristic observer datum fixed by the canonical
selected Schur bridge. By
Corollary~\ref{cor:intrinsic-package-implies-governing-datum-rigid},
\(\mathfrak R_{\mathrm{sel}}(h)\) is governing-datum rigid on the relevant
selected channels. Therefore
Theorem~\ref{cor:realized-local-internal-datum-forcing} applies to every
relevant selected channel of \(R_\ast\): on such a channel, the selected law
carries
\begin{enumerate}[label=(\roman*), leftmargin=2em, itemsep=0.2em]
\item a forced phase-carrier class modulo phase-equivalence,
\item a unique scalar quadratic recursion datum
      \(B_{h,\rho}^\ast B_{h,\rho}=\kappa_{h,\rho}I\) together with the forced
      reversible recursion law
      \(x_{n+1}=B_{h,\rho}x_n-x_{n-1}\),
\item a scalar closure class.
\end{enumerate}

Let \(\mathcal Y_\rho\subseteq \mathcal X_h\) be any relevant selected channel
fiber of \(R_\ast\), transported to the common selected bundle by the
horizon-preserving equivalence built into the canonical selected Schur bridge.
Because \(R_\ast\) lies in the bridge's admissible realized class,
\(\mathcal Y_\rho\) is faithful, reconstruction-compatible,
closure-admissible, and stable under admissible promotion on the selected cut.
The phase datum together with the forced reversible recursion and scalar
closure data therefore make \(\mathcal Y_\rho\) a witness for
local-state persistence. The same forced reversible recursion together with the
quadratic transport datum and the no-preferred-direction hypothesis already
present in
Corollary~\ref{cor:intrinsic-package-implies-governing-datum-rigid} make
\(\mathcal Y_\rho\) a witness for propagative persistence.

For comparative persistence, hypothesis \emph{(ii)} and
Corollary~\ref{cor:intrinsic-package-yields-gauge-compatible-local-datum}
apply on the same selected channels. Hence the selected law carries a unique
gauge-compatible local internal datum: in particular, the selected channel
transport system is locally gauge-covariant on the same visible fibers. The
comparison class determined by the transported channel system on
\(\mathcal Y_\rho\) is therefore a witness for comparative persistence on the
selected cut.

Thus \(R_\ast\) belongs to each of the three Schur-induced mode candidate
classes above. Hence those three classes are nonempty.
\end{proof}

\begin{proof}[Proof of Theorem~\ref{thm:schur-lifted-symmetric-square-response-forces-geometric-witness}]
\label{proof:thm:schur-lifted-symmetric-square-response-forces-geometric-witness}
Let \(R_\ast\), \(F_h\), and
\(\mathcal X_h^{\mathrm{sym}}\subseteq \mathcal X_h\) be as in
Definition~\ref{def:schur-lifted-symmetric-square-response-package}. By items
\emph{(ii)}--\emph{(iii)} of that definition,
\(\mathcal X_h^{\mathrm{sym}}\) is a faithful,
reconstruction-compatible, closure-admissible, admissible-promotion-stable
visible symmetric-response subcarrier on the selected bundle.

By item \emph{(iv)}, the selected visible effective law together with the
canonical return field induces a divergence-compatible
second-order response on \(\mathcal X_h^{\mathrm{sym}}\). Because the family
satisfies a canonical selected Schur bridge,
Theorem~\ref{thm:canonical-selected-schur-realization}(d)--(e) identifies that
return field with exact Schur residual on the same selected cut. Thus the
response on \(\mathcal X_h^{\mathrm{sym}}\) is genuinely Schur-induced rather
than an independently inserted constitutive term.

By item \emph{(v)} and Theorem~\ref{thm:coupled-exchange-balance}, visible
source on \(\mathcal X_h^{\mathrm{sym}}\) is admissible only as balanced
exchange with preceding selected carriers on the same selected bundle, encoded
here by a symmetric exchange tensor. This is exactly the causal-structural
source profile. Finally, item \emph{(vi)} is precisely the irreducibility
clause required of the symmetric-response witness.

Therefore \(\mathcal X_h^{\mathrm{sym}}\) satisfies all hypotheses of
Theorem~\ref{thm:macroscopic-witness-packages-force-two-schur-mode-classes}.
That theorem yields the nonemptiness of the causal-structural Schur-induced
mode candidate class.
\end{proof}

\begin{proof}[Proof of Corollary~\ref{cor:orthogonal-schur-lifted-symmetric-response-closes-mode-gap}]
\label{proof:cor:orthogonal-schur-lifted-symmetric-response-closes-mode-gap}
By
Theorem~\ref{thm:intrinsic-selected-channel-package-forces-three-schur-mode-classes},
the local-state, propagative, and comparative Schur-induced mode candidate
classes are nonempty. By
Theorem~\ref{thm:schur-lifted-symmetric-square-response-forces-geometric-witness},
the causal-structural Schur-induced mode candidate class is also nonempty.
Hence all four intrinsic Schur-induced mode candidate classes are nonempty.
Corollary~
\ref{cor:five-nonempty-schur-induced-mode-classes-force-intrinsic-selected-bundle-existence}
therefore yields intrinsic selected-bundle existence.

If the remaining hypotheses of
Theorem~\ref{thm:single-witness-intrinsic-flagship-forcing} are also satisfied,
then the Schur-lifted symmetric-square response package supplies exactly the
witness input required there via
Theorem~\ref{thm:schur-lifted-symmetric-square-response-forces-geometric-witness}.
Thus that theorem applies and forces the four intrinsic flagship law classes
on the selected bundle.
\end{proof}

\begin{proof}[Proof of Theorem~\ref{thm:first-three-intrinsic-law-classes-already-forced}]
\label{proof:thm:first-three-intrinsic-law-classes-already-forced}
By
Theorem~\ref{thm:intrinsic-selected-channel-package-forces-three-schur-mode-classes},
the local-state, propagative, and comparative Schur-induced mode candidate
classes on the selected bundle are nonempty. By
Theorem~\ref{thm:nonempty-schur-induced-mode-classes-have-singleton-reduced-selection},
each of those three classes already has singleton reduced selection data
modulo horizon-preserving equivalence and therefore yields a unique
least-motion visible realization of the corresponding persistence mode on the
selected cut.

Because a canonical selected Schur bridge fixes one selected observed bundle
\(\mathcal X_h=P_h\cH^\bullet\), these three unique least-motion visible
realizations live on one and the same selected bundle. By the defining mode
signatures in Definition~\ref{def:five-persistence-modes} and the minimality
clause in Definition~\ref{def:schur-induced-mode-candidate-class}, the
local-state realization is exactly the canonical minimal phase-generating
carrier, the propagative realization is exactly the canonical minimal
reversible second-order carrier, and the comparative realization is exactly
the canonical minimal exactness-compatible one-form carrier. This proves item
\emph{(a)}.

For item~\emph{(b)}, fix one of these three carriers. The singleton reduced
selection data just obtained are exactly the selected-cut forcing clause used
inside the proof of Theorem~\ref{thm:common-endpoint-rigidity-reduction}.
The scalarization surface required by that theorem is exactly
Corollary~\ref{cor:realized-self-adjoint-scalarization} on the irreducible
symmetry channels of the phase, wave, and gauge carriers. Because a canonical
selected Schur bridge fixes one exact selected visible law on the selected
bundle, Definition~\ref{def:selected-visible-effective-law} identifies the
selected endpoint closure class on each of these carriers as the
characteristic-boundary reduction of that same exact visible law. Therefore
Theorem~\ref{thm:common-endpoint-rigidity-reduction} applies on each of the
three carriers and yields the stated uniqueness modulo continuum equivalence.

Item~\emph{(c)} is then exactly the specialization of
Remark~\ref{rem:remaining-intrinsic-endpoint-obligations} to these three
already-forced carriers: once the carrier exists and its endpoint class is
already rigid modulo continuum equivalence, only the natural representative
identification remains.
\end{proof}

\begin{proof}[Proof of Theorem~\ref{thm:three-primitive-channels-already-forced}]
\label{proof:thm:three-primitive-channels-already-forced}
Item~\emph{(i)} is exactly item~\emph{(a)} of
Theorem~\ref{thm:first-three-intrinsic-law-classes-already-forced} together
with Corollary~\ref{cor:intrinsic-symmetry-forces-123}, which identifies the
underlying primitive channel profile as \(1+2+3\). Item~\emph{(ii)} is the
same identification read directly from
Theorem~\ref{thm:first-three-intrinsic-law-classes-already-forced}: local-state
persistence is realized by the phase carrier, propagative persistence by the
wave carrier, and comparative persistence by the gauge carrier. For item
\emph{(iii)}:
Corollary~\ref{cor:flagship-family-exhausts-five-persistence-modes} proves
that within the intrinsic flagship family there is no fourth primitive
persistence channel and no additional intrinsic lift. In particular there is
no further primitive channel beyond the three already listed.
\end{proof}

\begin{proof}[Proof of Corollary~\ref{cor:geometry-only-remaining-intrinsic-existence-lane}]
\label{proof:cor:geometry-only-remaining-intrinsic-existence-lane}
Item~\emph{(i)} is exactly
Theorem~\ref{thm:first-three-intrinsic-law-classes-already-forced}. Item
\emph{(ii)} is
Theorem~\ref{thm:primitive-three-symmetric-bilinear-lift-forces-schur-tensor-lift}
and
Theorem~\ref{thm:schur-lifted-symmetric-square-response-forces-geometric-witness}
combined with
Remark~\ref{rem:remaining-intrinsic-endpoint-obligations}: before that
bilinear/symmetric-square lift package is derived, the geometric lane is the
only intrinsic lane whose mode-existence theorem is still open; once it is
derived, the carrier-level step is collapse to the Einstein
representative.
\end{proof}

\begin{proof}[Proof of Theorem~\ref{thm:lane-by-lane-intrinsic-closure-frontier}]
\label{proof:thm:lane-by-lane-intrinsic-closure-frontier}
Items~\emph{(i)}--\emph{(iii)} are the lane-by-lane specialization of
Theorem~\ref{thm:first-three-intrinsic-law-classes-already-forced}: phase,
wave, and gauge already arise on the selected bundle as the canonical minimal
compatible realizations of the local-state, propagative, and comparative
persistence modes, their endpoint closure classes are rigid modulo continuum
equivalence, and therefore the only remaining burden on those three lanes is
the carrier-level representative collapse. The concrete representatives named
in items~\emph{(i)}--\emph{(iii)} are exactly the ones isolated in
Corollary~\ref{cor:recognizable-schrodinger-identity},
Corollary~\ref{cor:recognizable-wave-kg-identity}, and
Corollary~\ref{cor:recognizable-maxwell-identity}; the corresponding collapse
theorems are
Theorems~\ref{thm:phase-local-endpoint-collapse},
\ref{thm:wave-local-endpoint-collapse}, and
\ref{thm:gauge-local-endpoint-collapse}.

For item~\emph{(iv)}, the existence side is already closed by
Theorem~\ref{thm:geometry-unique-lift-of-primitive-three-channel}, which
forces the geometric law class on the selected bundle as the unique intrinsic
lift of the primitive \(3\)-channel. The local carrier-level representative
collapse is then discharged by
Theorem~\ref{thm:geometric-local-endpoint-collapse}, and
Corollary~\ref{cor:recognizable-einstein-identity} records that closed law in
standard Einstein notation.
\end{proof}

\begin{proof}[Proof of Theorem~\ref{thm:macroscopic-witness-packages-force-two-schur-mode-classes}]
\label{proof:thm:macroscopic-witness-packages-force-two-schur-mode-classes}
By hypothesis, the symmetric-response subcarrier
\(\mathcal X_h^{\mathrm{geom}}\subseteq \mathcal X_h\) is faithful,
reconstruction-compatible, closure-admissible, stable under admissible
promotion, and irreducible for the divergence-compatible second-order response
signature. By the temporal realization principle, visible source on that
subcarrier is admissible only as exchange or returned residual on the same
selected bundle, so \(\mathcal X_h^{\mathrm{geom}}\) has exactly the source
profile of causal-structural persistence. Therefore the realization carrying
\(\mathcal X_h^{\mathrm{geom}}\) belongs to the Schur-induced mode candidate
class for causal-structural persistence, proving that class is nonempty.
\end{proof}

\begin{proof}[Proof of Corollary~\ref{cor:five-nonempty-schur-induced-mode-classes-force-intrinsic-selected-bundle-existence}]
\label{proof:cor:five-nonempty-schur-induced-mode-classes-force-intrinsic-selected-bundle-existence}
By
Theorem~\ref{thm:nonempty-schur-induced-mode-classes-have-singleton-reduced-selection},
each of the four nonempty Schur-induced mode classes determines a
mode-compatible candidate class with singleton reduced selection data, and
hence yields a unique least-motion visible realization of the corresponding
primitive channel or geometric lift. Therefore the hypotheses of
Corollary~\ref{cor:mode-wise-candidate-forcing-yields-compatibility-surface}
are satisfied, so the inherited compatibility surface holds on the selected
cut. By
Corollary~\ref{cor:mode-wise-candidate-forcing-yields-mode-complete-realization},
the selected cut satisfies intrinsic-mode-complete visible realization.
By
Theorem~\ref{thm:four-nonempty-schur-induced-mode-classes-force-canonical-realization},
the same selected bundle also satisfies canonical realization of persistence
modes. Finally,
Theorem~\ref{thm:mode-complete-realization-forces-intrinsic-selected-bundle}
applies and yields intrinsic selected-bundle existence.
\end{proof}

\begin{proof}[Proof of Theorem~\ref{thm:four-nonempty-schur-induced-mode-classes-force-canonical-realization}]
\label{proof:thm:four-nonempty-schur-induced-mode-classes-force-canonical-realization}
For each of the four persistence roles of
Definition~\ref{def:five-persistence-modes}, the corresponding Schur-induced
mode candidate class is nonempty by hypothesis. Applying
Theorem~\ref{thm:nonempty-schur-induced-mode-classes-have-singleton-reduced-selection}
to each role yields a unique least-motion visible realization on the selected
cut for that role.

By Definition~\ref{def:schur-induced-mode-candidate-class}, each such
realization is minimal among visible subcarriers carrying the compatibility
signature and admissible source profile of its role. For the local-state role
this is the minimal phase-generating realization; for the propagative role it
is the minimal reversible second-order realization; for the comparative role
it is the minimal exactness-compatible one-form realization; and for the
causal-structural role it is the minimal divergence-compatible
symmetric-response lift of the primitive \(3\)-channel. These are exactly the
four clauses of Definition~\ref{def:canonical-sector-generation}. Therefore
the selected bundle satisfies canonical realization of persistence modes.
\end{proof}

\begin{proof}[Proof of Theorem~\ref{thm:geometric-witness-is-exact-remaining-intrinsic-obstruction}]
\label{proof:thm:geometric-witness-is-exact-remaining-intrinsic-obstruction}
The implication \emph{(i)} \(\Rightarrow\) \emph{(ii)} is exactly
Theorem~\ref{thm:macroscopic-witness-packages-force-two-schur-mode-classes}.

For \emph{(ii)} \(\Rightarrow\) \emph{(i)}, let
\(R\in \mathfrak R^{\mathrm{Schur}}_{\mathfrak m}(h)\) for the
causal-structural mode \(\mathfrak m\). By
Definition~\ref{def:schur-induced-mode-candidate-class}, after transporting the
selected datum of \(R\) to the common selected bundle there is a visible
subcarrier \(\mathcal Y_R\subseteq \mathcal X_h\) that is faithful,
reconstruction-compatible, closure-admissible, and stable under admissible
promotion. Because \(\mathfrak m\) is causal-structural persistence,
Definition~\ref{def:five-persistence-modes} identifies its compatibility
signature as divergence-compatible response on the time-realizing cut with
visible source carried by a preceding exchange tensor. Therefore the
restriction of the selected visible effective law to \(\mathcal Y_R\) carries a
divergence-compatible second-order response signature on a visible symmetric
response subcarrier. The irreducibility clause in
Definition~\ref{def:schur-induced-mode-candidate-class} says that no proper
visible subcarrier of \(\mathcal Y_R\) still supports that same mode signature.
Hence \(\mathcal Y_R\) satisfies exactly the symmetric-response witness
hypothesis of
Theorem~\ref{thm:macroscopic-witness-packages-force-two-schur-mode-classes}.

By Theorem~\ref{thm:intrinsic-selected-channel-package-forces-three-schur-mode-classes},
the local-state, propagative, and comparative Schur-induced mode classes are
already nonempty. The equivalence just proved gives nonemptiness of the
causal-structural class. Therefore
Corollary~\ref{cor:intrinsic-package-plus-macroscopic-witnesses-force-five-schur-nonempty}
and
Theorem~\ref{thm:four-nonempty-schur-induced-mode-classes-force-canonical-realization}
yield canonical realization of persistence modes. The same nonemptiness
package together with
Corollary~\ref{cor:five-nonempty-schur-induced-mode-classes-force-intrinsic-selected-bundle-existence}
then yields intrinsic selected-bundle existence.

Finally, if the hypotheses of
Theorem~\ref{thm:single-witness-intrinsic-flagship-forcing} are met, then the
first-three-modes package together with either equivalent condition above is
exactly its remaining witness input. Hence that theorem applies and forces the
four intrinsic flagship law classes on the selected bundle.
\end{proof}

\begin{proof}[Proof of Theorem~\ref{thm:single-witness-intrinsic-flagship-forcing}]
\label{proof:thm:single-witness-intrinsic-flagship-forcing}
By Corollary~
\ref{cor:intrinsic-package-plus-macroscopic-witnesses-force-five-schur-nonempty},
all four Schur-induced mode candidate classes on the selected bundle are
nonempty. Theorem~
\ref{thm:four-nonempty-schur-induced-mode-classes-force-canonical-realization}
therefore supplies canonical realization of persistence modes, and
Corollary~
\ref{cor:five-nonempty-schur-induced-mode-classes-force-intrinsic-selected-bundle-existence}
then yields item~\emph{(a)}: the phase, wave, gauge, and geometric carriers
arise directly as associated visible carriers of one physically realized
selected observed bundle on the least-motion cut.

The same canonical selected Schur bridge carries the intrinsic
\(1+2+3\) channel profile by
Theorem~\ref{thm:canonical-selected-schur-realization}(g). Together with the
derived canonical realization statement and item~\emph{(a)}, this is exactly
the intrinsic scope of Theorem~\ref{thm:scoped-intrinsic-sector-generation}.
Hence item~\emph{(b)} holds: the intrinsic flagship sectors are generated on
that one selected bundle as the canonical minimal compatible realizations of
the three primitive channels together with the geometric lift of the primitive
\(3\)-channel.

For item~\emph{(c)}, fix one of the four intrinsic carriers. Because its
Schur-induced mode class is nonempty, Theorem~
\ref{thm:nonempty-schur-induced-mode-classes-have-singleton-reduced-selection}
already supplies singleton reduced selection data modulo horizon-preserving
equivalence for that carrier's mode-compatible class. The channelwise
scalarization surface of
Theorem~\ref{thm:common-endpoint-rigidity-reduction} is then inherited from the
structural and continuum-bridge theorem surface used throughout the law-side
dynamics of Part~II: on phase, wave, and
gauge-side irreducible symmetry channels it is exactly
Corollary~\ref{cor:realized-self-adjoint-scalarization}, while on the geometric
carrier it is exactly
Corollary~\ref{cor:realized-isotropic-symmetric-response-scalarization}. The
continuum-equivalence uniqueness clause is supplied on that same selected cut
by Definition~\ref{def:selected-visible-effective-law}: each intrinsic carrier
inherits the exact carrier-side closure class of the one selected visible law
under characteristic boundary reduction. Therefore the hypotheses of
Theorem~\ref{thm:common-endpoint-rigidity-reduction} are satisfied on each of
the four intrinsic carriers, and item~\emph{(c)} follows.

Item~\emph{(d)} is now exactly
Remark~\ref{rem:remaining-intrinsic-endpoint-obligations}: once those four
intrinsic endpoint classes are already forced, the remaining burden is not
intrinsic bundle existence or sector generation, but only the carrier-specific
identification of the natural representative on each forced class.
\end{proof}

\begin{proof}[Proof of Theorem~\ref{thm:canonical-selected-schur-bridge-forces-primitive-three-bilinear-lift-package}]
\label{proof:thm:canonical-selected-schur-bridge-forces-primitive-three-bilinear-lift-package}
By hypothesis \emph{(ii)}, the local balanced channel algebra on the selected
bundle is the visible realization of the canonical six-slot carrier of
Definition~\ref{def:local-slot-carrier}, together with its intrinsic
\(1+2+3\)-channel decomposition. Let \(F_h\subseteq \mathcal X_h\) denote the
primitive \(3\)-channel fiber.

By Theorem~\ref{thm:six-slot-carrier-is-symmetric-square-of-u3}, the canonical
six-slot carrier is the symmetric square of that primitive \(3\)-channel.
Transporting that intrinsic identification to the common selected bundle by the
horizon-preserving equivalence built into the canonical selected Schur bridge,
we obtain a visible subcarrier
\[
\mathcal X_h^{\mathrm{bil}}\subseteq \mathcal X_h
\]
and a symmetric bilinear visible coupling
\[
\beta_h:F_h\times F_h\to \mathcal X_h^{\mathrm{bil}}
\]
whose image generates \(\mathcal X_h^{\mathrm{bil}}\). The same theorem gives
\(\dim F_h=3\) and \(\dim \mathcal X_h^{\mathrm{bil}}=6\). This proves items
\emph{(i)}, \emph{(ii)}, and \emph{(iv)} of
Definition~\ref{def:primitive-three-symmetric-bilinear-lift-package}.

Because \(\mathcal X_h^{\mathrm{bil}}\) is read on the same selected observed
bundle and the underlying realization lies in the canonical admissible realized
class \(\mathfrak R_{\mathrm{sel}}(h)\), the structural clauses already used in
Theorem~\ref{thm:intrinsic-selected-channel-package-forces-three-schur-mode-classes}
apply to it as well: \(\mathcal X_h^{\mathrm{bil}}\) is faithful,
reconstruction-compatible, closure-admissible, and stable under admissible
promotion on the selected cut. This is item \emph{(iii)}.

By Theorem~\ref{thm:first-three-intrinsic-law-classes-already-forced}, the
primitive \(3\)-channel is already realized on the selected bundle as the
canonical minimal one-form comparison realization of comparative persistence.
By Definition~\ref{def:five-persistence-modes}, that realization carries
exactness and local gauge covariance, with visible source admissible only as an
exchange current or returned residual. The carrier
\(\mathcal X_h^{\mathrm{bil}}\) is generated by symmetric bilinear couplings of
that comparative datum. Therefore the restriction of the selected visible
effective law to \(\mathcal X_h^{\mathrm{bil}}\) is second order in the
underlying comparison data and symmetric by construction.

Exactness on the primitive \(3\)-channel forces compatibility of the induced
bilinear response with divergence on the lifted carrier: the visible source on
each comparison factor is already an exchange current, and by
Theorem~\ref{thm:temporal-realization-forbids-primitive-internal-creation} no
primitive creation is admissible on the selected cut. Because the family
satisfies a canonical selected Schur bridge,
Theorem~\ref{thm:canonical-selected-schur-realization}(c)--(e) identifies any
remaining hidden-side contribution with returned residual, namely exact Schur
correction. Hence the induced response on
\(\mathcal X_h^{\mathrm{bil}}\) is a divergence-compatible second-order visible
response signature on the selected cut. This is item \emph{(v)}.

Finally, the image of \(\beta_h\) generates all of
\(\mathcal X_h^{\mathrm{bil}}\). A proper visible subcarrier omits some
symmetric comparison directions \(u\odot v\), and therefore cannot support the
full response signature induced by the selected law on all bilinear couplings of
the primitive \(3\)-channel. Hence no proper visible subcarrier still carries
that same divergence-compatible second-order response signature at the observed
scale. This is item \emph{(vi)}.
\end{proof}

\begin{proof}[Proof of Theorem~\ref{thm:primitive-three-symmetric-bilinear-lift-forces-schur-tensor-lift}]
\label{proof:thm:primitive-three-symmetric-bilinear-lift-forces-schur-tensor-lift}
Let \(R_\ast\), \(F_h\), \(\mathcal X_h^{\mathrm{bil}}\), and \(\beta_h\) be
as in
Definition~\ref{def:primitive-three-symmetric-bilinear-lift-package}. Because
\(\beta_h\) is symmetric bilinear, the universal property of the symmetric
square yields a unique linear map
\[
\overline{\beta}_h:\mathrm{Sym}^2(F_h)\to \mathcal X_h^{\mathrm{bil}}
\]
whose value on the class of \(u\otimes v\) is \(\beta_h(u,v)\). Since the image
of \(\beta_h\) generates \(\mathcal X_h^{\mathrm{bil}}\) by hypothesis,
\(\overline{\beta}_h\) is surjective. Because \(\dim F_h=3\), standard
finite-dimensional linear algebra gives
\[
\dim \mathrm{Sym}^2(F_h)=\frac{3(3+1)}2=6.
\]
Hypothesis~\emph{(iii)}(iv) gives \(\dim \mathcal X_h^{\mathrm{bil}}=6\). A
surjective linear map between finite-dimensional vector spaces of the same
dimension is an isomorphism. Therefore \(\overline{\beta}_h\) identifies
\(\mathcal X_h^{\mathrm{bil}}\) canonically with \(\mathrm{Sym}^2(F_h)\).
This proves item~\emph{(a)}.

Now compare the hypotheses of
Definition~\ref{def:primitive-three-symmetric-bilinear-lift-package} with
those of
Theorem~\ref{thm:three-channel-schur-tensor-lift-forces-geometric-lane}. Item
\emph{(i)} of the present theorem gives the first-three-lanes package. Item
\emph{(ii)} gives the intrinsic \(1+2+3\) decomposition. Item~\emph{(a)} just
proved supplies the symmetric-square identification required there. The
faithfulness, reconstruction compatibility, closure admissibility, and
promotion stability clauses are exactly hypothesis~\emph{(iv)} of that
theorem; the divergence-compatible second-order response signature is exactly
hypothesis~\emph{(v)}; and irreducibility is exactly hypothesis~\emph{(vi)}.
Hence the hypotheses of
Theorem~\ref{thm:three-channel-schur-tensor-lift-forces-geometric-lane} are
satisfied. This proves item~\emph{(b)}.

Items~\emph{(c)} and \emph{(d)} are now exactly items~\emph{(b)} and
\emph{(c)} of
Theorem~\ref{thm:three-channel-schur-tensor-lift-forces-geometric-lane}.
\end{proof}

\begin{proof}[Proof of Theorem~\ref{thm:geometric-symmetric-square-factorization}]
\label{proof:thm:geometric-symmetric-square-factorization}
Because \(\gamma_h:F_h\times F_h\to \mathcal Y_h\) is symmetric bilinear, the
universal property of the symmetric square yields a unique linear map
\[
\overline{\gamma}_h:\mathrm{Sym}^2(F_h)\to \mathcal Y_h
\]
satisfying \(\overline{\gamma}_h(u\odot v)=\gamma_h(u,v)\). This proves
item~\emph{(a)}.

Now compare the hypotheses on \(\mathcal Y_h\) with those already used in
Theorem~\ref{thm:primitive-three-symmetric-bilinear-lift-forces-schur-tensor-lift}.
The common primitive \(3\)-channel fiber \(F_h\), the faithfulness /
reconstruction / admissibility / promotion-stability clauses, the
divergence-compatible second-order response signature, and the irreducibility
of \(\mathcal Y_h\) place \(\mathcal Y_h\) in the same symmetric-bilinear lift
problem as the canonical geometric carrier. The theorem just cited identifies
the canonical lift with \(\mathrm{Sym}^2(F_h)\) and forces the geometric lane
on the selected cut. Therefore \(\overline{\gamma}_h\) lands in the same
horizon-preserving equivalence class as that canonical lift. This proves
item~\emph{(b)}.

Item~\emph{(c)} is the reformulation of items~\emph{(a)}--\emph{(b)}: every
admissible geometric lift generated from the primitive \(3\)-channel first
factors uniquely through \(\mathrm{Sym}^2(F_h)\), and the resulting image is
horizon-preserving equivalent to the canonical geometric lift already forced
on the selected bundle.
\end{proof}

\begin{proof}[Proof of Theorem~\ref{thm:three-channel-schur-tensor-lift-forces-geometric-lane}]
\label{proof:thm:three-channel-schur-tensor-lift-forces-geometric-lane}
Items~\emph{(iii)}--\emph{(vi)} now suffice to produce the package of
Definition~\ref{def:schur-lifted-symmetric-square-response-package}. Item
\emph{(iii)} gives the selected channel fiber of dimension at least two, item
\emph{(iv)} identifies the visible symmetric-square carrier on the common
selected bundle and its faithful/reconstruction-compatible/promotion-stable
status, and item \emph{(vi)} gives the irreducibility clause. It remains only
to identify the response as Schur-induced and its source profile as balanced
visible exchange.

Because the family satisfies a canonical selected Schur bridge,
Theorem~\ref{thm:canonical-selected-schur-realization}(c)--(e) applies on the
same selected cut: the selected visible effective law is exactly visible and
intrinsically closed there, no exogenous constitutive completion is admitted,
the canonical return field is exact Schur residual, and the selected visible
law is visible closure plus returned residual. Therefore the apparent
divergence-compatible second-order response of item~\emph{(v)} cannot be an
independent constitutive insertion; it is induced by the selected visible
effective law together with the canonical Schur return field. This is exactly
item~\emph{(iv)} of
Definition~\ref{def:schur-lifted-symmetric-square-response-package}.

Next, by
Theorem~\ref{thm:temporal-realization-forbids-primitive-internal-creation}, no
visible source on the selected cut is primitive creation. Any explicit visible
source term on \(\mathcal X_h^{\mathrm{sym}}\) must therefore be exchange with
another visible sector or with the hidden complement of the same coherent
bundle. The hidden-complement contribution is already absorbed into the
Schur-return part of the response by
Theorem~\ref{thm:canonical-selected-schur-realization}(d)--(e). Hence any
remaining explicit visible source on \(\mathcal X_h^{\mathrm{sym}}\) is
balanced exchange with the other selected visible carriers on the same bundle.
By Theorem~\ref{thm:coupled-exchange-balance}, this gives an exchange-balanced
coupled family; and because \(\mathcal X_h^{\mathrm{sym}}\) is a symmetric
square carrier, that visible exchange datum is represented on it by a
symmetric exchange tensor furnished by preceding selected carriers. This is
exactly item~\emph{(v)} of
Definition~\ref{def:schur-lifted-symmetric-square-response-package}.

Therefore the selected cut carries a Schur-lifted symmetric-square response
package, proving item~\emph{(a)}.

By
Theorem~\ref{thm:schur-lifted-symmetric-square-response-forces-geometric-witness},
the package of item~\emph{(a)} forces the geometric witness and hence
nonemptiness of the causal-structural Schur-induced mode class on the selected
bundle. Since item~\emph{(i)} already places phase, wave, and gauge beyond the
existence stage,
Corollary~\ref{cor:geometry-only-remaining-intrinsic-existence-lane} now shows
that geometry is the last intrinsic lane to close. The same theorem also says
what remains after that closure: only the carrier-level collapse to the
trace-adjusted-Ricci / Einstein representative. This proves item~\emph{(b)}.

Finally, item~\emph{(i)} contains the first-three-lanes package together with
canonical realization of persistence modes. Item~\emph{(a)} supplies exactly
the remaining witness input isolated by
Theorem~\ref{thm:geometric-witness-is-exact-remaining-intrinsic-obstruction}.
Hence that theorem applies and yields item~\emph{(c)}: the four intrinsic
flagship law classes are forced on one selected bundle.
\end{proof}

\begin{proof}[Proof of Theorem~\ref{thm:geometry-unique-lift-of-primitive-three-channel}]
\label{proof:thm:geometry-unique-lift-of-primitive-three-channel}
Item~\emph{(i)} is exactly item~\emph{(b)} of
Theorem~\ref{thm:three-channel-schur-tensor-lift-forces-geometric-lane}. For
item~\emph{(ii)}, the same theorem already shows that the geometric lane is
forced by a symmetric-square lift of a selected \(3\)-channel fiber. By
Definition~\ref{def:five-persistence-modes} and
Definition~\ref{def:canonical-sector-generation}, causal-structural
persistence is not a fourth primitive channel but the canonical
symmetric-response lift of the primitive \(3\)-channel. Therefore the forced
geometric law class is realized exactly as that lift. Item~\emph{(iii)} is now
Corollary~\ref{cor:flagship-family-exhausts-five-persistence-modes}: within the
scoped intrinsic flagship family there is no additional intrinsic lift beyond
the geometric one. Finally, item~\emph{(iv)} is the remaining-burden clause of
item~\emph{(b)} in
Theorem~\ref{thm:three-channel-schur-tensor-lift-forces-geometric-lane}.
\end{proof}

\begin{proof}[Proof of Theorem~\ref{thm:three-forcing-bridges-from-bedrock-to-primitive-physics}]
\label{proof:thm:three-forcing-bridges-from-bedrock-to-primitive-physics}
Item~\emph{(a)} is exactly
Theorem~\ref{thm:temporal-interface-from-d3} together with
Corollary~\ref{cor:temporal-interface-precedes-physical-time}: the earlier
structural chapters together with the continuum bridge force the unique
temporal interface but do not yet read it as physical time.

For item~\emph{(b)}, item \emph{(ii)} is exactly
Corollary~\ref{cor:derived-selected-schur-realization-forces-three-primitive-channels}:
once the temporal interface and the continuum-bridge realized horizon tower force the
selected Schur realization on the least-motion cut, the remaining intrinsic
hypotheses already force the three primitive visible channels there, with
phase, wave, and gauge as their canonical direct realizations. Read through
the intrinsic \(1+2+3\) decomposition, these are the three successive direct
bands of the carrier tower itself: local-state persistence on the first band,
propagative persistence on the second, and comparative persistence on the
third.

For item~\emph{(c)}, item \emph{(iii)} adds exactly the intrinsic \(1+2+3\)
decomposition needed by
Theorem~\ref{thm:derived-selected-schur-realization-forces-primitive-three-bilinear-lift-package}.
That theorem yields the faithful six-dimensional symmetric bilinear visible
self-coupling of
Definition~\ref{def:primitive-three-symmetric-bilinear-lift-package}, and
Corollary~\ref{cor:derived-selected-schur-realization-forces-geometric-lift-and-four-intrinsic-laws}
then yields the Schur tensor lift, the geometric law class, and the four
intrinsic flagship law classes on the same selected bundle. Geometry is
therefore identified as the unique intrinsic lift of the primitive
\(3\)-channel rather than a fourth primitive channel.

Finally, item~\emph{(d)} is the combination of
Theorem~\ref{thm:derived-selected-schur-realization-forces-first-three-intrinsic-law-classes},
Corollary~\ref{cor:derived-selected-schur-realization-forces-geometric-lift-and-four-intrinsic-laws},
and Remark~\ref{rem:remaining-intrinsic-endpoint-obligations}: once those
three forcing bridges are complete, the remaining work is no longer existence
or taxonomy, but only the canonical representative collapse on each already
forced intrinsic lane. In particular, the endpoint-rigidity packages of the
flagship sections are no longer carrying the burden of channel existence; they
are compatibility/readout packages for lanes already fixed by the common
selected-bundle forcing theorem.
\end{proof}

\begin{proof}[Proof of Corollary~\ref{cor:current-forcing-frontier-for-primitive-physics}]
\label{proof:cor:current-forcing-frontier-for-primitive-physics}
The first bridge is already complete by
Theorem~\ref{thm:temporal-interface-from-d3} and
Corollary~\ref{cor:temporal-interface-precedes-physical-time}. The second is
already complete by
Corollary~\ref{cor:derived-selected-schur-realization-forces-three-primitive-channels}.
The third is complete by
Theorem~\ref{thm:derived-selected-schur-realization-forces-primitive-three-bilinear-lift-package}
and
Corollary~\ref{cor:derived-selected-schur-realization-forces-geometric-lift-and-four-intrinsic-laws}.
The phase and wave closures are then discharged by
Theorems~\ref{thm:phase-local-endpoint-collapse} and
\ref{thm:wave-local-endpoint-collapse}, and the gauge closure is discharged by
Theorem~\ref{thm:gauge-local-endpoint-collapse}, while the geometric closure is
discharged by Theorem~\ref{thm:geometric-local-endpoint-collapse}. The optional
constant-mass refinement is discharged by
Theorem~\ref{thm:homogeneous-wave-scalar-profile-forces-constant-mass}. The
theorem just cited therefore yields the stated frontier. The remaining
carrier-specific endpoint packages are therefore to be read as canonical
representative collapses and compatibility checks, not as independent sectoral
existence assumptions.
\end{proof}

\begin{proof}[Proof of Theorem~\ref{thm:mode-compatible-selection-forces-least-motion-realization}]
\label{proof:thm:mode-compatible-selection-forces-least-motion-realization}
Because \(\mathfrak R_{\mathfrak m}(h)\) is an admissible realized class,
Theorem~\ref{thm:selection-interface} applies to every datum in that class.
Thus the realization problem on \(\mathfrak R_{\mathfrak m}(h)\) reduces to a
finite and checkable reduced selection problem. By
Definition~\ref{def:mode-compatible-candidate-class}(iv), this reduced problem
carries the observer-motion design objective used to define least motion on
that class. When the projection choice is part of the datum, existence of a
minimizer follows from Theorem~\ref{thm:existence-minimizers}; when the
selected cut is already fixed, the same clause simply says that the finite
reduced parameter space already carries the observer-motion minimization
problem. In either case, at least one observer datum in
\(\mathfrak R_{\mathfrak m}(h)\) minimizes observer motion on that class.

By hypothesis, the admissible reduced selection data modulo horizon-preserving
equivalence form a singleton. Corollary~\ref{cor:selection-to-forcing}
therefore forces one selected observed law on
\(\mathfrak R_{\mathfrak m}(h)\), unique up to horizon-preserving equivalence.
Hence any two observer-motion minimizers in the class are
horizon-preserving equivalent and determine the same selected visible law.

Let \(\mathcal Y_{\mathfrak m}\subseteq \mathcal X_h\) be the visible
subcarrier generated by one such minimizing datum. By
Definition~\ref{def:mode-compatible-candidate-class}(i),
\(\mathcal Y_{\mathfrak m}\) is faithful, reconstruction-compatible,
closure-admissible, and stable under admissible promotion. By
Definition~\ref{def:mode-compatible-candidate-class}(ii), it preserves the
compatibility signature and admissible source profile of the mode
\(\mathfrak m\). By
Definition~\ref{def:mode-compatible-candidate-class}(iii), no proper visible
subcarrier of \(\mathcal Y_{\mathfrak m}\) still lies in the same
mode-compatible class, so no proper visible subcarrier still preserves that
mode signature together with the same faithful reconstruction,
closure-admissibility, and promotion stability. By construction,
\(\mathcal Y_{\mathfrak m}\) minimizes observer motion on the class. These are
exactly the clauses of Definition~\ref{def:least-motion-characteristic-observer}.
Therefore the datum generated by \(\mathcal Y_{\mathfrak m}\) is a
least-motion characteristic observer for the mode \(\mathfrak m\), unique up
to horizon-preserving equivalence.
\end{proof}

\begin{proof}[Proof of Corollary~\ref{cor:mode-wise-candidate-forcing-yields-compatibility-surface}]
\label{proof:cor:mode-wise-candidate-forcing-yields-compatibility-surface}
Apply
Theorem~\ref{thm:mode-compatible-selection-forces-least-motion-realization} to
each of the four candidate classes named in hypotheses
\emph{(ii)}--\emph{(v)}. This yields:
\begin{enumerate}[label=(\alph*), leftmargin=2em, itemsep=0.2em]
\item a unique least-motion scalar visible realization of local-state
      persistence carrying phase-compatible reversible transport together with
      scalar recursive closure;
\item a unique least-motion scalar visible realization of propagative
      persistence carrying reversible second-order transport together with
      scalar endpoint closure;
\item a unique least-motion one-form comparison realization of comparative
      persistence carrying exactness and local gauge covariance;
\item a unique least-motion symmetric-response realization of the geometric
      lift of the primitive \(3\)-channel carrying divergence-compatible
      second-order response;
\end{enumerate}
Each realization is faithful, closure-admissible, stable under admissible
promotion, and unique up to horizon-preserving equivalence by the theorem just
proved. These are exactly the subcarrier-level hypotheses listed in items
\emph{(ii)}--\emph{(v)} of
Theorem~\ref{thm:inherited-compatibility-surface-forces-mode-complete-realization}.
Hence that inherited compatibility surface holds on the selected cut.
\end{proof}

\begin{proof}[Proof of Theorem~\ref{thm:inherited-compatibility-surface-forces-mode-complete-realization}]
\label{proof:thm:inherited-compatibility-surface-forces-mode-complete-realization}
Item~\emph{(ii)} supplies a faithful visible realization of local-state
persistence on \(\mathcal X_h\): the relevant compatibility signature is
phase-compatible reversible scalar transport, the visible source profile is
source-free up to exchange or returned residual by the temporal realization
principle, and the stated minimality condition makes the chosen scalar
subcarrier irreducible for that signature. Item~\emph{(iii)} does the same for
propagative persistence, now with reversible second-order scalar transport and
scalar endpoint closure. Item~\emph{(iv)} gives a faithful visible realization
of comparative persistence on a one-form comparison subcarrier carrying the
exactness and local gauge-covariance signature. Item~\emph{(v)} gives a faithful
visible realization of the geometric lift on a symmetric-response subcarrier
carrying divergence-compatible response.

In each case, the subcarrier is faithful, closure-admissible, and stable under
admissible promotion by assumption, and item~\emph{(vi)} supplies uniqueness
up to horizon-preserving equivalence inside the corresponding mode-compatible
class. Therefore each of the three primitive channels together with the
geometric lift of Definition~\ref{def:five-persistence-modes} admits at least
one faithful, reconstruction-compatible, admissible-promotion-stable visible
realization on the one selected observed bundle \(\mathcal X_h\). This is
exactly Definition~\ref{def:mode-complete-visible-realization}.
\end{proof}

\begin{proof}[Proof of Theorem~\ref{thm:mode-complete-realization-forces-intrinsic-selected-bundle}]
\label{proof:thm:mode-complete-realization-forces-intrinsic-selected-bundle}
Because a canonical selected Schur bridge fixes one least-motion characteristic
observer datum, it fixes one selected projector and hence one selected observed
bundle
\[
\mathcal X_h=P_h\cH^\bullet
\]
on the physically realized cut. Thus every faithful visible realization named
in Definition~\ref{def:mode-complete-visible-realization} is already realized
on one and the same selected observed bundle.

By hypothesis~\emph{(i)}, each of the three primitive persistence channels
together with the geometric lift admits at least one faithful visible
realization on that bundle. By hypothesis~\emph{(iii)}, each such role has a
canonical minimal compatible realization on that same bundle: phase for
local-state persistence, wave for propagative persistence, gauge for
comparative persistence, and geometric for the symmetric-response lift of the
primitive \(3\)-channel. Therefore the four intrinsic flagship carrier classes
arise on the selected bundle itself.

By Theorem~\ref{thm:persistence-mode-separation}, these structural roles are
pairwise distinct. By
Corollary~\ref{cor:flagship-family-exhausts-five-persistence-modes}, within the
present intrinsic flagship scope there is no fourth primitive channel and no
additional intrinsic lift. Hence the four canonical realizations above are
exactly the intrinsic flagship carrier classes on that selected bundle. Since
all of them arise on the same physically realized selected observed bundle,
this is precisely
Definition~\ref{def:intrinsic-selected-bundle-existence}.
\end{proof}

\begin{proof}[Proof of Theorem~\ref{thm:scoped-intrinsic-sector-generation}]
\label{proof:thm:scoped-intrinsic-sector-generation}
By Theorem~\ref{thm:common-selected-bundle-assembly}, the four intrinsic
flagship carriers are associated visible carriers of one and the same selected
observed bundle, and each induced visible law is obtained by restriction of
the same selected visible effective law together with its inherited endpoint
closure class. Thus the intrinsic flagship sectors already live on one common
physical bundle and are not independent carrier-wise laws.

By Corollary~\ref{cor:intrinsic-symmetry-forces-123}, the local balanced
channel algebra on that bundle carries a unique intrinsic channel profile
\(1+2+3\). This removes any remaining ambiguity in the underlying local
channel algebra from which visible sectors may be generated.

By Theorem~\ref{thm:canonical-sector-generation-reduction}, the intrinsic
flagship carriers are exactly the canonical realizations of the three
primitive persistence channels together with the canonical symmetric-response
lift of the primitive \(3\)-channel on that selected observed bundle. By
Theorem~\ref{thm:existence-of-five-persistence-modes}, every one of those modes
is realized there. By Theorem~\ref{thm:persistence-mode-separation}, the modes
are pairwise distinct. By
Corollary~\ref{cor:flagship-family-exhausts-five-persistence-modes}, within the
scoped intrinsic flagship family there is no fourth primitive channel and no
additional intrinsic lift.
Therefore the intrinsic flagship sectors are determined intrinsically by the
distinct compatibility signatures and admissible source profiles carried by the
common selected law on the intrinsic \(1+2+3\) channel algebra.

It follows that, on this scope, sector generation is not a primitive dictionary
from channel dimensions to names. It is the intrinsic production of the unique
minimal compatible visible realizations of the three primitive persistence
channels together with the canonical geometric lift carried by the selected
bundle. This is exactly
Definition~\ref{def:intrinsic-sector-generation}.
\end{proof}

\begin{proof}[Proof of Theorem~\ref{thm:scoped-bundle-first-coherence-decomposition}]
\label{proof:thm:scoped-bundle-first-coherence-decomposition}
Item \emph{(i)} is exactly
Theorem~\ref{thm:scoped-intrinsic-sector-generation}: on the common selected
bundle the intrinsic sectors are generated from the same selected visible
effective law and the intrinsic \(1+2+3\) channel algebra, rather than chosen
carrier-by-carrier. Item \emph{(ii)} is exactly
Theorem~\ref{thm:flagship-mode-classification}. Item \emph{(iii)} is exactly
Corollary~\ref{cor:flagship-family-exhausts-five-persistence-modes}, which
shows that within the present intrinsic flagship scope there is no fourth
primitive channel and no additional intrinsic lift. For item \emph{(iv)},
Theorem~\ref{thm:resolution_integrated_balance_closure_not_intrinsic_mode}
shows that a hydrodynamic law arising from a resolution-integrated balance
hierarchy is not an additional primitive channel or intrinsic lift, but only
an effective closure of integrated unresolved exchange. The final sentence is
the joint restatement of these four clauses.
\end{proof}

\begin{proof}[Proof of Theorem~\ref{thm:no-auxiliary-exact-visible-restriction-forces-general-bundle-trichotomy}]
\label{proof:thm:no-auxiliary-exact-visible-restriction-forces-general-bundle-trichotomy}
Let \(\mathcal Y\subseteq \mathcal X_h\) be an irreducible faithful visible law
class satisfying clauses \emph{(i)}--\emph{(iii)} of
Definition~\ref{def:general-bundle-direct-lift-closure-trichotomy}. If
\(\mathcal Y\) appears only through a resolution-integrated effective balance
closure, then clause \emph{(c)} of that definition already holds. So assume
from now on that \(\mathcal Y\) is not merely an effective closure.

Because the restriction of the selected visible effective law to \(\mathcal Y\)
is exact-visible and self-sufficient, and because the selected cut satisfies
the no-auxiliary internal shadow elimination package, the restricted law on
\(\mathcal Y\) introduces no hidden labels, quotient choice, gauge slice, or
exogenous constitutive datum beyond the same coherent bundle. Hence the visible
role of \(\mathcal Y\) must be read entirely from the common selected visible
effective law together with the local balanced channel algebra already carried
by \(\mathcal X_h\).

By Theorem~\ref{thm:scoped-bundle-first-coherence-decomposition}, that local
channel algebra is intrinsically the \(1+2+3\) algebra, primitive visible
coherence on the selected bundle appears only as the three direct channels
together with the canonical geometric lift, and there is no fourth primitive
direct channel and no additional intrinsic lift on the present scope.
Therefore an irreducible exact-visible self-sufficient class that is not merely
an effective closure cannot carry any extra primitive visible role. Since
\(\mathcal Y\) is irreducible and minimal for its compatibility signature and
admissible source profile, it cannot be a composite of several primitive roles
inside the same exact-visible law. It must therefore be carried either by one
of the three primitive direct channels or by the canonical symmetric-response
lift of the primitive \(3\)-channel.

The direct-channel and lift cases are mutually exclusive by the same scoped
bundle-first decomposition, and the effective-closure case was separated at the
start. Hence exactly one of clauses \emph{(a)}--\emph{(c)} of
Definition~\ref{def:general-bundle-direct-lift-closure-trichotomy} holds for
\(\mathcal Y\). This is precisely the general-bundle direct/lift/closure
trichotomy.
\end{proof}

\begin{proof}[Proof of Theorem~\ref{thm:direct-lift-closure-trichotomy-forces-intrinsic-role-dichotomy}]
\label{proof:thm:direct-lift-closure-trichotomy-forces-intrinsic-role-dichotomy}
Let \(\mathcal Y\subseteq \mathcal X_h\) be an irreducible faithful visible law
class satisfying clauses \emph{(i)}--\emph{(iii)} of
Definition~\ref{def:general-bundle-intrinsic-role-dichotomy}. By the
direct/lift/closure trichotomy, exactly one of cases \emph{(a)}--\emph{(c)} of
Definition~\ref{def:general-bundle-direct-lift-closure-trichotomy} holds.

If case \emph{(c)} holds, then \(\mathcal Y\) already appears only through a
resolution-integrated effective balance closure, which is exactly clause
\emph{(b)} of the intrinsic-role dichotomy.

Assume next that case \emph{(a)} holds. Because the selected bundle carries
the intrinsic \(1+2+3\) channel algebra and satisfies the scoped bundle-first
coherence decomposition of
Theorem~\ref{thm:scoped-bundle-first-coherence-decomposition}, each primitive
direct channel on that bundle is realized by exactly one of the three direct
intrinsic roles of
Definition~\ref{def:five-persistence-modes}: local-state persistence,
propagative persistence, or comparative persistence. Therefore \(\mathcal Y\)
preserves the compatibility signature and admissible source profile of one of
those three intrinsic roles, which is clause \emph{(a)} of
Definition~\ref{def:general-bundle-intrinsic-role-dichotomy}.

Finally, assume case \emph{(b)} holds. By
Definition~\ref{def:five-persistence-modes}, the canonical
divergence-compatible symmetric-response lift of the primitive \(3\)-channel is
exactly the causal-structural intrinsic role. Hence \(\mathcal Y\) preserves
the compatibility signature and admissible source profile of that fourth
intrinsic role, again yielding clause \emph{(a)} of
Definition~\ref{def:general-bundle-intrinsic-role-dichotomy}.

Thus every irreducible exact-visible self-sufficient faithful visible law class
on the selected bundle either preserves one of the four intrinsic role
signatures or appears only through a resolution-integrated effective balance
closure. This is precisely the general-bundle intrinsic-role dichotomy.
\end{proof}

\begin{proof}[Proof of Theorem~\ref{thm:intrinsic-role-compatible-exact-visible-law-classes-are-schur-induced}]
\label{proof:thm:intrinsic-role-compatible-exact-visible-law-classes-are-schur-induced}
Because the family satisfies a canonical selected Schur bridge, one least-motion
characteristic observer datum is fixed together with one common selected
observed bundle \(\mathcal X_h=P_h\cH^\bullet\) and one admissible realized
class \(\mathfrak R_{\mathrm{sel}}(h)\). Let \(R_\ast\in
\mathfrak R_{\mathrm{sel}}(h)\) be a realization whose selected datum is that
fixed least-motion datum.

We claim that \(R_\ast\) belongs to the Schur-induced mode candidate class
\(\mathfrak R^{\mathrm{Schur}}_{\mathfrak m}(h)\) of
Definition~\ref{def:schur-induced-mode-candidate-class}. Since
\(\mathcal Y\subseteq \mathcal X_h\) already lives on the common selected
bundle of \(R_\ast\), no further transport is needed: we may take the visible
subcarrier \(\mathcal Y_{R_\ast}:=\mathcal Y\).

Hypothesis \emph{(ii)} gives clause \emph{(i)} of
Definition~\ref{def:schur-induced-mode-candidate-class}: \(\mathcal Y\) is
faithful, reconstruction-compatible, closure-admissible, and stable under
admissible promotion. Hypothesis \emph{(iv)} gives clause \emph{(ii)} of that
definition: the restriction of the selected visible effective law to
\(\mathcal Y\) preserves the compatibility signature and admissible source
profile of the role \(\mathfrak m\). Hypothesis \emph{(iii)} gives clause
\emph{(iii)}: no proper visible subcarrier of \(\mathcal Y\) still satisfies
those same role conditions. Therefore \(R_\ast\in
\mathfrak R^{\mathrm{Schur}}_{\mathfrak m}(h)\), so the Schur-induced mode
candidate class for \(\mathfrak m\) is nonempty.
\end{proof}

\begin{proof}[Proof of Corollary~\ref{cor:intrinsic-role-dichotomy-forces-general-bundle-exact-visible-exhaustion-hypothesis}]
\label{proof:cor:intrinsic-role-dichotomy-forces-general-bundle-exact-visible-exhaustion-hypothesis}
Let \(\mathcal Y\subseteq \mathcal X_h\) be an irreducible faithful visible law
class satisfying the three clauses of
Definition~\ref{def:general-bundle-exact-visible-exhaustion-hypothesis}. By
the intrinsic-role dichotomy, either \(\mathcal Y\) preserves the
compatibility signature and admissible source profile of one intrinsic role
\(\mathfrak m\), or \(\mathcal Y\) appears only through a
resolution-integrated effective balance closure.

In the first case,
Theorem~\ref{thm:intrinsic-role-compatible-exact-visible-law-classes-are-schur-induced}
shows that the Schur-induced mode candidate class for \(\mathfrak m\) is
nonempty. Hence \(\mathcal Y\) satisfies clause \emph{(a)} of
Definition~\ref{def:general-bundle-exact-visible-exhaustion-hypothesis}. In the
second case, \(\mathcal Y\) satisfies clause \emph{(b)} of that same
definition. Therefore the general-bundle exact-visible exhaustion hypothesis
holds.
\end{proof}

\begin{proof}[Proof of Theorem~\ref{thm:general-bundle-exact-visible-exhaustion}]
\label{proof:thm:general-bundle-exact-visible-exhaustion}
Let \(\mathcal Y\subseteq \mathcal X_h\) be an irreducible faithful visible law
class whose restricted law is exact-visible and self-sufficient. By
Definition~\ref{def:general-bundle-exact-visible-exhaustion-hypothesis},
\(\mathcal Y\) either belongs to a nonempty Schur-induced mode candidate class
for one of the four intrinsic roles of
Definition~\ref{def:five-persistence-modes}, or appears only through a
resolution-integrated effective balance closure.

If the second case holds, then \(\mathcal Y\) is exactly an effective closure
of the type listed in item \emph{(v)}. By
Theorem~\ref{thm:resolution_integrated_balance_closure_not_intrinsic_mode}, it
is not an additional primitive direct channel or intrinsic lift.

Assume therefore that the first case holds. Then there exists one of the four
intrinsic roles \(\mathfrak m\) such that \(\mathcal Y\) lies in a nonempty
Schur-induced mode candidate class for \(\mathfrak m\). By
Theorem~\ref{thm:nonempty-schur-induced-mode-classes-have-singleton-reduced-selection},
that class has singleton reduced selection data modulo horizon-preserving
equivalence, so \(\mathcal Y\) determines the unique least-motion visible
realization of the role \(\mathfrak m\) on the selected cut.

By
Theorem~\ref{thm:scoped-bundle-first-coherence-decomposition}(i), the common
selected bundle generates intrinsically the canonical realizations of the
three primitive direct channels together with the canonical geometric lift.
Therefore the unique realization of \(\mathfrak m\) is respectively the phase,
wave, gauge, or geometric carrier according to which of the four intrinsic
roles \(\mathfrak m\) is. Hence \(\mathcal Y\) is equivalent to exactly one of
the four canonical intrinsic realizations listed in items
\emph{(i)}--\emph{(iv)}.

Finally, uniqueness follows from
Theorem~\ref{thm:scoped-bundle-first-coherence-decomposition}(ii)--(iii): each
intrinsic visible law on the selected bundle belongs to exactly one of the four
roles, and there is no additional primitive direct channel or intrinsic lift on
the present scope. Therefore no irreducible exact-visible self-sufficient law
class on the selected bundle can realize a primitive role outside items
\emph{(i)}--\emph{(iv)}, and any remaining visible macroscopic law belongs to
the effective-closure case \emph{(v)}.
\end{proof}

\begin{proof}[Proof of Theorem~\ref{thm:part-iv-law-completion}]
\label{proof:thm:part-iv-law-completion}
Item~\emph{(i)} unifies the bridge and residual-return requirements into one
exact Schur realization statement: the least-motion minimizer, the realized
horizon cut, the selected observed law, the canonical residual-return field,
and the HFT-2 irreversibility accounting are all read on the same selected
datum. Item~\emph{(ii)} supplies the common selected observed bundle
intrinsically, together with the stronger conclusion that the four intrinsic
flagship carriers inherit the same selected visible effective law and differ
only at the final carrier-level identification step, while the hydrodynamic
lane is read on that same bundle as a resolution-integrated effective balance
closure. Item~\emph{(iii)} is exactly the scoped intrinsic-generation
statement of Theorem~\ref{thm:scoped-intrinsic-sector-generation}: the four
intrinsic visible carriers are canonical minimal compatible realizations of
the three primitive channels together with the canonical geometric lift of the
primitive \(3\)-channel, not externally assigned labels, and that intrinsic
channel profile is itself already forced from the closure-balance/symmetry
surface; the finite promotion carrier used there is read directly from the
continuum bridge's closure-typed promotion bridge. Finally,
item~\emph{(iv)} resolves the carrier-level identification step isolated by
Theorem~\ref{thm:common-endpoint-rigidity-reduction} by identifying the unique
canonical representative on each intrinsic generated carrier together with the
canonical effective hydrodynamic normal form on the balance hierarchy, so the
completion step is exactly carrier-level classification together with those
representative witnesses. The abstract exchange law
of Theorem~\ref{thm:coupled-exchange-balance} then already gives the common
interaction grammar, so what remains downstream is only the concrete coupled
realization on realized carriers together with quantitative remainder control.
\end{proof}

\begin{proof}[Proof of Theorem~\ref{thm:returned-remainders-form-predictive-correction-surface}]
\label{proof:thm:returned-remainders-form-predictive-correction-surface}
By Theorem~\ref{thm:part-iv-law-completion}, the selected law on every flagship
lane is read on one common selected bundle, all observable forcing factors
through the canonical residual return field there, and returned residual is
exactly the Schur correction. Therefore every recognizable endpoint law on the
present scope is already the visible leading law together with the returned
residual correction carried by the same coherent bundle.

On the phase lane, Theorem~\ref{thm:phase-local-endpoint-collapse} gives the
selected Schr\"odinger-side representative with remainder
\(\mathsf R_{\mathrm{ph}}\), and
Theorem~\ref{thm:induced-phase-hamiltonian-law} identifies that remainder as an
exchange remainder rather than a primitive source. On the wave lane,
Theorem~\ref{thm:wave-local-endpoint-collapse} gives the selected wave
representative with remainder \(\mathsf R_{\mathrm{wav}}\), while
Theorem~\ref{thm:induced-wave-law} identifies it as an exchange remainder on
the selected cut. On the gauge lane,
Theorem~\ref{thm:gauge-local-endpoint-collapse} gives the selected Maxwell-side
representative with remainder \(\mathsf R_{\mathrm g}\), and
Theorem~\ref{thm:induced-gauge-compatible-law} identifies it as the remaining
local gauge-invariant exchange remainder. On the geometric lane,
Theorem~\ref{thm:geometric-local-endpoint-collapse} gives the selected
Einstein-side representative with remainder
\(\mathsf N_{\mathrm{grav}}(g,T)\), while the induced coefficients
\(\Lambda_{\mathrm{ind}}\) and \(\kappa_{\mathrm{ind}}\) are part of the same
selected realization. On the effective hydrodynamic lane,
Corollary~\ref{cor:hydrodynamic-closure-law} gives the selected closure normal
form with remainder \(\mathsf R_{\mathrm{hyd}}\) and induced transport
coefficients \(\mu_{\mathrm{ind}}\) and \(\zeta_{\mathrm{ind}}\).

Each of those flagship theorems further states that in the lossless
leading-order regime the corresponding remainder vanishes, leaving the
recognizable representative. Since
Theorem~\ref{thm:temporal-realization-forbids-primitive-internal-creation}
forbids primitive internal source creation on the selected cut, every nonzero
departure from the leading representatives must be carried either by returned
remainder or by the induced scales of that same realization. This is precisely
the stated predictive correction surface.
\end{proof}

\begin{proof}[Proof of Theorem~\ref{thm:least-motion-observer-equivalence-class-is-admissible-selected-class}]
\label{proof:thm:least-motion-observer-equivalence-class-is-admissible-selected-class}
Let \(\mathcal O_h\) be the given least-motion characteristic observer datum,
and let \(\mathfrak R_{\mathrm{sel}}(h)\) be its horizon-preserving
equivalence class.

Because \(\mathcal O_h\) is a characteristic observer datum in the sense of
Definition~\ref{def:characteristic-observer-datum}, it already carries the
selection datum, reduced candidate family, normalization package, and
reconstruction/promotion data needed to place it inside the continuum-bridge
selection interfaces. Horizon-preserving equivalence preserves the observed
differential, defect, lower-bound operator, candidate family, symmetry action,
and normalization package by
Definition~\ref{def:horizon-preserving-equivalence}. Therefore every datum in
\(\mathfrak R_{\mathrm{sel}}(h)\) inherits a transported selection datum
satisfying the admissibility conditions of
Definition~\ref{def:admissible-realized-class}. This proves item \emph{(i)}.

Items \emph{(i)}--\emph{(iv)} of
Definition~\ref{def:least-motion-characteristic-observer} are invariant under
horizon-preserving equivalence: faithfulness and reconstruction compatibility
are preserved by transport of the selected datum, closure-admissibility is
preserved because the reduced admissible problem is intertwined by the same
equivalence, and admissible-promotion stability is preserved because the
promotion data are part of the transported observer datum. Hence every datum in
\(\mathfrak R_{\mathrm{sel}}(h)\) is faithful, closure-admissible, and stable
under admissible promotion on the same selected bundle after transport. This
proves item \emph{(ii)}.

Clause \emph{(v)} of
Definition~\ref{def:least-motion-characteristic-observer} is exactly the
observer-motion design objective on the reduced selection problem. Since the
entire reduced selection surface is transported by horizon-preserving
equivalence, the same design objective is carried by
\(\mathfrak R_{\mathrm{sel}}(h)\). This proves item \emph{(iii)}.

Finally, every two data in \(\mathfrak R_{\mathrm{sel}}(h)\) are
horizon-preserving equivalent by construction. Therefore their admissible
reduced selection data represent the same class modulo
horizon-preserving equivalence. This proves item \emph{(iv)}.
\end{proof}

\begin{proof}[Proof of Theorem~\ref{thm:classwise-horizon-equivalence-collapses-reduced-selection}]
\label{proof:thm:classwise-horizon-equivalence-collapses-reduced-selection}
Let \(R,R'\in\mathfrak R_{\mathrm{adm}}(h)\). By hypothesis, they are related
by a horizon-preserving equivalence intertwining the admissible reduced
operator data and the normalization package. Therefore the reduced selection
datum carried by \(R\) represents the same class as the reduced selection
datum carried by \(R'\) once both are read modulo horizon-preserving
equivalence. Since \(R\) and \(R'\) were arbitrary, every admissible reduced
selection datum on \(\mathfrak R_{\mathrm{adm}}(h)\) lies in one and the same
equivalence class.
\end{proof}

\begin{proof}[Proof of Theorem~\ref{thm:temporal-realization-forces-selected-cut-forcing-package}]
\label{proof:thm:temporal-realization-forces-selected-cut-forcing-package}
Let \(\mathfrak R_{\mathrm{sel}}(h)\) be the horizon-preserving equivalence
class of the given least-motion observer datum \(\mathcal O_h\). By
Theorem~\ref{thm:least-motion-observer-equivalence-class-is-admissible-selected-class},
\(\mathfrak R_{\mathrm{sel}}(h)\) is an admissible realized class whose
reduced selection problem carries the observer-motion design objective and
whose admissible reduced selection data modulo horizon-preserving equivalence
form a singleton. Therefore the selection interface of
Theorem~\ref{thm:selection-interface} applies to every datum in that class, and
Corollary~\ref{cor:selection-to-forcing} forces one selected observed law
across \(\mathfrak R_{\mathrm{sel}}(h)\), unique up to horizon-preserving
equivalence.

Because \(\mathcal O_h\) is already a least-motion characteristic observer
datum, it is the least-motion minimizer on this selection surface. Hypothesis
\emph{(i)} identifies its finite selected projector with the realized horizon
cut. Hence clauses \emph{(i)}--\emph{(iii)} of
Definition~\ref{def:selected-cut-forcing-package} all hold on one and the same
datum.
\end{proof}

\begin{proof}[Proof of Theorem~\ref{thm:part-iii-realized-horizon-tower-forces-selected-schur-hft2-return-package}]
\label{proof:thm:part-iii-realized-horizon-tower-forces-selected-schur-hft2-return-package}
By Theorem~\ref{thm:hilbert-completion}, the stage projectors
\((\CCPi{n})_{n\in\bbN}\) form a faithful transitive horizon tower on the
Hilbert completion of the chosen directed horizon family. By the displayed
definition, the persisting selected visible law on that same horizon family is
read stagewise as the exact effective operator of the one coherent operator
\(A_{\mathrm{coh}}\). Thus the stagewise family
\((\mathcal L_{\mathrm{vis},n}(z))_{n\in\bbN}\) is exactly the Schur tower of
\(A_{\mathrm{coh}}\) on that realized horizon tower. Applying
Theorem~\ref{thm:exact-schur-coherence-general} yields item \emph{(i)} of
Definition~\ref{def:selected-schur-hft2-return-package}: the selected observed
law carries an exact Schur tower on the realized horizon tower. This proves
item \emph{(a)}.

For the HFT-2 side, the same tower arises from orthogonal projection onto
nested stage subspaces of the Hilbert completion, so adjacent refinement
layers split orthogonally. Therefore the hypotheses of
Theorem~\ref{thm:HFT2-general} are satisfied on this realized horizon tower,
and the HFT-2 energy accounting holds there. Equivalently,
Corollary~\ref{cor:horizon-error-budget} gives the exact shadow-energy budget
on that same tower. The visible forcing term is therefore read exactly as the
returned visible trace of hidden/interface residual carried by the same Schur
tower, and observable irreversibility is exactly the unrecovered stock in the
HFT-2 accounting. This proves items \emph{(b)} and \emph{(c)}.

Items \emph{(a)} and \emph{(c)} are precisely the two clauses of
Definition~\ref{def:selected-schur-hft2-return-package}. Hence item
\emph{(d)} follows immediately.
\end{proof}

\begin{proof}[Proof of Corollary~\ref{cor:temporal-realization-on-part-iii-horizon-tower-forces-grounded-selected-schur-package}]
\label{proof:cor:temporal-realization-on-part-iii-horizon-tower-forces-grounded-selected-schur-package}
By hypothesis, the datum satisfies the selected-cut forcing package. By
Theorem~\ref{thm:part-iii-realized-horizon-tower-forces-selected-schur-hft2-return-package},
the same datum satisfies the selected Schur/HFT-2 return package. Therefore
Corollary~\ref{cor:selected-cut-plus-schur-return-forces-grounded-package}
gives the grounded selected Schur package. The equivalence with the canonical
selected Schur bridge is then
Theorem~\ref{thm:canonical-selected-schur-bridge-iff-grounded-package}.
\end{proof}

\begin{proof}[Proof of Theorem~\ref{thm:derived-temporal-interface-and-part-iii-horizon-tower-force-canonical-selected-schur-bridge}]
\label{proof:thm:derived-temporal-interface-and-part-iii-horizon-tower-force-canonical-selected-schur-bridge}
By Theorem~\ref{thm:temporal-realization-forces-selected-cut-forcing-package},
the least-motion datum \((P_h,A_{\mathrm{coh}})\) furnished by the stated
selection hypotheses satisfies the selected-cut forcing package. The added
continuum-bridge horizon-tower hypothesis is exactly the remaining hypothesis of
Corollary~\ref{cor:temporal-realization-on-part-iii-horizon-tower-forces-grounded-selected-schur-package}.
Therefore that corollary yields the grounded selected Schur package on the same
datum, hence also the canonical selected Schur bridge.
\end{proof}

\begin{proof}[Proof of Theorem~\ref{thm:derived-selected-schur-realization-forces-first-three-intrinsic-law-classes}]
\label{proof:thm:derived-selected-schur-realization-forces-first-three-intrinsic-law-classes}
The setting of
Corollary~\ref{cor:derived-temporal-interface-and-part-iii-horizon-tower-force-canonical-selected-schur-realization}
is exactly the setting of
Theorem~\ref{thm:derived-temporal-interface-and-part-iii-horizon-tower-force-canonical-selected-schur-bridge}.
Hence the same datum satisfies the canonical selected Schur bridge. The
present additional hypotheses are exactly the remaining hypotheses of
Theorem~\ref{thm:first-three-intrinsic-law-classes-already-forced}. Therefore
that theorem applies directly.
\end{proof}

\begin{proof}[Proof of Corollary~\ref{cor:derived-selected-schur-realization-forces-three-primitive-channels}]
\label{proof:cor:derived-selected-schur-realization-forces-three-primitive-channels}
The present setting includes the setting of
Theorem~\ref{thm:derived-selected-schur-realization-forces-first-three-intrinsic-law-classes},
hence in particular the setting of
Theorem~\ref{thm:derived-temporal-interface-and-part-iii-horizon-tower-force-canonical-selected-schur-bridge}
and the remaining hypotheses of
Theorem~\ref{thm:first-three-intrinsic-law-classes-already-forced}. Therefore
Theorem~\ref{thm:three-primitive-channels-already-forced} applies directly.
\end{proof}

\begin{proof}[Proof of Theorem~\ref{thm:derived-selected-schur-realization-forces-primitive-three-bilinear-lift-package}]
\label{proof:thm:derived-selected-schur-realization-forces-primitive-three-bilinear-lift-package}
The present setting includes the setting of
Theorem~\ref{thm:derived-selected-schur-realization-forces-first-three-intrinsic-law-classes},
hence the setting of
Theorem~\ref{thm:derived-temporal-interface-and-part-iii-horizon-tower-force-canonical-selected-schur-bridge}
together with the remaining hypotheses of
Theorem~\ref{thm:first-three-intrinsic-law-classes-already-forced}. The
additional intrinsic \(1+2+3\) decomposition hypothesis is exactly the only
extra input required by
Theorem~\ref{thm:canonical-selected-schur-bridge-forces-primitive-three-bilinear-lift-package}.
Therefore that theorem applies directly.
\end{proof}

\begin{proof}[Proof of Corollary~\ref{cor:derived-selected-schur-realization-forces-geometric-lift-and-four-intrinsic-laws}]
\label{proof:cor:derived-selected-schur-realization-forces-geometric-lift-and-four-intrinsic-laws}
By
Theorem~\ref{thm:derived-selected-schur-realization-forces-primitive-three-bilinear-lift-package},
the same selected cut carries the primitive \(3\)-channel symmetric bilinear
lift package. The common setting already includes the canonical selected Schur
bridge together with the remaining first-three-lanes hypotheses. Therefore
Theorems~\ref{thm:primitive-three-symmetric-bilinear-lift-forces-schur-tensor-lift}
and~\ref{thm:geometry-unique-lift-of-primitive-three-channel} apply in
sequence and yield the stated conclusions.
\end{proof}

\begin{proof}[Proof of Theorem~\ref{thm:canonical-selected-schur-bridge-iff-grounded-package}]
\label{proof:thm:canonical-selected-schur-bridge-iff-grounded-package}
The two conditions share clauses \emph{(i)}--\emph{(iii)} and
\emph{(v)}--\emph{(vi)} verbatim. The only difference is clause
\emph{(iv)}. In
Definition~\ref{def:canonical-selected-schur-bridge}, clause \emph{(iv)}
requires a grounded physical bridge on the selected datum and immediately
states that the self-sufficient realization law is already available there.
Definition~\ref{def:grounded-selected-schur-package} replaces that bridge-side
wording by the corresponding law-side content itself, namely the
self-sufficient realization principle on the same selected cut. Since
Definition~\ref{def:canonical-selected-schur-bridge} further identifies exact
visibility, intrinsic closure, no exogenous constitutive completion, and
constitutive appearance only as returned residual as the operative content of
that grounded bridge, the two formulations are equivalent.
\end{proof}

\begin{proof}[Proof of Theorem~\ref{thm:canonical-selected-schur-realization}]
\label{proof:thm:canonical-selected-schur-realization}
The selected-cut layer already identifies the finite minimizer with the
realized horizon cut, exports forcing across the admissible class by
horizon-preserving transport, and supplies the HFT-2 stock/flux
irreversibility accounting on that same selected datum. The grounded physical
bridge then contributes the law-side core of self-sufficient realization on
that cut: exact visibility, intrinsic closure, no exogenous constitutive
completion, and constitutive appearance only as returned residual. The local
facts of observable support and cut commutation are derived from the exact-
visible clause, and the exact Schur identity is then a direct consequence of
the general projection/target calculus. The residual return field is the
canonical return field induced by
the selected cut and selected-horizon shadow propagator, so returned residual
equals the Schur correction by construction and the selected law is visible
closure plus returned residual. The \(D=3\) barrier together with the intrinsic
\(1+2+3\) channel profile are already inherited from Parts~I--II.
\end{proof}

\begin{proof}[Proof of Corollary~\ref{cor:canonical-selected-schur-realization-yields-temporal-realization}]
\label{proof:cor:canonical-selected-schur-realization-yields-temporal-realization}
By Theorem~\ref{thm:canonical-selected-schur-realization}(a)--(b), the
realization-selection clause of
Definition~\ref{def:physical-realization-principle} holds on the selected cut:
the least-motion minimizer and realized cut coincide, and the selected visible
law is forced across the admissible class by horizon-preserving equivalence.

By Theorem~\ref{thm:canonical-selected-schur-realization}(c), the
self-sufficient realization clause holds there: the selected law is exactly
visible and intrinsically closed on that cut, no exogenous constitutive
completion is admitted, and any constitutive appearance is returned residual
from the same coherent bundle.

By Theorem~\ref{thm:canonical-selected-schur-realization}(d)--(f), the
exchange-return clause also holds: the canonical return field is exactly the
Schur correction, visible forcing is returned residual on that same selected
datum, and observable irreversibility is measured by unrecovered HFT-2
residual. These are exactly the three clauses of
Definition~\ref{def:physical-realization-principle}.
\end{proof}

\begin{proof}[Proof of Theorem~\ref{thm:canonical-selected-schur-realization-forces-coherent-law-exact-elimination}]
\label{proof:thm:canonical-selected-schur-realization-forces-coherent-law-exact-elimination}
By
Theorem~\ref{thm:canonical-selected-schur-realization}(c), the selected law is
exactly visible and intrinsically closed on the selected cut, with no exogenous
constitutive completion and all constitutive appearance interpreted as returned
residual. By
Theorem~\ref{thm:canonical-selected-schur-realization}(d), the canonical
residual-return field on that cut is the return field induced directly by the
selected law and the selected-horizon shadow propagator, and its returned
residual is exactly the Schur correction. By
Theorem~\ref{thm:canonical-selected-schur-realization}(e), the selected
observed law is exactly visible closure plus returned residual on that cut.
Using Definition~\ref{def:selected-visible-effective-law}, this means precisely
that the selected observed law is the selected visible effective law
\(\mathcal L_{\mathrm{vis},h}(z)\), that the selected return field is the
residual return field of the same coherent operator and selected projector, and
that every hidden-side visible contribution is read only through the exact
Schur correction / returned residual of the same coherent bundle. These are
exactly clauses \emph{(i)}--\emph{(iii)} of
Definition~\ref{def:coherent-law-exact-elimination-package}.
\end{proof}

\begin{proof}[Proof of Corollary~\ref{cor:grounded-selected-schur-package-closes-intrinsic-surface}]
\label{proof:cor:grounded-selected-schur-package-closes-intrinsic-surface}
By
Theorem~\ref{thm:canonical-selected-schur-bridge-iff-grounded-package}, the
grounded selected Schur package is exactly the canonical selected Schur bridge.
Hence Theorem~\ref{thm:canonical-selected-schur-realization} applies, so the
selected datum already carries one exact realization problem in which physical
visibility, residual return, irreversibility bookkeeping, and the intrinsic
\(1+2+3\) channel profile are read together. The three forcing bridges are
then complete by
Theorem~\ref{thm:three-forcing-bridges-from-bedrock-to-primitive-physics}, and
their present manuscript frontier is exactly
Corollary~\ref{cor:current-forcing-frontier-for-primitive-physics}. That
corollary yields items \emph{(a)}--\emph{(d)} directly. Item \emph{(e)} is the
closing clause of the same corollary.
\end{proof}

\begin{proof}[Proof of Theorem~\ref{thm:layered-selected-schur-realization-closes-present-scope-trichotomy}]
\label{proof:thm:layered-selected-schur-realization-closes-present-scope-trichotomy}
By
Corollary~\ref{cor:layered-selected-schur-realization-closes-intrinsic-surface},
the intrinsic flagship surface is already closed on the selected bundle: the
three forcing bridges are complete, the phase, wave, gauge, and geometric law
classes are forced on that bundle, their recognizable representative
identities are closed, and no intrinsic channel/lift classification problem
remains open on the present flagship scope.

By Theorem~\ref{thm:scoped-bundle-first-coherence-decomposition}, the same
selected bundle decomposes bundle-first into exactly the three primitive direct
channels together with the canonical geometric lift, while any hydrodynamic
law arising there is only a resolution-integrated effective closure and not an
additional primitive channel or intrinsic lift. Thus every irreducible
exact-visible self-sufficient faithful visible law class in the generated
family falls into one of those three cases.

The cases are mutually exclusive on the present scope: the direct-channel and
geometric-lift roles are pairwise distinct by the bundle-first decomposition,
and the effective-closure case is disjoint from the intrinsic roles by
Theorem~\ref{thm:resolution_integrated_balance_closure_not_intrinsic_mode}.
Hence exactly one of \emph{(a)}--\emph{(c)} holds for each such law class.
\end{proof}

% --- Proofs from bridge_program.tex ---

\begin{proof}[Proof of Proposition~\ref{prop:observer-triad}]
\label{proof:prop:observer-triad}
The pullback identity
\[
\frac{d}{dt}\bigl(\Phi_{t,0}^*\Theta_t\bigr)
=
\Phi_{t,0}^*(\partial_t\Theta_t + \mathcal L_{V_t}\Theta_t)
=
\Phi_{t,0}^*(\mathcal D_\chi \Theta_t)
\]
shows that the Eulerian and Lagrangian formulations are equivalent. For the
characteristic presentation, fix $a \in M$ and write $\gamma_a(t)=\Phi_{t,0}(a)$.
The chain rule gives
\[
\frac{d}{dt}\Theta_t(\gamma_a(t))
=
(\partial_t\Theta_t + \mathcal L_{V_t}\Theta_t)(\gamma_a(t))
=
(\mathcal D_\chi \Theta_t)(\gamma_a(t)).
\]
Thus the Eulerian equation vanishes identically if and only if the transported
quantity is constant along every characteristic. The three descriptions are
therefore equivalent.
\end{proof}

\begin{proof}[Proof of Lemma~\ref{lem:flux-differential}]
\label{proof:lem:flux-differential}
By Cartan's formula and the conservative hypothesis,
\[
d(\iota_{\widetilde V}\nu)=\mathcal L_{\widetilde V}\nu=0.
\]
Hence
\[
dJ_\Theta
= d\Theta \wedge \iota_{\widetilde V}\nu
  + \Theta\, d(\iota_{\widetilde V}\nu)
= d\Theta \wedge \iota_{\widetilde V}\nu.
\]
For any volume form $\nu$ and vector field $\widetilde V$,
$d\Theta \wedge \iota_{\widetilde V}\nu = (\widetilde V\Theta)\nu$, so
\[
dJ_\Theta = (\widetilde V\Theta)\nu
          = (\partial_t\Theta + \mathcal L_{V_t}\Theta)\nu
          = (\mathcal D_\chi \Theta)\nu.
\]
\end{proof}

\begin{proof}[Proof of Theorem~\ref{thm:cut-conservation}]
\label{proof:thm:cut-conservation}
By Lemma~\ref{lem:flux-differential}, $\mathcal D_\chi \Theta = 0$ implies
$dJ_\Theta=0$. Stokes' theorem therefore yields
\[
0 = \int_U dJ_\Theta = \int_{\partial U} J_\Theta
  = \int_{\Sigma_1} J_\Theta - \int_{\Sigma_0} J_\Theta
    + \int_{\Sigma_{\mathrm{side}}} J_\Theta.
\]
Since $\widetilde V$ is tangent to $\Sigma_{\mathrm{side}}$,
$\iota_{\widetilde V}\nu$ restricts trivially there, so
$\int_{\Sigma_{\mathrm{side}}} J_\Theta=0$. The claimed equality follows.
\end{proof}

\begin{proof}[Proof of Theorem~\ref{thm:characteristic-ibp}]
\label{proof:thm:characteristic-ibp}
Apply Stokes' theorem to the form $fg\,\iota_{\widetilde V}\nu$. Using again
$d(\iota_{\widetilde V}\nu)=0$,
\begin{align*}
d\bigl(fg\,\iota_{\widetilde V}\nu\bigr)
&= d(fg)\wedge \iota_{\widetilde V}\nu \\
&= \bigl(\widetilde V(fg)\bigr)\nu \\
&= \bigl((\mathcal D_\chi f)g + f(\mathcal D_\chi g)\bigr)\nu.
\end{align*}
Integrating over $U$ gives the result.
\end{proof}

\begin{proof}[Proof of Proposition~\ref{prop:cellular-cut-conservation}]
\label{proof:prop:cellular-cut-conservation}
By the defining relation between coboundary and boundary,
\[
\langle J,\Gamma_1-\Gamma_0\rangle
=
\langle J,\partial C\rangle
=
\langle \delta J,C\rangle.
\]
Since $\delta J$ vanishes on every $2$-cell of $C$, the right-hand side is
zero. Therefore $\langle J,\Gamma_1\rangle = \langle J,\Gamma_0\rangle$.
\end{proof}

\begin{proof}[Proof of Lemma~\ref{lem:through-horizon-decomposition}]
\label{proof:lem:through-horizon-decomposition}
Since $u=\CCPi{h}u$ and $\id=\CCPi{h}+\CCPibar{h}$,
\[
du = (\CCPi{h}+\CCPibar{h})d\CCPi{h}u
   = \CCPi{h}d\CCPi{h}u + \CCPibar{h}d\CCPi{h}u
   = \CCObs{d}{h}u + \CCLeak{d}{h}u.
\]
Apply $d$ and use $d^2=0$:
\[
0 = d^2u = d\,\CCObs{d}{h}u + d\,\CCLeak{d}{h}u.
\]
Projecting to the observed sector yields
\[
0 = \CCPi{h}d\,\CCObs{d}{h}u + \CCPi{h}d\,\CCLeak{d}{h}u.
\]
By definition,
\[
\CCPi{h}d\,\CCObs{d}{h}u
= \CCPi{h}d\CCPi{h}d\CCPi{h}u
= \CCObs{d}{h}^2u,
\]
and
\[
\CCPi{h}d\,\CCLeak{d}{h}u
= \CCPi{h}d\CCPibar{h}d\CCPi{h}u
= \CCDefect{d}{h}u.
\]
The claim follows.
\end{proof}

\begin{proof}[Proof of Theorem~\ref{thm:characteristic-realization}]
\label{proof:thm:characteristic-realization}
The identity is exactly the operator statement obtained in
Lemma~\ref{lem:through-horizon-decomposition}. The commutation relations with
$U_t$ follow directly from the hypotheses
$U_t d = d U_t$ and $U_t \CCPi{h} = \CCPi{h} U_t$:
\[
U_t \CCObs{d}{h}
= U_t \CCPi{h} d \CCPi{h}
= \CCPi{h} d \CCPi{h} U_t
= \CCObs{d}{h} U_t,
\]
and the same calculation gives the defect commutation relation.
\end{proof}

\begin{proof}[Proof of Theorem~\ref{thm:dynamics-triad}]
\label{proof:thm:dynamics-triad}
Regularity of $L$ makes the Legendre transform a local diffeomorphism, so the
Euler--Lagrange equations are equivalent to Hamilton's equations in the
standard way. For the Hamilton--Jacobi step, set
$p(q,t):=\partial_q S(q,t)$. Differentiating the Hamilton--Jacobi equation with
respect to $q$ gives
\[
\partial_t p + \partial_q H(q,p,t) + \partial_p H(q,p,t)\,\partial_q p = 0.
\]
Along a curve satisfying $\dot q=\partial_p H(q,p,t)$, the chain rule yields
\[
\dot p = \partial_t p + \dot q\,\partial_q p = -\partial_q H(q,p,t),
\]
which is Hamilton's equation for $p$. Thus every Hamilton--Jacobi
characteristic is a Hamiltonian trajectory. Conversely, on a simply connected
domain, any Hamiltonian trajectory lying in an exact Lagrangian graph
$p=\partial_q S$ integrates to a local generating function $S$, and that
function satisfies the Hamilton--Jacobi equation. The three descriptions
therefore determine the same local dynamics.
\end{proof}

\begin{proof}[Proof of Proposition~\ref{prop:hj-projection}]
\label{proof:prop:hj-projection}
If $(q(t),p(t))$ is a Hamiltonian trajectory contained in $\Gamma_t$, then
$p(t)=\partial_q S(q(t),t)$ and Hamilton's equation gives
\[
\dot q(t) = \partial_p H\bigl(q(t),p(t),t\bigr)
          = \partial_p H\bigl(q(t),\partial_q S(q(t),t),t\bigr)
          = V_t(q(t)).
\]
Thus the configuration-space projection of the Hamiltonian trajectory is an
integral curve of $V_t$.

Conversely, let $q(t)$ be an integral curve of $V_t$ and define
$p(t):=\partial_q S(q(t),t)$. Differentiate along the curve:
\[
\dot p(t)=\partial_t(\partial_q S)(q(t),t)+\partial_q^2S(q(t),t)\,\dot q(t).
\]
Using $\dot q(t)=V_t(q(t))=\partial_p H(q(t),p(t),t)$ and differentiating the
Hamilton--Jacobi equation with respect to $q$ yields
\[
\dot p(t) = -\,\partial_q H(q(t),p(t),t).
\]
Therefore $(q(t),p(t))$ satisfies Hamilton's equations and lies in $\Gamma_t$.
The final statement follows by Definition~\ref{def:characteristic-presentation}.
\end{proof}

\begin{proof}[Proof of Proposition~\ref{prop:selection-problem}]
\label{proof:prop:selection-problem}
By Corollary~\ref{cor:transfer-part-one}, the Part~I calculus applies to every
member of $\mathfrak R$. By Theorem~\ref{thm:dynamics-triad} and
Proposition~\ref{prop:hj-projection}, the Lagrangian, Hamiltonian, and
characteristic descriptions of the same realized dynamics are synchronized on
that class. The realized operator algebra is therefore fixed uniformly across
$\mathfrak R$. What remains undecided is which members of $\mathfrak R$ satisfy
any further admissibility, coercivity, symmetry, or normalization conditions.
Those additional restrictions are exactly selection hypotheses on the realized
class.
\end{proof}

\begin{proof}[Proof of Theorem~\ref{thm:selection-interface}]
\label{proof:thm:selection-interface}
By Corollary~\ref{cor:transfer-part-one}, every
$R\in\mathfrak R_{\mathrm{adm}}(h)$ carries the structural calculus of Part~I
on its observed sector, proving (i).

For (ii), the finite-dimensionality of $\mathcal Q_h(R)$ is part of the
definition of admissible realized class. Since both $D_h$ and $\mathcal Q_h(R)$
are $G_h$-equivariant, Theorem~\ref{thm:normal-form} block-diagonalizes the
problem across isotypic components, Corollary~\ref{cor:isotypic-block} identifies
the free variables with multiplicity-space operators, and
Corollary~\ref{cor:parameter-count} gives the resulting finite parameter count.

For (iii), Theorem~\ref{thm:symmetry-reduction-admissibility} reduces the
completion constraint and its certificates to independent blockwise problems.
Theorem~\ref{thm:compiler-contract} then gives exactly the three stated outputs:
a unique admissible completion, a minimal face certifying residual freedom, or
inconsistency of the constraint set.

For (iv), if projection choice is also part of the realized datum, then
Theorem~\ref{thm:existence-minimizers} supplies existence of optimal horizons
for continuous design objectives in finite dimensions, while
Theorem~\ref{thm:stationarity-condition} gives the stated stationarity
criterion.
\end{proof}

\begin{proof}[Proof of Theorem~\ref{thm:intrinsic-equivariant-forcing}]
\label{proof:thm:intrinsic-equivariant-forcing}
By Corollary~\ref{cor:selection-to-forcing}, the selected observed law is
forced on $\mathfrak R_{\mathrm{adm}}(h)$ up to horizon-preserving
equivalence. Transport each selected law to the canonical carrier by $W_R$.
Because the transported datum is $S_4$-equivariant by
Definition~\ref{def:intrinsically-equivariant-realized-class}, the transported
selected law commutes with the intrinsic relabeling action of $S_4$ on
$V_{\mathrm{slot}}$. Theorem~\ref{thm:intrinsic-slot-normal-form} therefore
gives unique scalars $a,b,c$ producing the displayed formula. Since the
selected law is forced up to horizon-preserving equivalence and all realizations
have been transported into the same canonical carrier, those scalars are the
same for every realization in the class. The eigenvalue formula is exactly
Corollary~\ref{cor:intrinsic-channel-eigenvalues}.
\end{proof}

\begin{proof}[Proof of Corollary~\ref{cor:realized-phase-class-forcing}]
\label{proof:cor:realized-phase-class-forcing}
By Corollary~\ref{cor:selection-to-forcing}, the selected observed law is
forced on $\mathfrak R_{\mathrm{adm}}(h)$ up to horizon-preserving
equivalence. For each realization, the selected-law sector carries a
complex-type multiplicity package by
Definition~\ref{def:complex-type-realized-class}, so
Theorem~\ref{thm:scalar-commutant-collapses-phase-gauge} supplies a unique
phase-carrier class modulo phase-equivalence on that sector. Because the
selected law itself is forced up to horizon-preserving equivalence, these
phase-carrier classes agree across the realized class up to the same
equivalence.
\end{proof}

\begin{proof}[Proof of Corollary~\ref{cor:realized-finite-observable-rigidity}]
\label{proof:cor:realized-finite-observable-rigidity}
By Corollary~\ref{cor:selection-to-forcing}, the selected observed law is
forced on $\mathfrak R_{\mathrm{adm}}(h)$ up to horizon-preserving
equivalence. On each channel,
Theorem~\ref{thm:finite-observable-rigidity} upgrades the stated observable
hypothesis to the full finite-rigidity package, and then
Corollary~\ref{cor:finite-observable-rigidity-phase} collapses the residual
phase gauge to a unique class modulo phase-equivalence. These channelwise
phase classes therefore agree across the realized class up to the same
equivalence.
\end{proof}

\begin{proof}[Proof of Corollary~\ref{cor:realized-phase-generation-implies-rigidity}]
\label{proof:cor:realized-phase-generation-implies-rigidity}
For every relevant selected-law channel, item (ii) of
Definition~\ref{def:phase-generated-realized-class} says that the selected-law
family is phase-generating in the sense of
Definition~\ref{def:phase-generating-observable-family}. By
Corollary~\ref{cor:phase-generation-forces-phase-rigidity}, the corresponding
multiplicity-space family is phase-rigid. The remaining items of
Definition~\ref{def:phase-rigid-realized-class} are already part of the
phase-generated hypothesis. Therefore
$\mathfrak R_{\mathrm{adm}}(h)$ is phase-rigid on the selected channels.
\end{proof}

\begin{proof}[Proof of Corollary~\ref{cor:observable-complete-implies-observable-rigid}]
\label{proof:cor:observable-complete-implies-observable-rigid}
For each relevant selected-law channel, item (ii) and
Theorem~\ref{thm:observable-completeness-irreducible} give
observable-irreducibility, while item (ii) and
Theorem~\ref{thm:observable-completeness-scalar-commutant} give phase rigidity
in the sense of Definition~\ref{def:phase-rigid-observable-family}. Together
with items (i) and (iii), this shows that
$\mathfrak R_{\mathrm{adm}}(h)$ is phase-rigid on the selected channels in the
sense of Definition~\ref{def:phase-rigid-realized-class}, and also
observable-irreducible on the selected channels in the sense of
Definition~\ref{def:observable-irreducible-realized-class}. Therefore both
clauses of Definition~\ref{def:observable-rigid-realized-class} hold.
\end{proof}

\begin{proof}[Proof of Corollary~\ref{cor:scalar-commutant-complex-type-implies-observable-rigid}]
\label{proof:cor:scalar-commutant-complex-type-implies-observable-rigid}
For each relevant selected-law channel, item (i) gives the complex-type
multiplicity package together with scalar multiplicity commutant, which is
exactly the phase-rigid condition of
Definition~\ref{def:phase-rigid-observable-family} on the corresponding
multiplicity space. Together with item (iii), this shows that
$\mathfrak R_{\mathrm{adm}}(h)$ is phase-rigid on the selected channels in the
sense of Definition~\ref{def:phase-rigid-realized-class}. Item (ii) is exactly
the observable-irreducible clause of
Definition~\ref{def:observable-rigid-realized-class}. Therefore both clauses
of Definition~\ref{def:observable-rigid-realized-class} hold.
\end{proof}

\begin{proof}[Proof of Corollary~\ref{cor:transport-rigid-implies-scalar-recursive}]
\label{proof:cor:transport-rigid-implies-scalar-recursive}
Fix a relevant selected channel fiber $F_h$. By item (ii) of
Definition~\ref{def:transport-rigid-realized-class} and
Corollary~\ref{cor:realized-reversible-recursion-forcing}, every bi-infinite
selected channel transport sequence on $F_h$ satisfies the recursion law
\[
x_{n+1}=B_hx_n-x_{n-1}.
\]
By items (ii) and (iii) of
Definition~\ref{def:transport-rigid-realized-class} and
Corollary~\ref{cor:realized-quadratic-transport-rigidity}, there exists a
unique scalar $\kappa_h\ge 0$ such that
\[
B_h^*B_h=\kappa_h\,I_{F_h}.
\]
These are exactly the clauses of
Definition~\ref{def:scalar-recursive-selected-channel-system}.
\end{proof}

\begin{proof}[Proof of Corollary~\ref{cor:transport-generation-yields-scalar-recursive}]
\label{proof:cor:transport-generation-yields-scalar-recursive}
Corollary~\ref{cor:realized-transport-generation-forcing} gives both the unique
scalar $\kappa_h$ with
\[
B_h^*B_h=\kappa_h\,I_{F_h}
\]
and the recursion law
\[
x_{n+1}=B_hx_n-x_{n-1}.
\]
These are exactly the clauses of
Definition~\ref{def:scalar-recursive-selected-channel-system}.
\end{proof}

\begin{proof}[Proof of Corollary~\ref{cor:transport-generated-implies-transport-rigid}]
\label{proof:cor:transport-generated-implies-transport-rigid}
Item (i) of Definition~\ref{def:transport-rigid-realized-class} is part of the
assumption. For each relevant selected channel fiber, items (i) and (ii) of
Definition~\ref{def:transport-generated-selected-channel-system} give the
order-one forcing data and bounded transport base required in item (ii) of
Definition~\ref{def:transport-rigid-realized-class}. Item (iii) of
Definition~\ref{def:transport-generated-selected-channel-system} says exactly
that the selected channel is without preferred direction at quadratic order,
which is a quadratic transport rigidity package and therefore a quadratically
rigid selected channel system. Hence item (iii) of
Definition~\ref{def:transport-rigid-realized-class} also holds.
\end{proof}

\begin{proof}[Proof of Corollary~\ref{cor:observable-orbit-frame-implies-symmetry-generated}]
\label{proof:cor:observable-orbit-frame-implies-symmetry-generated}
Item (i) places $\mathfrak R_{\mathrm{adm}}(h)$ in case (v) of
Definition~\ref{def:observable-rigidity-generated-realized-class}, so clause
(i) of Definition~\ref{def:symmetry-generated-realized-class} holds. Item
(ii) is exactly clause (ii) of
Definition~\ref{def:symmetry-generated-realized-class}.

Fix a relevant selected channel fiber $F_h$. By item (iii) and
Corollary~\ref{cor:realized-orbit-direction-frame}, the selected direction
family is a tight isotropic direction frame on $F_h$. Together with item (iv),
this places the selected channel system in the scope of
Definition~\ref{def:direction-frame-selected-channel-system}. Combining that
with item (v) gives clauses (iii) and (iv) of
Definition~\ref{def:symmetry-generated-realized-class}.

Finally, items (vi)--(viii) are exactly clauses (v)--(vii) of
Definition~\ref{def:symmetry-generated-realized-class}. Therefore every clause
of that definition holds.
\end{proof}

\begin{proof}[Proof of Theorem~\ref{cor:realized-local-internal-datum-forcing}]
\label{proof:cor:realized-local-internal-datum-forcing}
By item (i) of
Definition~\ref{def:governing-datum-rigid-realized-class} and
Corollaries~\ref{cor:realized-phase-rigidity-forcing} and
\ref{cor:realized-self-adjoint-scalarization}, the selected law carries a
forced phase-carrier class modulo phase-equivalence on every relevant channel,
and every selected self-adjoint closure operator on the channel is scalar on
the corresponding multiplicity space. By item (ii) of
Definition~\ref{def:governing-datum-rigid-realized-class}, each relevant
selected channel is transport-rigid; therefore
Corollary~\ref{cor:transport-rigid-implies-scalar-recursive} supplies the
unique scalar quadratic recursion datum of
Definition~\ref{def:scalar-recursive-selected-channel-system}.
These are exactly the data packaged in
Definition~\ref{def:local-internal-datum-selected-channels}.
\end{proof}

% --- Proofs from flagship_hydrodynamics.tex ---

\begin{proof}[Proof of Theorem~\ref{thm:kinetic-moment-hierarchy}]
\label{proof:thm:kinetic-moment-hierarchy}
Multiply the kinetic law of
Definition~\ref{def:kinetic-transport-branch} by \(\varphi(\xi)\) and
integrate over \(\Xi\). This gives
\[
\partial_t M_\varphi + \nabla\!\cdot F_\varphi
=
\int_\Xi \varphi(\xi)\int_\Xi K(x,t;\xi,\eta)f(x,\eta,t)\,d\mu(\eta)\,d\mu(\xi)
-\int_\Xi \varphi(\xi)\sigma(x,t;\xi)f(x,\xi,t)\,d\mu(\xi).
\]
Swap the order of integration in the gain term and then rewrite the loss term
using the definition of \(\sigma\):
\begin{align*}
\partial_t M_\varphi + \nabla\!\cdot F_\varphi
&=
\int_\Xi\int_\Xi \varphi(\xi)K(x,t;\xi,\eta)f(x,\eta,t)\,d\mu(\eta)\,d\mu(\xi) \\
&\qquad
-\int_\Xi\int_\Xi \varphi(\eta)K(x,t;\xi,\eta)f(x,\eta,t)\,d\mu(\eta)\,d\mu(\xi) \\
&=
\int_\Xi\int_\Xi
\bigl(\varphi(\xi)-\varphi(\eta)\bigr)
K(x,t;\xi,\eta)f(x,\eta,t)\,d\mu(\eta)\,d\mu(\xi) \\
&=
Q_\varphi.
\end{align*}
If \(\varphi\) is a collision invariant, then the gain integral equals the
loss integral pointwise, so \(Q_\varphi=0\).
\end{proof}

\begin{proof}[Proof of Theorem~\ref{thm:least-motion-hydrodynamic-observer}]
\label{proof:thm:least-motion-hydrodynamic-observer}
By hypothesis \emph{(ii)}, the carrier \(\mathcal X_{\mathrm{hyd}}\) is
faithful, closure-admissible, and stable under admissible promotion. By
hypothesis \emph{(iii)}, no proper subcarrier still supports exact visible
mass-momentum balance together with those same structural properties.
Hypothesis \emph{(iv)} says that any competing characteristic observer datum
with those properties is horizon-preserving equivalent to the datum generated
by \(\mathcal X_{\mathrm{hyd}}\). Therefore the datum generated by
\(\mathcal X_{\mathrm{hyd}}\) satisfies the clauses of
Definition~\ref{def:least-motion-characteristic-observer} and is unique up to
horizon-preserving equivalence. Hence \(\mathcal X_{\mathrm{hyd}}\) is the
unique least-motion characteristic observer for the realized kinetic family.
\end{proof}

\begin{proof}[Proof of Corollary~\ref{cor:hydrodynamic-moment-balance}]
\label{proof:cor:hydrodynamic-moment-balance}
Apply Theorem~\ref{thm:kinetic-moment-hierarchy} first with
\(\varphi(\xi)=1\), which gives
\[
M_\varphi=U_0,
\qquad
F_\varphi=\int_\Xi \xi f\,d\mu = U_1.
\]
Then apply it componentwise with \(\varphi(\xi)=\xi_i\), which gives
\[
M_{\xi_i}=(U_1)_i,
\qquad
F_{\xi_i,j}=\int_\Xi \xi_i\xi_j f\,d\mu=(U_2)_{ij}.
\]
The collision-invariance hypothesis makes the production terms vanish in both
cases.
\end{proof}

\begin{proof}[Proof of Theorem~\ref{cor:hydrodynamic-closure-law}]
\label{proof:cor:hydrodynamic-closure-law}
The first balance law follows from
Corollary~\ref{cor:hydrodynamic-moment-balance} and the identity
\(U_1=U_0 w\). For the second, decompose
\[
U_2 = U_0 w\otimes w + sI + \Theta.
\]
Substituting this into
\[
\partial_t U_1 + \nabla\!\cdot U_2 = 0
\]
and using \(U_1=U_0 w\) gives
\[
\partial_t(U_0 w)
+
\nabla\!\cdot(U_0 w\otimes w)
+
\nabla s
=
-\nabla\!\cdot \Theta.
\]
By hypothesis, \(\Theta\) is the exact characteristic return on the selected
cut. Because this cut is least-motion, no smaller visible carrier remains both
faithful and closure-admissible at this scale, so unresolved higher moments can
re-enter the visible law only through that exact shadow return. By
Theorem~\ref{thm:temporal-realization-forbids-primitive-internal-creation},
this return is an exact balanced exchange with the unresolved kinetic sector,
not a primitive source insertion.
Corollary~\ref{cor:characteristic-boundary-reduction} removes the bulk part of
the characteristic integral, so after the stated slow-variation assumption only
the first endpoint operator on \(\nabla w\) remains at leading order.
Objectivity removes the antisymmetric contribution, hence the endpoint operator
acts on the symmetric gradient. The symmetric tensor representation splits into
the scalar trace channel \(\mathbb RI\) and the traceless channel
\(\mathrm{Sym}^2_0(\mathbb R^d)\). The no-preferred-direction hypothesis makes
the endpoint operator orthogonally equivariant, so each channel is invariant,
and symmetry-scalar rigidity forces a scalar action on each channel. This
produces unique coefficients \(\zeta_{\mathrm{ind}}\) and
\(2\mu_{\mathrm{ind}}\) on the trace and traceless channels respectively, with
all higher-order endpoint terms absorbed into \(\mathsf R_{\mathrm{hyd}}\).
Substituting the resulting expression for \(\Theta\) gives the displayed law.
\end{proof}

\begin{proof}[Proof of Corollary~\ref{cor:recognizable-navier-stokes-form}]
\label{proof:cor:recognizable-navier-stokes-form}
Expand
\[
2\mu_{\mathrm{ind}}\,\mathsf D_0(w)
+
\zeta_{\mathrm{ind}}(\nabla\!\cdot w)I
\]
using the definition of \(\mathsf D_0(w)\). The traceless part contributes
\(\mu_{\mathrm{ind}}(\nabla w+\nabla w^\top)\), and the remaining scalar part is
\((\zeta_{\mathrm{ind}}-\frac{2}{d}\mu_{\mathrm{ind}})(\nabla\!\cdot w)I\),
which is exactly \(\lambda_{\mathrm{ind}}(\nabla\!\cdot w)I\).
\end{proof}

\begin{proof}[Proof of Corollary~\ref{cor:hydrodynamic-schur-remainder-identity}]
\label{proof:cor:hydrodynamic-schur-remainder-identity}
By Theorem~\ref{thm:canonical-selected-schur-realization}, every hidden-side
visible contribution on the selected cut is carried by the common returned
residual field, and that field is exactly the Schur correction
\(\CCSigma{h}{z}\). By Corollary~\ref{cor:hydrodynamic-closure-law}, the
selected hydrodynamic endpoint map on the resolution-integrated balance
hierarchy yields unique induced transport coefficients and a unique higher-order
remainder term. Therefore the common Schur correction has a unique image on the
selected hydrodynamic endpoint; denote it by \(\Sigma_{h,\mathrm{hyd}}(z)\). By
Theorem~\ref{thm:returned-remainders-form-predictive-correction-surface}, the
hydrodynamic remainder is exactly that returned correction on the
resolution-integrated hydrodynamic lane, so
\(\mathsf R_{\mathrm{hyd}}=\Sigma_{h,\mathrm{hyd}}(z)\). Substituting this into
Corollary~\ref{cor:hydrodynamic-closure-law} gives the displayed identities.
\end{proof}

% --- Proofs from flagship_thermodynamics.tex ---

\begin{proof}[Proof of Theorem~\ref{thm:selected-history-second-law}]
\label{proof:thm:selected-history-second-law}
For \(N=0\), both displayed formulas are empty-sum identities. For \(N>0\),
apply Theorem~\ref{thm:effective-thermodynamic-closure}(iii) successively at
horizons \(h_\ast,h_\ast+1,\dots,h_\ast+N-1\):
\[
\mathrm{Stock}_{\mathrm{unrec}}(x;h_\ast+j)
=
\mathrm{Flux}_{\mathrm{return}}(x;h_\ast+j)
+
\mathrm{Stock}_{\mathrm{unrec}}(x;h_\ast+j+1).
\]
Summing these identities over \(j=0,\dots,N-1\) telescopes all intermediate
stock terms and yields item \emph{(i)}:
\[
\mathrm{Stock}_{\mathrm{unrec}}(x;h_\ast)
-
\mathrm{Stock}_{\mathrm{unrec}}(x;h_\ast+N)
=
\sum_{j=0}^{N-1}\mathrm{Flux}_{\mathrm{return}}(x;h_\ast+j).
\]
Item \emph{(ii)} is the difference of the \(N+1\) and \(N\) cases of item
\emph{(i)}. By Theorem~\ref{thm:effective-thermodynamic-closure}(iv), the
unrecovered stock is monotone under one-step refinement, so each increment in
item \emph{(ii)} is nonnegative. Therefore
\(\mathcal S_{\mathrm{th}}^{(N)}(x)\) is nonnegative and nondecreasing, which
is item \emph{(iii)}. Since
\(\mathrm{Stock}_{\mathrm{unrec}}(x;h_\ast+N)\ge 0\), the defining formula of
Definition~\ref{def:selected-history-entropy-production} gives
\[
\mathcal S_{\mathrm{th}}^{(N)}(x)
\le
\mathrm{Stock}_{\mathrm{unrec}}(x;h_\ast).
\]
Equality holds exactly when the remainder term
\(\mathrm{Stock}_{\mathrm{unrec}}(x;h_\ast+N)\) vanishes. This is item
\emph{(iv)}.
\end{proof}

\begin{proof}[Proof of Corollary~\ref{cor:selected-history-equilibrium-criterion}]
\label{proof:cor:selected-history-equilibrium-criterion}
By Theorem~\ref{thm:selected-history-second-law}(ii),
\[
\mathcal S_{\mathrm{th}}^{(N+1)}(x)-\mathcal S_{\mathrm{th}}^{(N)}(x)
=
\mathrm{Flux}_{\mathrm{return}}(x;h_\ast+N).
\]
Therefore item \emph{(i)} is equivalent to item \emph{(ii)}. Using the
one-step balance law of
Theorem~\ref{thm:effective-thermodynamic-closure}(iii),
\[
\mathrm{Stock}_{\mathrm{unrec}}(x;h_\ast+N)
=
\mathrm{Flux}_{\mathrm{return}}(x;h_\ast+N)
+
\mathrm{Stock}_{\mathrm{unrec}}(x;h_\ast+N+1),
\]
the vanishing of the returned flux is equivalent to constancy of the
unrecovered stock. Hence item \emph{(ii)} is equivalent to item \emph{(iii)}.
\end{proof}

\begin{proof}[Proof of Theorem~\ref{thm:selected-tower-thermodynamic-state-law}]
\label{proof:thm:selected-tower-thermodynamic-state-law}
Item \emph{(i)} is just the selected-horizon stock/flux balance law of
Theorem~\ref{thm:effective-thermodynamic-closure}(iii), together with the
one-step increment law of Theorem~\ref{thm:selected-history-second-law}(ii).
For item \emph{(ii)}, combine Definition~\ref{def:selected-history-entropy-production}
with the notation above:
\[
U_N(x)+\Sigma_N(x)
=
\mathrm{Stock}_{\mathrm{unrec}}(x;h_\ast+N)
\bigl(
\mathrm{Stock}_{\mathrm{unrec}}(x;h_\ast)
-
\mathrm{Stock}_{\mathrm{unrec}}(x;h_\ast+N)
\bigr)
=
U_0(x).
\]
Item \emph{(iii)} follows from
Theorem~\ref{thm:effective-thermodynamic-closure}(iv) and
Theorem~\ref{thm:selected-history-second-law}(iii). Item \emph{(iv)} follows
from item \emph{(ii)} together with the nonnegativity of both \(U_N(x)\) and
\(\Sigma_N(x)\).
\end{proof}

\begin{proof}[Proof of Theorem~\ref{thm:canonical-selected-tower-entropy-functional}]
\label{proof:thm:canonical-selected-tower-entropy-functional}
We argue by induction on \(N\). The case \(N=0\) is immediate from
\(E_0(x)=0=\mathcal S_{\mathrm{th}}^{(0)}(x)\). Assume
\(E_N(x)=\mathcal S_{\mathrm{th}}^{(N)}(x)\). Then
\[
E_{N+1}(x)
=
E_N(x)+\Phi_N(x)
=
\mathcal S_{\mathrm{th}}^{(N)}(x)+\Phi_N(x).
\]
By Theorem~\ref{thm:selected-history-second-law}(ii),
\[
\mathcal S_{\mathrm{th}}^{(N+1)}(x)-\mathcal S_{\mathrm{th}}^{(N)}(x)
=
\Phi_N(x),
\]
so
\[
E_{N+1}(x)
=
\mathcal S_{\mathrm{th}}^{(N+1)}(x).
\]
The induction closes for all \(N\).
\end{proof}

% --- Proofs from interface_portability.tex ---

\begin{proof}[Proof of Proposition~\ref{prop:monograph-factors-through-interface-package}]
\label{proof:prop:monograph-factors-through-interface-package}
Corollary~\ref{cor:temporal-interface-precedes-physical-time} gives the
temporal interface datum, which is item \emph{(i)}. The admissible faithful
visible cuts and their horizon-preserving equivalence are exactly the observer
comparison surface used in the selected-cut dynamics chapters, especially in
the observer-covariance clause of
Definition~\ref{def:physical-realization-principle} and in the scoped
temporal-completion class of
Definition~\ref{def:scoped-temporal-completion-class}; this gives item
\emph{(ii)}. On each such admissible cut, the selected visible effective law is
already defined by exact internal elimination on one coherent bundle by
Definition~\ref{def:selected-visible-effective-law}, giving item \emph{(iii)}.
Finally, the same selected-cut surface carries the exact residual-return /
Schur grammar and the exact carrier-side endpoint closure class, by
Definition~\ref{def:selected-visible-effective-law} and
Proposition~\ref{prop:return-residual-is-schur}. This is item \emph{(iv)}.
\end{proof}

\begin{proof}[Proof of Corollary~\ref{cor:portability-of-temporal-realization-and-common-grammar}]
\label{proof:cor:portability-of-temporal-realization-and-common-grammar}
The preceding Part~II law-side arguments above the carrier-specific readout
stage use exactly the items listed in
Definition~\ref{def:pre-dynamical-temporal-interface-package} together with
the Temporal Realization Law and its operational clauses in
Definition~\ref{def:physical-realization-principle}. They do not additionally
use the counting derivation of the package once the package itself is already
given. Therefore the same arguments apply verbatim on any external host
realization of that package. The final sentence is immediate: every later
carrier theorem whose hypotheses are stated only in terms of that package and
its carrier-local compatibility data may be read on the same host without
changing the law-side logic.
\end{proof}

\begin{proof}[Proof of Theorem~\ref{thm:microscopic-selected-bundle-readout-forces-carrier-closure-transport}]
\label{proof:thm:microscopic-selected-bundle-readout-forces-carrier-closure-transport}
Definition~\ref{def:microscopic-selected-bundle-readout} packages four carrier
readout families on one selected bundle together with one common transport law
for the induced continuum closure class on that same selected cut. In the Lean
formalization this is recorded by the detached structure
\lean{ClosureCurrent.BundleReadout}, sitting above one selected-bundle assembly
and one detached bridge datum. The generic bridge theorem
\lean{ClosureCurrent.SelectedBundleForcing.carrierClosureReadout} identifies any
carrier observer read from that bundle with the same selected physical law and
then invokes the active carrier closure readout theorem on that observer. The
packaged surface theorem \lean{ClosureCurrent.SelectedBundleForcing.readoutSurface}
simply applies that argument to the phase, wave, gauge, and geometric readout
families. This is exactly the theorem statement.
\end{proof}

\begin{proof}[Proof of Theorem~\ref{thm:detached-microscopic-forcing-interface}]
\label{proof:thm:detached-microscopic-forcing-interface}
The detached microscopic forcing interface is formalized by
\lean{ClosureCurrent.IntrinsicSectorForcing}. Its selected-bundle component
\lean{ClosureCurrent.SelectedBundleForcing} supplies, by
\lean{ClosureCurrent.SelectedBundleForcing.surface}, a nonempty intrinsic
selected-bundle datum together with the matched bedrock-carrying generator,
autonomy on the selected cut, and fiber invariance there. Its sector-generation
component supplies, by
\lean{ClosureCurrent.IntrinsicSectorForcing.surface}, the active Step~3
generation package: intrinsic generation, minimal compatible realizations,
closure-typed finite-carrier control, and the canonical \([3,2,1]\) logical
profile. The theorem \lean{ClosureCurrent.forcingSurface} is the conjunction of
those two outputs, with no further hidden hypotheses. Hence once the detached
microscopic law furnishes exactly those two interface layers, the present
upstream law-completion seam closes at the current active surface.
\end{proof}

\begin{proof}[Proof of Theorem~\ref{thm:role-indexed-microscopic-law-forces-selected-bundle-closure-readout}]
\label{proof:thm:role-indexed-microscopic-law-forces-selected-bundle-closure-readout}
Definition~\ref{def:role-indexed-microscopic-carrier-family} is formalized by
the pair \lean{ClosureCurrent.VisibleCarrierRole} and
\lean{ClosureCurrent.RoleReadoutFamily}. The detached microscopic selected-bundle
law itself is formalized by \lean{ClosureCurrent.MicroscopicBundleLaw}. That law
first collapses to the earlier detached bundle-readout interface by the
constructor \lean{ClosureCurrent.MicroscopicBundleLaw.toSelectedBundleForcing}.
The active selected-bundle existence and autonomy statements then follow from
\lean{ClosureCurrent.SelectedBundleForcing.surface}. The role-indexed carrier
closure inheritance is packaged separately by
\lean{ClosureCurrent.MicroscopicBundleLaw.readoutSurface}, which applies the
generic carrier closure readout theorem uniformly to each visible role in the
indexed family. Combining these two outputs yields exactly the theorem
statement, and this conjunction is formalized by
\lean{ClosureCurrent.MicroscopicBundleLaw.surface}.
\end{proof}

\begin{proof}[Proof of Theorem~\ref{thm:role-indexed-microscopic-sector-law-closes-current-forcing-seam}]
\label{proof:thm:role-indexed-microscopic-sector-law-closes-current-forcing-seam}
The role-indexed microscopic sector law is formalized by
\lean{ClosureCurrent.MicroscopicSectorLaw}. Its bundle component is the
role-indexed microscopic law of
Theorem~\ref{thm:role-indexed-microscopic-law-forces-selected-bundle-closure-readout};
its additional data are precisely the intrinsic non-dictionary sector-generation
clauses already isolated in the detached forcing layer. The constructor
\lean{ClosureCurrent.MicroscopicSectorLaw.toIntrinsicSectorForcing} recovers the
earlier detached forcing package. Applying the compact forcing theorem
\lean{ClosureCurrent.forcingSurface} gives intrinsic selected-bundle existence,
intrinsic sector generation, bedrock-carrying selected generator, selected-cut
autonomy, fiber invariance, and the canonical logical profile \([3,2,1]\).
Applying \lean{ClosureCurrent.MicroscopicBundleLaw.readoutSurface} to the bundle
component adds the role-indexed carrier closure transport. These are exactly the
claims grouped together by \lean{ClosureCurrent.MicroscopicSectorLaw.surface}.
\end{proof}

\begin{proof}[Proof of Theorem~\ref{thm:two-law-detached-microscopic-package-forces-role-indexed-selected-bundle-surface}]
\label{proof:thm:two-law-detached-microscopic-package-forces-role-indexed-selected-bundle-surface}
Definitions~\ref{def:detached-exact-exchange-current-law} and
\ref{def:detached-event-exactification-autonomy-law} are formalized by
\lean{ClosureCurrent.ExactExchangeCurrentLaw} and
\lean{ClosureCurrent.EventExactificationLaw}. The exact exchange current law
supplies the current-side grammar witness, the bundle-first assembly law, the
matched selected bridge datum, and the role-indexed visible readout family.
The event exactification/autonomy law adds the local lossless exactification,
relabeling-equivariance, autonomous visible readout, returned-defect clause,
and the intrinsic selected-bundle clauses. The constructor
\lean{ClosureCurrent.EventExactificationLaw.toMicroscopicBundleLaw} collapses
that two-law package to the previously isolated role-indexed microscopic
selected-bundle law. The resulting bundle existence, autonomy, fiber
invariance, and role-indexed carrier closure transport are then exactly the
surface grouped by \lean{ClosureCurrent.EventExactificationLaw.surface}.
\end{proof}

\begin{proof}[Proof of Corollary~\ref{cor:two-law-detached-microscopic-package-plus-intrinsic-sector-clauses-closes-current-forcing-seam}]
\label{proof:cor:two-law-detached-microscopic-package-plus-intrinsic-sector-clauses-closes-current-forcing-seam}
The corollary is formalized by
\lean{ClosureCurrent.EventExactificationSectorLaw}. This extends the two-law
detached microscopic package by precisely the intrinsic non-dictionary
sector-generation clauses needed for Step~3. The constructor
\lean{ClosureCurrent.EventExactificationSectorLaw.toMicroscopicSectorLaw}
recovers the previously isolated role-indexed microscopic sector law, so the
current forcing seam then closes by
\lean{ClosureCurrent.EventExactificationSectorLaw.surface}. Concretely, that
surface combines the two-law selected-bundle result of
Theorem~\ref{thm:two-law-detached-microscopic-package-forces-role-indexed-selected-bundle-surface}
with intrinsic sector generation, minimal compatible realizations, and the
canonical logical profile \([3,2,1]\).
\end{proof}

\begin{proof}[Proof of Theorem~\ref{thm:local-event-microscopic-law-recovers-selected-bundle-surface}]
\label{proof:thm:local-event-microscopic-law-recovers-selected-bundle-surface}
Definitions~\ref{def:detached-local-event-closure-current-law} and
\ref{def:detached-local-event-exactification-autonomy-law} are formalized by
\lean{ClosureCurrent.LocalEventCurrentLaw} and
\lean{ClosureCurrent.LocalEventExactificationLaw}. The first local-event law
stores the current-side data in a more primitive form than the earlier
two-law package: the selected visible physical principle, the selected bundle
and its five sectors, the hidden exact exchange current on oriented
pair-comparison slots, the six-slot clause for the first closure-stable bulk
event, the role-indexed visible readout family, and the common
continuum-closure transport law all appear directly there. The constructors
\lean{ClosureCurrent.LocalEventCurrentLaw.toBundleAssemblyLaw},
\lean{ClosureCurrent.LocalEventCurrentLaw.toRoleReadoutFamily}, and
\lean{ClosureCurrent.LocalEventCurrentLaw.toClosureClassTransport} recover the
bundle-first assembly, the role-indexed readout family, and the common
closure-transport law. The constructor
\lean{ClosureCurrent.LocalEventExactificationLaw.toEventExactificationLaw} then
recovers the earlier detached event exactification/autonomy package. Applying
\lean{ClosureCurrent.LocalEventExactificationLaw.surface} yields exactly the
selected-bundle existence, bedrock-carrying generator, selected-cut autonomy,
fiber invariance, and role-indexed carrier closure transport stated in the
theorem.
\end{proof}

\begin{proof}[Proof of Corollary~\ref{cor:local-event-microscopic-sector-law-closes-current-forcing-seam}]
\label{proof:cor:local-event-microscopic-sector-law-closes-current-forcing-seam}
The corollary is formalized by \lean{ClosureCurrent.LocalEventSectorLaw}. This
extends the local-event exactification/autonomy package by the same intrinsic
non-dictionary sector-generation clauses used earlier. The constructor
\lean{ClosureCurrent.LocalEventSectorLaw.toEventExactificationSectorLaw}
recovers the previous detached event exactification / sector package, so the
present forcing seam closes by
\lean{ClosureCurrent.LocalEventSectorLaw.surface}. Concretely, that surface
combines the local-event selected-bundle result of
Theorem~\ref{thm:local-event-microscopic-law-recovers-selected-bundle-surface}
with intrinsic sector generation, minimal compatible realizations, and the
canonical logical profile \([3,2,1]\).
\end{proof}

\begin{proof}[Proof of Theorem~\ref{thm:detached-pair-exchange-current-object-recovers-local-event-law}]
\label{proof:thm:detached-pair-exchange-current-object-recovers-local-event-law}
Definition~\ref{def:detached-pair-exchange-current-object} is formalized by the
detached structures \lean{ClosureCurrent.PairExchangeCarrier},
\lean{ClosureCurrent.PairExchangeCurrent},
\lean{ClosureCurrent.EventExactifier},
\lean{ClosureCurrent.VisibleQuotient}, and
\lean{ClosureCurrent.LocalEventObject}. The first three give an actual local
current object: a carrier with reversal and stock, a local pair current on the
event legs, and a local exactifier with returned defect. The visible quotient
objects supply autonomous visible readouts role by role. The theorem
\lean{ClosureCurrent.LocalEventObject.surface} records those object-level
facts directly. The constructors
\lean{ClosureCurrent.LocalEventObject.toLocalEventCurrentLaw} and
\lean{ClosureCurrent.LocalEventObject.toLocalEventExactificationLaw} then read
the selected bundle, role-indexed visible family, and common continuum closure
transport from that object-level package and recover the detached local-event
exactification/autonomy law. This is exactly the statement of the theorem.
\end{proof}

\begin{proof}[Proof of Corollary~\ref{cor:detached-pair-exchange-current-object-closes-current-forcing-seam}]
\label{proof:cor:detached-pair-exchange-current-object-closes-current-forcing-seam}
The corollary is formalized by \lean{ClosureCurrent.LocalEventSectorObject},
which extends the detached pair-exchange current object by the same intrinsic
non-dictionary sector-generation clauses used in the earlier local-event and
two-law packages. The constructor
\lean{ClosureCurrent.LocalEventSectorObject.toLocalEventSectorLaw} recovers the
detached local-event sector law, so the forcing seam closes by
\lean{ClosureCurrent.LocalEventSectorObject.surface}. Concretely, that surface
combines the object-level current facts of
Theorem~\ref{thm:detached-pair-exchange-current-object-recovers-local-event-law}
with intrinsic selected-bundle existence, intrinsic sector generation, and the
canonical logical profile \([3,2,1]\).
\end{proof}

\begin{proof}[Proof of Theorem~\ref{thm:observer-realized-detached-current-object-derives-selected-bundle-surface}]
\label{proof:thm:observer-realized-detached-current-object-derives-selected-bundle-surface}
Definition~\ref{def:observer-realized-detached-current-object} is formalized by
\lean{ClosureCurrent.VisibleObserverRealization} and
\lean{ClosureCurrent.ObservedCurrentObject}. The key point is that the selected
bundle sectors and role-indexed selected-bundle readout predicates are no
longer primitive fields. The direct visible sectors are derived as nonempty
quotient-value types, and the kinetic sector is derived from stock on a
distinguished local leg by
\lean{ClosureCurrent.ObservedCurrentObject.derivedSectors}. The role-indexed
readout predicates are likewise derived by
\lean{ClosureCurrent.ObservedCurrentObject.readsRoleOnSelectedBundle}, which
feeds the current quotient value for each role into the observer realization.
The constructor \lean{ClosureCurrent.ObservedCurrentObject.toLocalEventObject}
then recovers the earlier detached local-event object, and the packaged
surface theorem \lean{ClosureCurrent.ObservedCurrentObject.surface} yields the
nonempty local pair current, direct/return compatibility, autonomous minimal
visible quotients, and nonempty intrinsic selected-bundle existence claimed in
the theorem.
\end{proof}

\begin{proof}[Proof of Corollary~\ref{cor:observer-realized-detached-current-sector-object-closes-the-reduced-seam}]
\label{proof:cor:observer-realized-detached-current-sector-object-closes-the-reduced-seam}
The corollary is formalized by \lean{ClosureCurrent.ObservedCurrentSectorObject}.
This extends the observer-realized detached current object by the same
intrinsic non-dictionary sector-generation clauses used earlier. The
constructor \lean{ClosureCurrent.ObservedCurrentSectorObject.toLocalEventSectorObject}
recovers the detached local-event sector object, and the packaged surface
theorem \lean{ClosureCurrent.ObservedCurrentSectorObject.surface} then yields
nonempty intrinsic selected-bundle existence, nonempty intrinsic sector
generation, and the canonical logical profile \([3,2,1]\), together with the
object-side current, direct/return, and autonomous-quotient facts inherited
from Theorem~\ref{thm:observer-realized-detached-current-object-derives-selected-bundle-surface}.
\end{proof}

\begin{proof}[Proof of Theorem~\ref{thm:selected-observer-current-anchor-forces-exact-anchored-closure}]
\label{proof:thm:selected-observer-current-anchor-forces-exact-anchored-closure}
Definition~\ref{def:selected-observer-current-anchor} is formalized by
\lean{ClosureCurrent.CurrentAnchor}. This is smaller than the transported
anchored package: it keeps only the hidden current object, the visible quotient
objects, and the continuum-class identity at the selected bridge observer
itself. The packaged theorem \lean{ClosureCurrent.CurrentAnchor.surface}
combines the object-side current facts with the existing Lean PartIV
selected-cut closure lemma \lean{SelectedCut.carrierClosureReadout} at that
one anchored
observer. Concretely, it yields a nonempty local pair current, direct/return
compatibility, autonomous minimal visible quotients, exact visibility on the
selected bridge observer, and the forced continuum singleton on the anchored
continuum class. The auxiliary theorem
\lean{ClosureCurrent.AnchoredCurrentObject.anchorSurface} records that the
later anchored transported package still carries this smaller anchored surface
before its full selected-cut closure transport is used elsewhere.
\end{proof}

\begin{proof}[Proof of Theorem~\ref{thm:anchored-current-transport-extends-exact-anchored-closure-to-the-full-selected-cut}]
\label{proof:thm:anchored-current-transport-extends-exact-anchored-closure-to-the-full-selected-cut}
Definition~\ref{def:anchored-current-transport} is formalized by
\lean{ClosureCurrent.AnchoredCurrentTransport}. This extends the selected-
observer current anchor by one additional datum only: the common selected-cut
closure transport package \lean{ClosureClassTransport}. The packaged theorem
\lean{ClosureCurrent.AnchoredCurrentTransport.surface} therefore retains the
entire anchored closure surface of
Theorem~\ref{thm:selected-observer-current-anchor-forces-exact-anchored-closure}
and adds exactly the quantified closure-transport conclusion across the full
selected cut. The auxiliary theorem
\lean{ClosureCurrent.AnchoredCurrentObject.transportSurface} records that the
later anchored detached current object still carries this transported surface
before its remaining bundle-intrinsicity clauses are used to recover the
observer-realized package.
\end{proof}

\begin{proof}[Proof of Theorem~\ref{thm:anchored-detached-current-object-recovers-the-observer-realized-surface}]
\label{proof:thm:anchored-detached-current-object-recovers-the-observer-realized-surface}
Definition~\ref{def:anchored-detached-current-object} is formalized by
\lean{ClosureCurrent.AnchoredCurrentObject}. The reduced observer side is now
the selected-observer current anchor together with the generic closure-class
transport package already isolated in the detached forcing surface. The
constructor
\lean{ClosureCurrent.AnchoredCurrentObject.toObservedCurrentObject} reads each
visible quotient value at that one selected observer and packages the result as
the earlier predicate-style observer realization. The packaged theorem
\lean{ClosureCurrent.AnchoredCurrentObject.surface} then yields exactly the
same reduced consequences as before: nonempty local pair current,
direct/return compatibility, autonomous minimal visible quotients, and
nonempty intrinsic selected-bundle existence.
\end{proof}

\begin{proof}[Proof of Corollary~\ref{cor:anchored-detached-current-sector-object-closes-the-reduced-seam}]
\label{proof:cor:anchored-detached-current-sector-object-closes-the-reduced-seam}
The corollary is formalized by
\lean{ClosureCurrent.AnchoredCurrentSectorObject}. This extends the anchored
detached current object by the same intrinsic non-dictionary sector-generation
clauses used earlier. The constructor
\lean{ClosureCurrent.AnchoredCurrentSectorObject.toObservedCurrentSectorObject}
recovers the observer-realized sector object, and the packaged theorem
\lean{ClosureCurrent.AnchoredCurrentSectorObject.surface} then yields
nonempty intrinsic selected-bundle existence, nonempty intrinsic sector
generation, and the canonical logical profile \([3,2,1]\), together with the
same current-side direct/return and autonomous-quotient facts.
\end{proof}

\begin{proof}[Proof of Theorem~\ref{thm:origin-readout-law-forces-origin-closure}]
\label{proof:thm:origin-readout-law-forces-origin-closure}
The exact minimal constructive layer is now formalized by
\lean{ClosureCurrent.OriginReadoutLaw}. This structure is indexed by one fixed
\lean{ClosureCurrent.MicroscopicClosureLaw}; it does not try to reconstruct
that law. Instead it records only the constructive selected-anchor data still
missing above it. Those missing data are now split explicitly into two smaller
packages:
\begin{enumerate}
\item \lean{ClosureCurrent.OriginReadoutAnchor}, which keeps one
      selected-anchor current and one visible quotient family;
\item \lean{ClosureCurrent.OriginReadoutClosureTarget}, which asks for one
      common selected-cut closure transport package;
\item \lean{ClosureCurrent.OriginReadoutClauses}, which keeps one common
      grammar package and the two bundle-intrinsicity clauses.
\end{enumerate}

The remaining exact descent targets are therefore
\lean{ClosureCurrent.OriginReadoutAnchorTarget},
\lean{ClosureCurrent.OriginReadoutClosureTarget}, and
\lean{ClosureCurrent.OriginReadoutClausesTarget}. Any actual readout law
realizes them immediately by
\lean{ClosureCurrent.OriginReadoutLaw.anchorTarget},
\lean{ClosureCurrent.OriginReadoutLaw.closureTarget}, and
\lean{ClosureCurrent.OriginReadoutLaw.clausesTarget}. The first two are
reassembled into \lean{ClosureCurrent.OriginReadoutTransportTarget} by
\lean{ClosureCurrent.OriginReadoutTransport.ofParts} together with
\lean{ClosureCurrent.MicroscopicClosureLaw.originReadoutTransportTargetOfParts}.
The full \lean{ClosureCurrent.OriginReadoutTarget} is then recovered from that
transport target and the clauses target by
\lean{ClosureCurrent.OriginReadoutLaw.ofParts} together with the reduction
theorem \lean{ClosureCurrent.MicroscopicClosureLaw.originReadoutTargetOfParts}.
So the first origin seam is no longer one opaque existence claim; it is one
selected-anchor current/quotient target, one selected-cut closure-transport
target, and one common grammar/intrinsicity target.

The resulting selected-bundle/readout witness is built directly by
\lean{ClosureCurrent.OriginReadoutLaw.toOriginClosureWitness}. No exactifier,
no leg-count bookkeeping, and no distinguished leg are used there. The
selected bundle is the coherent bundle already carried by the physical
realization principle; the phase, wave, gauge, and geometric sectors are read
from the visible quotient values of the selected-anchor current; the kinetic
sector is witnessed by stock at the zero current value; the role-indexed
selected-bundle readout is the selected bridge observer itself; and the common
continuum-closure identity is exactly the stored selected-cut closure
transport.

Therefore the general reduction theorem
\lean{ClosureCurrent.MicroscopicClosureLaw.originClosureTargetOfReadout} is
immediate once \lean{ClosureCurrent.OriginReadoutTarget} is met: the required
origin-closure witness already exists. Equivalently, once the two smaller
transport-plus-clauses targets are met, the same conclusion follows through
\lean{ClosureCurrent.MicroscopicClosureLaw.originReadoutTargetSurfaceOfParts}.
This yields exactly the detached selected-bundle forcing surface in the
theorem statement. So the precise current origin seam is now completely
explicit above a fixed microscopic law.
\end{proof}

\begin{proof}[Proof of Theorem~\ref{thm:local-event-exactification-law-reduces-origin-seam-to-transport}]
\label{proof:thm:local-event-exactification-law-reduces-origin-seam-to-transport}
The key point is that the stronger local-event exactification package already
contains every proposition-level clause on the origin side except the
constructive selected-anchor current/quotient data themselves. On the Lean side, the
forgetful map
\lean{ClosureCurrent.LocalEventExactificationLaw.toMicroscopicClosureLaw}
simply drops the selected-bundle/readout fields and retains exactly the
smaller proposition-level microscopic closure law on the same bridge and
selected physical law.

The common grammar and the two bundle-intrinsicity clauses are not lost under
that forgetting step. They are repackaged directly by
\lean{ClosureCurrent.LocalEventExactificationLaw.toOriginReadoutClauses},
which immediately yields the exact
\lean{ClosureCurrent.OriginReadoutClausesTarget} via
\lean{ClosureCurrent.LocalEventExactificationLaw.clausesTarget}. So, once one
starts from the stronger local-event exactification law, the clauses side of
the origin seam is already closed.

The selected-cut closure transport is not lost either. It is already available
from the local-event exactification package through the inherited selected-cut
closure law, and it is repackaged as the exact
\lean{ClosureCurrent.OriginReadoutClosureTarget} by
\lean{ClosureCurrent.LocalEventExactificationLaw.closureTarget}. So above the
stronger local-event exactification law, the closure side of the selected-anchor
transport seam is already closed as well.

The same exactification package also already carries the common grammar and
bundle-intrinsicity clauses needed for the sharper constructive microscopic
layer. On the Lean side this is the map
\lean{ClosureCurrent.LocalEventExactificationLaw.toConstructiveMicroscopicClauses},
which immediately yields the exact
\lean{ClosureCurrent.ConstructiveMicroscopicClausesTarget} by
\lean{ClosureCurrent.LocalEventExactificationLaw.constructiveMicroscopicClausesTarget}.
So if one phrases the remaining origin-side descent through the sharper
constructive split from Theorem~\ref{thm:constructive-microscopic-law-forces-origin-closure},
then the clauses side is still already closed above the stronger
local-event exactification lane.

What remains is therefore exactly the selected-anchor current/quotient target
\lean{ClosureCurrent.OriginReadoutAnchorTarget} for the smaller microscopic law
just extracted. If that target is supplied, then the selected-anchor transport
target follows immediately by
\lean{ClosureCurrent.LocalEventExactificationLaw.originReadoutTransportTargetOfAnchor},
because the closure target is already available. The full origin-readout target
then follows by
\lean{ClosureCurrent.LocalEventExactificationLaw.originReadoutTargetOfAnchor}.
The corresponding origin-closure target is then obtained by
\lean{ClosureCurrent.LocalEventExactificationLaw.originClosureTargetOfAnchor},
and the detached selected-bundle forcing surface follows by
\lean{ClosureCurrent.LocalEventExactificationLaw.originTargetSurfaceOfAnchor}.
So above the stronger local-event exactification lane, the immediate
origin-readout theorem is reduced to one precise constructive datum: the
selected-anchor current/quotient target.

If instead one works with the sharper constructive split, then the same
exactification package already supplies the exact selected-cut closure target
and the exact grammar/intrinsicity target for the induced microscopic law.
The only remaining constructive datum is therefore
\lean{ClosureCurrent.ConstructiveAnchorCurrentTarget}. Once that target is
supplied, the constructive current target follows by
\lean{ClosureCurrent.LocalEventExactificationLaw.constructiveCurrentTargetOfConstructiveAnchorCurrentTarget},
the bundled constructive microscopic target follows by
\lean{ClosureCurrent.LocalEventExactificationLaw.constructiveMicroscopicTargetOfConstructiveAnchorCurrentTarget},
and the same origin-readout, origin-closure, and forcing-surface conclusions
follow directly by
\lean{ClosureCurrent.LocalEventExactificationLaw.originReadoutTargetOfConstructiveAnchorCurrentTarget},
\lean{ClosureCurrent.LocalEventExactificationLaw.originClosureTargetOfConstructiveAnchorCurrentTarget},
and
\lean{ClosureCurrent.LocalEventExactificationLaw.originTargetSurfaceOfConstructiveAnchorCurrentTarget}.
So above the stronger exactification lane, the sharper constructive origin
seam is now reduced to the constructive anchor-current target alone.
\end{proof}

\begin{proof}[Proof of Theorem~\ref{thm:constructive-microscopic-law-forces-origin-closure}]
\label{proof:thm:constructive-microscopic-law-forces-origin-closure}
The canonical-origin law is the first theorem-bearing microscopic package in
which the origin data are internal rather than externally supplied. It fixes,
on one selected datum, a selected anchor, the visible quotient/readout data,
the exactifier, the coherent selected-cut closure transport, and the common
grammar together with the remaining bundle-intrinsicity clauses. It is smaller
than the anchored detached current object because the latter contains one
further distinguished-leg choice that is irrelevant to the present theorem.

Forgetting the extra structure does not change the underlying microscopic law:
the same hidden pair-exchange current, conservation law, relabeling symmetry,
local losslessness, autonomous visible readout, and returned-defect clause are
already determined by the selected-anchor current data carried by the
canonical-origin package itself. The constructive law then separates naturally
into three exact pieces.

First, it carries the anchor-current realization: the selected anchor and the
visible quotient family are already fixed on the chosen datum. Second, it
carries the selected-cut closure transport: the closure identity holding at
the selected anchor is transported coherently across the full selected cut.
Third, it carries the common grammar and bundle-intrinsicity clauses needed to
read the visible sectors on that same datum. The first two pieces assemble the
selected-anchor transport law, and adjoining the clauses package assembles the
full origin-readout law.

Once that origin-readout law is present, the selected-bundle/readout witness is
explicit. The selected bundle is the coherent bundle already carried by the
law; the phase, wave, gauge, and geometric sectors are read from the visible
quotients on the selected anchor; the kinetic role is read from stock at zero
current; the role-indexed readout is taken on the same selected datum; and the
common closure identity is exactly the transported selected-cut closure law.
Hence the origin-closure target follows, and with it the detached
selected-bundle forcing surface.

The same three exact pieces also recover the bundled constructive microscopic
target, so that bundled target adds no additional mathematical content beyond
the anchor-current realization, the selected-cut closure transport, and the
common grammar/intrinsicity clauses. Therefore the origin seam has genuinely
collapsed: once the law itself carries canonical origin realization, origin
closure and the selected-bundle forcing surface are forced with no further
origin-side witness input.
\end{proof}

\begin{proof}[Proof of Corollary~\ref{cor:constructive-sector-law-forces-sector-closure}]
\label{proof:cor:constructive-sector-law-forces-sector-closure}
The corollary is formalized by \lean{ClosureCurrent.ConstructiveSectorLaw}.
This is now a genuinely smaller self-contained sector law: it extends the
constructive microscopic law by exactly
\lean{SelectedBundle.CanonicalGeneration},
\lean{generatedFromIntrinsicChannelAlgebra}, and
\lean{minimalCompatibleRealizationsOnly}, and nothing else. The stronger
anchored current sector object forgets to it through
\lean{ClosureCurrent.AnchoredCurrentSectorObject.toConstructiveSectorLaw}, so
no distinguished leg remains visible at this sector seam.

The induced map
\lean{ClosureCurrent.ConstructiveSectorLaw.toConstructiveMicroscopicLaw} then
forgets back to the constructive microscopic law from the previous theorem. So
the underlying origin-closure theorem is inherited unchanged from
Theorem~\ref{thm:constructive-microscopic-law-forces-origin-closure}; on the
Lean side this is
\lean{ClosureCurrent.ConstructiveSectorLaw.originClosureTarget}. The intrinsic
sector witness is then read directly by
\lean{ClosureCurrent.ConstructiveSectorLaw.toSectorOriginWitness}, so the
sector-closure target itself is forced by
\lean{ClosureCurrent.ConstructiveSectorLaw.sectorClosureTarget}. Applying the
general target-to-surface theorem yields
\lean{ClosureCurrent.ConstructiveSectorLaw.sectorTargetSurface}, which is
exactly the detached Step~3 forcing surface stated in the corollary.

Again the same reduction is now exposed at the exact law level. The target
\lean{ClosureCurrent.ConstructiveSectorTarget} says only that one fixed
proposition-level microscopic law admits one such self-contained constructive
sector realization on the same selected datum. Any actual constructive sector
law realizes that target by \lean{ClosureCurrent.ConstructiveSectorLaw.target}.
The abstract reductions
\lean{ClosureCurrent.MicroscopicClosureLaw.constructiveMicroscopicTargetOfConstructiveSectorTarget},
\lean{ClosureCurrent.MicroscopicClosureLaw.originClosureTargetOfConstructiveSectorTarget},
\lean{ClosureCurrent.MicroscopicClosureLaw.sectorClosureTargetOfConstructiveSectorTarget},
and
\lean{ClosureCurrent.MicroscopicClosureLaw.sectorTargetSurfaceOfConstructiveSectorTarget}
then show that once this exact sector target is met, the constructive
microscopic target, the origin-closure target, the sector-closure target, and
the detached Step~3 forcing surface all follow formally. So the remaining
law-level Step~3 descent target is now isolated exactly as the constructive
sector target itself.
\end{proof}

\begin{proof}[Proof of Theorem~\ref{thm:autonomy-forces-the-detached-flagship-equation}]
\label{proof:thm:autonomy-forces-the-detached-flagship-equation}
Definition~\ref{def:autonomous-realized-detached-current-object} is formalized
by \lean{ClosureCurrent.AutonomousStateField} and
\lean{ClosureCurrent.AutonomousAnchoredCurrentObject}. The compact state field
stores one realization of the hidden pair current as a selected-cut state,
together with one role-indexed realization of visible quotient values as
projection-fixed states and one autonomous state law on the same selected cut
for each visible role. The expanded current-state and per-role visible-law
objects are recovered by
\lean{ClosureCurrent.AutonomousStateField.toCurrentState} and
\lean{ClosureCurrent.AutonomousStateField.toVisibleStateLaw}. The key forcing
step is the theorem \lean{ClosureCurrent.VisibleStateLaw.transportCompatible}:
for each visible role, the autonomy factorization on the selected cut and the
realized exactifier/readout/evolution identities imply the boxed equation
\[
C_r(\Gamma j)=T_r(C_r(j)).
\]
The residual identity \(\Delta_r(j)=C_r(R_e(j))\) remains the same
returned-defect readout used in the earlier local and anchored law forms. The
pointwise agreement
\[
L_r(s_r(v))=X(s_r(v))
\]
is the theorem \lean{ClosureCurrent.VisibleStateLaw.law_eq_generator_on_realized},
obtained from the selected-cut autonomy uniqueness lemma already isolated in
the detached bridge surface. The packaged theorem
\lean{ClosureCurrent.AutonomousAnchoredCurrentObject.surface} combines these
autonomy-forced channel equations with the anchored exact-visibility and
selected-cut closure-transport conclusions already carried by the underlying
anchored detached current object.
\end{proof}

\begin{proof}[Proof of Theorem~\ref{thm:autonomy-collapses-to-a-unified-state-field}]
\label{proof:thm:autonomy-collapses-to-a-unified-state-field}
Theorem~\ref{thm:autonomy-forces-the-detached-flagship-equation} still carries a
role-indexed family of autonomous visible laws \(L_r\). Theorem
\ref{thm:autonomy-collapses-to-a-unified-state-field} is the statement that
this extra indexing is not genuinely new dynamics. On the Lean side this is
formalized by
\lean{ClosureCurrent.AutonomousAnchoredCurrentObject.unifiedSurface}. The key
constructor is
\lean{ClosureCurrent.AutonomousStateField.toUnifiedStateField}: for each visible
role, the theorem
\lean{ClosureCurrent.VisibleStateLaw.law_eq_generator_on_realized} identifies
the role-indexed law with the same matched selected generator on realized
visible states. This yields the unified package
\lean{ClosureCurrent.UnifiedAnchoredCurrentObject}, whose surface theorem then
supplies the same flagship equation and selected-cut closure conclusions with a
single common ambient generator.
\end{proof}

\begin{proof}[Proof of Theorem~\ref{thm:unified-state-field-forces-the-detached-flagship-equation}]
\label{proof:thm:unified-state-field-forces-the-detached-flagship-equation}
Definition~\ref{def:unified-detached-state-field} is formalized by
\lean{ClosureCurrent.UnifiedStateField} and
\lean{ClosureCurrent.UnifiedAnchoredCurrentObject}. This is a genuine reduction
of the previous autonomy package: the role-indexed visible laws are no longer
stored separately. Instead every visible role is realized on the same ambient
selected-cut state space and transported there by the same matched selected
generator. The constructor
\lean{ClosureCurrent.UnifiedStateField.toAutonomousStateField} recovers the
earlier role-indexed autonomous package by taking each visible law to be the
same matched selected generator. The key direct theorem is
\lean{ClosureCurrent.UnifiedStateField.transportCompatible}, which proves the
boxed equation
\[
C_r(\Gamma j)=T_r(C_r(j))
\]
from the current realization, the quotient readout factorization, and the
selected-cut autonomy law of the matched selected generator itself. The pointwise
common-generator statement
\[ 
X(s_r(v))=s_r(T_r(v))
\]
is the theorem \lean{ClosureCurrent.UnifiedStateField.generator_on_visible}. The
packaged selected-cut conclusion is then exactly
\lean{ClosureCurrent.UnifiedAnchoredCurrentObject.surface}.
\end{proof}

\begin{proof}[Proof of Theorem~\ref{thm:explicit-current-state-assembly-forces-the-unified-flagship-surface}]
\label{proof:thm:explicit-current-state-assembly-forces-the-unified-flagship-surface}
Definition~\ref{def:explicit-detached-current-state-assembly} is formalized by
\lean{ClosureCurrent.SignedExchangeReadout},
\lean{ClosureCurrent.FourLegFrame},
\lean{ClosureCurrent.CurrentStateAssembly},
\lean{ClosureCurrent.AssembledUnifiedStateField}, and
\lean{ClosureCurrent.AssembledUnifiedAnchoredCurrentObject}. The new content is
that the hidden current realization is no longer stored as an arbitrary map
\[
\rho : \texttt{PairExchangeCurrent}\to\texttt{State}.
\]
Instead, \lean{ClosureCurrent.assembledState} defines it explicitly from the
pair current itself: a signed readout on carrier values is applied to the
current incident from one framed leg to each of the four framed legs, and the
four resulting signed counts are packaged as one state in \texttt{State}. The
theorem \lean{ClosureCurrent.CurrentStateAssembly.surface} records exactly this
explicit formula together with the remaining nontrivial compatibility clause,
namely that the local exactifier intertwines this assembled state with the
matched selected generator.

The constructor
\lean{ClosureCurrent.AssembledUnifiedStateField.toUnifiedStateField} then
reads the earlier unified detached state field from the explicit assembly,
without re-introducing any opaque hidden-state map. Finally,
\lean{ClosureCurrent.AssembledUnifiedAnchoredCurrentObject.surface} is obtained
by transporting the existing theorem
\lean{ClosureCurrent.UnifiedAnchoredCurrentObject.surface} across that
constructor. Hence the same unified selected-cut flagship surface holds, but
now with a concrete hidden current-state formula rather than a black-box
realization.
\end{proof}

\begin{proof}[Proof of Theorem~\ref{thm:explicit-current-assembly-is-balanced}]
\label{proof:thm:explicit-current-assembly-is-balanced}
The statement is formalized first by
\lean{ClosureCurrent.assembledState_total_zero}. Write the four framed legs as
\(\lambda_0,\lambda_1,\lambda_2,\lambda_3\). By construction, the total
assembled state is the sum of the four leg-flux values
\[
\Phi_{\lambda_i}(j_e)
=
\sum_{k=0}^3 \chi_e(j_e(\lambda_i,\lambda_k)).
\]
The diagonal terms vanish because the pair current vanishes on the diagonal.
Every off-diagonal ordered term appears exactly twice in the total sum:
once as \(\chi_e(j_e(\lambda_i,\lambda_k))\) and once as
\(\chi_e(j_e(\lambda_k,\lambda_i))\). The pair-current law gives
\[
j_e(\lambda_k,\lambda_i)=\rho_e(j_e(\lambda_i,\lambda_k)),
\]
and the signed readout law gives
\[
\chi_e(\rho_e x)=-\chi_e(x).
\]
Thus each off-diagonal pair contributes a signed count together with its
negation, and the total signed sum is zero. The packaged selected-cut version
is \lean{ClosureCurrent.AssembledUnifiedAnchoredCurrentObject.currentBalance},
which applies the same cancellation theorem to every hidden current state
produced by the explicit assembled anchored object.
\end{proof}

\begin{proof}[Proof of Corollary~\ref{thm:explicit-current-assembly-factors-through-balanced-sector}]
\label{proof:thm:explicit-current-assembly-factors-through-balanced-sector}
The factorization is formalized first by
\lean{ClosureCurrent.CurrentStateAssembly.realizeBalanced}: the explicit current
assembly is upgraded from a map
\[
\rho_{\mathrm{asm}} : \texttt{PairExchangeCurrent}\to\texttt{State}
\]
to a map
\[
\rho_{\mathrm{bal}} : \texttt{PairExchangeCurrent}\to\texttt{BalancedState},
\]
where \texttt{BalancedState} packages a state together with the certificate that
its total signed content vanishes. The theorem
\lean{ClosureCurrent.CurrentStateAssembly.balancedSurface} then records two
facts: first, forgetting \(\rho_{\mathrm{bal}}(j_e)\) back to
\texttt{State} recovers exactly the assembled state \(\rho_{\mathrm{asm}}(j_e)\);
second, the exactifier/generator intertwining law still holds after this
forgetful map. The packaged anchored-object versions are
\lean{ClosureCurrent.AssembledUnifiedAnchoredCurrentObject.currentBalancedState}
and
\lean{ClosureCurrent.AssembledUnifiedAnchoredCurrentObject.currentBalancedSurface}.
So the detached hidden current is not merely pointwise balanced; it genuinely
takes values in the balanced zero-sum sector.
\end{proof}

\begin{proof}[Proof of Corollary~\ref{thm:explicit-hidden-current-law-closes-on-the-balanced-sector}]
\label{proof:thm:explicit-hidden-current-law-closes-on-the-balanced-sector}
The factorization corollary already identifies the realized hidden current with
a balanced-state-valued field
\[
\rho_{\mathrm{bal}} : \texttt{PairExchangeCurrent}\to\texttt{BalancedState}
\]
whose forgetful image is the explicit assembled state. Apply this to the
exactified current \(\Gamma j_e\). Since
\(\rho_{\mathrm{bal}}(\Gamma j_e)\in\texttt{BalancedState}\), its forgetful
state has total signed content zero. But
\lean{ClosureCurrent.CurrentStateAssembly.balancedSurface} identifies that
forgetful state with
\[
X(\rho_{\mathrm{bal}}(j_e)_{\mathrm{forget}}).
\]
Hence the matched selected generator preserves balance on every realized hidden
current state. This is formalized directly by
\lean{ClosureCurrent.CurrentStateAssembly.generator_preserves_balance}, and the
packaged selected-cut form is
\lean{ClosureCurrent.AssembledUnifiedAnchoredCurrentObject.currentBalancedClosure}.
\end{proof}

\begin{proof}[Proof of Corollary~\ref{thm:explicit-hidden-current-reduces-to-three-coordinates}]
\label{proof:thm:explicit-hidden-current-reduces-to-three-coordinates}
The reduced carrier itself is formalized by
\lean{ClosureCurrent.BalancedCoordinates}. Its reconstruction map
\lean{ClosureCurrent.BalancedCoordinates.toBalancedState} sends a triple
\((x_0,x_1,x_2)\) to the balanced four-state
\((x_0,x_1,x_2,-(x_0+x_1+x_2))\). The theorem
\lean{ClosureCurrent.BalancedCoordinates.toBalancedState_ofBalancedState}
shows that this loses no information on the balanced hidden sector, while
\lean{ClosureCurrent.BalancedCoordinates.ofBalancedState_toBalancedState} shows
that the three-coordinate carrier reconstructs exactly.

For the current law, the reduction is expressed by
\lean{ClosureCurrent.CurrentStateAssembly.realizeCoordinates} together with
\lean{ClosureCurrent.CurrentStateAssembly.coordinateSurface}. The latter states
that the balanced hidden current state is exactly the reconstruction of the
three visible reduced coordinates, and that the exactifier/generator law still
holds after passing through this reduced carrier and reconstructing. The
selected-cut packaged form is
\lean{ClosureCurrent.AssembledUnifiedAnchoredCurrentObject.currentCoordinates}
together with
\lean{ClosureCurrent.AssembledUnifiedAnchoredCurrentObject.currentCoordinateSurface}.
\end{proof}

\begin{proof}[Proof of Corollary~\ref{thm:reduced-hidden-current-transport-law}]
\label{proof:thm:reduced-hidden-current-transport-law}
The reduced transport operator is formalized by
\lean{ClosureCurrent.BalancedCoordinates.transport}: given a reduced balanced
coordinate package, it reconstructs the balanced four-state, applies the ambient
matched selected generator, and projects the result back to the first three
coordinates. The exact current-side law is then
\lean{ClosureCurrent.CurrentStateAssembly.coordinateTransport}, which proves
\[
\rho_3(\Gamma j_e)=\Pi_3\!\bigl(X(\iota(\rho_3(j_e)))\bigr)
\]
for the explicit detached current assembly. The packaged selected-cut version is
\lean{ClosureCurrent.AssembledUnifiedAnchoredCurrentObject.currentCoordinateTransport}.
So the detached hidden current now has a genuine reduced three-coordinate
transport law, not merely a balanced-state factorization.
\end{proof}

\begin{proof}[Proof of Corollary~\ref{thm:reduced-hidden-current-discrete-trajectory-law}]
\label{proof:thm:reduced-hidden-current-discrete-trajectory-law}
The two discrete iterates are formalized by
\lean{ClosureCurrent.AnchoredCurrentObject.exactifyIterate} on the hidden pair
current and \lean{ClosureCurrent.CurrentStateAssembly.coordinateIterate} on the
reduced balanced carrier. Both are defined by the same recursive endomap
iteration scheme.

The key step is the one-step transport theorem
\lean{ClosureCurrent.CurrentStateAssembly.coordinateTransport}. The finite-step
statement then follows by induction on \(n\): the base case is the identity
iterate, and the inductive step rewrites one exactification step by the
one-step reduced transport law and then applies the induction hypothesis.
This is exactly
\lean{ClosureCurrent.CurrentStateAssembly.coordinateTrajectory}. The packaged
selected-cut form is
\lean{ClosureCurrent.AssembledUnifiedAnchoredCurrentObject.currentCoordinateTrajectory}.
\end{proof}

\begin{proof}[Proof of Corollary~\ref{thm:selected-observed-state-discrete-trajectory-law}]
\label{proof:thm:selected-observed-state-discrete-trajectory-law}
The selected observed state is formalized by
\lean{ClosureCurrent.CurrentStateAssembly.observedState}:
\[
\rho_{\mathrm{sel}}(j_e)
=
\Pi_{\mathrm{sel}}(\rho(j_e)).
\]
The one-step observed transport identity is
\lean{ClosureCurrent.CurrentStateAssembly.observedTransport}. Its proof uses
only three ingredients:
\[
\rho(\Gamma j_e)=X(\rho(j_e)),
\qquad
\Pi_{\mathrm{sel}}(Xx)=X(\Pi_{\mathrm{sel}}x),
\qquad
X(\Pi_{\mathrm{sel}}x)\in\operatorname{im}(\Pi_{\mathrm{sel}}),
\]
namely the explicit exactifier/generator assembly law together with autonomy of
the matched selected generator on the selected cut.

The finite-step statement then follows by induction on \(n\), exactly as in the
reduced hidden-current case, and is formalized by
\lean{ClosureCurrent.CurrentStateAssembly.observedTrajectory}. The packaged
selected-cut form is
\lean{ClosureCurrent.AssembledUnifiedAnchoredCurrentObject.currentObservedTrajectory}.
\end{proof}

\begin{proof}[Proof of Corollary~\ref{thm:detached-law-terms-already-have-concrete-coordinate-expressions}]
\label{proof:thm:detached-law-terms-already-have-concrete-coordinate-expressions}
The reduced hidden current formula is literally the definition proved in
\lean{ClosureCurrent.CurrentStateAssembly.coordinateFormula}: the three reduced
coordinates are exactly the first three framed leg-flux sums of the explicit
assembled current state. The packaged selected-cut form is
\lean{ClosureCurrent.AssembledUnifiedAnchoredCurrentObject.currentCoordinateFormula}.

Likewise, the selected observed state formula is literally
\lean{ClosureCurrent.CurrentStateAssembly.observedFormula}: it is the selected
projection of the concrete assembled four-leg flux state. The packaged
selected-cut form is
\lean{ClosureCurrent.AssembledUnifiedAnchoredCurrentObject.currentObservedFormula}.

So at this stage the law terms themselves are already concrete once the
detached assembly datum is fixed. The only remaining abstraction is upstream:
the signed carrier readout, the genuine four-leg frame, and the matched
selected generator are still carried as detached data rather than yet derived
from the bare pair-current/event law.
\end{proof}

\begin{proof}[Proof of Corollary~\ref{thm:selected-observed-trajectory-is-read-directly-from-reduced-hidden-coordinates}]
\label{proof:thm:selected-observed-trajectory-is-read-directly-from-reduced-hidden-coordinates}
The basic readout identity is
\lean{ClosureCurrent.CurrentStateAssembly.observedFromCoordinates}. It says
that on every realized detached current state,
\[
\rho_{\mathrm{sel}}(j_e)
=
\Pi_{\mathrm{sel}}(\iota(\rho_3(j_e))).
\]
This is stronger than merely having both a hidden reduced trajectory and an
observed trajectory: it identifies the observed state as the exact selected
readout of the reduced hidden coordinates.

The finite-step version is then
\lean{ClosureCurrent.CurrentStateAssembly.observedCoordinateTrajectory}. Its
proof is immediate from two earlier ingredients:
\[
\rho_3(\Gamma^{[n]} j_e)=T_3^{[n]}(\rho_3(j_e)),
\qquad
\rho_{\mathrm{sel}}(j_e)=\Pi_{\mathrm{sel}}(\iota(\rho_3(j_e))).
\]
Substituting the reduced hidden trajectory into the readout identity yields the
claimed exact observed trajectory formula. The packaged selected-cut forms are
\lean{ClosureCurrent.AssembledUnifiedAnchoredCurrentObject.currentObservedFromCoordinates}
and
\lean{ClosureCurrent.AssembledUnifiedAnchoredCurrentObject.currentObservedCoordinateTrajectory}.
\end{proof}

\begin{proof}[Proof of Corollary~\ref{thm:detached-current-trajectories-factor-through-the-reduced-observed-system}]
\label{proof:thm:detached-current-trajectories-factor-through-the-reduced-observed-system}
The system object itself is formalized by
\lean{ClosureCurrent.ReducedObservedSystem}. It keeps exactly two maps on the
reduced hidden carrier: the hidden step map and the selected observed readout.
The canonical detached instance is
\lean{ClosureCurrent.AnchoredCurrentObject.reducedObservedSystem}, whose hidden
step is \(T_3\) and whose observed readout is \(O_{\mathrm{sel}}\).

The hidden trajectory factorization is
\lean{ClosureCurrent.CurrentStateAssembly.systemHiddenTrajectory}, which is just
the reduced hidden trajectory theorem restated as a trajectory of the system
hidden step map. The observed trajectory factorization is
\lean{ClosureCurrent.CurrentStateAssembly.systemObservedTrajectory}, which uses
the earlier theorem that the observed trajectory is read directly from the
reduced hidden coordinates. The packaged selected-cut forms are
\lean{ClosureCurrent.AssembledUnifiedAnchoredCurrentObject.reducedObservedSystem},
\lean{ClosureCurrent.AssembledUnifiedAnchoredCurrentObject.systemHiddenTrajectory},
and
\lean{ClosureCurrent.AssembledUnifiedAnchoredCurrentObject.systemObservedTrajectory}.
\end{proof}

\begin{proof}[Proof of Corollary~\ref{thm:reduced-hidden-coordinates-are-already-the-autonomous-detached-state-variable}]
\label{proof:thm:reduced-hidden-coordinates-are-already-the-autonomous-detached-state-variable}
The key point is that the previous corollary already factors both hidden and
observed trajectories through one reduced observed system on
\texttt{BalancedCoordinates}. Once that factorization is available, equal
initial reduced coordinates force equal iterated hidden states simply because
the same iterate of the same endomap is being applied to equal starting points.
This is formalized by the generic system congruence theorem
\lean{ClosureCurrent.ReducedObservedSystem.hiddenIterate_congr} together with
the current-side consequence
\lean{ClosureCurrent.CurrentStateAssembly.sameCoordinates_sameHiddenTrajectory}.

The observed statement is identical after composing with the canonical selected
observed readout of the same reduced system. The generic congruence theorem is
\lean{ClosureCurrent.ReducedObservedSystem.observedIterate_congr}, and the
detached current-side consequence is
\lean{ClosureCurrent.CurrentStateAssembly.sameCoordinates_sameObservedTrajectory}.
The packaged selected-cut forms are
\lean{ClosureCurrent.AssembledUnifiedAnchoredCurrentObject.sameCoordinates_sameHiddenTrajectory}
and
\lean{ClosureCurrent.AssembledUnifiedAnchoredCurrentObject.sameCoordinates_sameObservedTrajectory}.
Thus the reduced hidden coordinates already determine both the hidden and the
selected observed discrete trajectories.
\end{proof}

\begin{proof}[Proof of Corollary~\ref{thm:explicit-detached-assemblies-matter-dynamically-only-through-the-reduced-initial-state}]
\label{proof:thm:explicit-detached-assemblies-matter-dynamically-only-through-the-reduced-initial-state}
Fix one anchored detached current object \(data\) and two explicit current
assemblies \(A_1,A_2 : \texttt{CurrentStateAssembly}\;data\). The crucial point
is that the canonical reduced observed system
\lean{ClosureCurrent.AnchoredCurrentObject.reducedObservedSystem} depends only
on \(data\), not on the particular assembly. So both assemblies factor their
hidden and observed trajectories through the same reduced hidden step map and
the same canonical selected observed readout.

If the two assemblies assign the same reduced hidden coordinates to the chosen
initial currents, then the same iterate of that common reduced hidden step map
is being applied to equal starting points. Hence the iterated hidden reduced
states agree. This is formalized by
\lean{ClosureCurrent.CurrentStateAssembly.sameCoordinates_acrossAssemblies_sameHiddenTrajectory}.
Applying the common observed readout to those equal reduced hidden states then
gives equal observed states, formalized by
\lean{ClosureCurrent.CurrentStateAssembly.sameCoordinates_acrossAssemblies_sameObservedTrajectory}.
Thus the explicit assembly data contribute no additional dynamical memory once
the reduced hidden initial state is fixed.
\end{proof}

\begin{proof}[Proof of Corollary~\ref{thm:the-reduced-detached-hidden-state-assembly-already-carries-the-detached-dynamics}]
\label{proof:thm:the-reduced-detached-hidden-state-assembly-already-carries-the-detached-dynamics}
The reduced package is formalized by \lean{ClosureCurrent.ReducedCurrentAssembly}.
It keeps only two ingredients: the reduced hidden-state map
\[
\rho_3 : \texttt{PairExchangeCurrent}\to\texttt{BalancedCoordinates}
\]
and the exact one-step transport identity
\[
\rho_3(\Gamma j_e)=T_3(\rho_3(j_e)).
\]
Every explicit detached current assembly determines such a reduced package by
\lean{ClosureCurrent.CurrentStateAssembly.toReducedCurrentAssembly}; this is
just the passage from the full readout/frame realization to the reduced hidden
coordinates and their transport law.

Once that smaller package is present, the finite-step hidden trajectory law no
longer needs the rest of the explicit assembly. It is proved directly by
induction from the one-step transport law as
\lean{ClosureCurrent.ReducedCurrentAssembly.coordinateTrajectory}. The hidden
and observed factorization through the canonical reduced observed system are
then
\lean{ClosureCurrent.ReducedCurrentAssembly.systemHiddenTrajectory} and
\lean{ClosureCurrent.ReducedCurrentAssembly.systemObservedTrajectory}. The
packaged selected-cut forms are
\lean{ClosureCurrent.AssembledUnifiedAnchoredCurrentObject.toReducedCurrentAssembly},
\lean{ClosureCurrent.AssembledUnifiedAnchoredCurrentObject.reducedAssemblyHiddenTrajectory},
and
\lean{ClosureCurrent.AssembledUnifiedAnchoredCurrentObject.reducedAssemblyObservedTrajectory}.
So at this point the full explicit assembly is upstream data only; the
detached discrete dynamics themselves already live on the smaller reduced
hidden-state assembly.
\end{proof}

\begin{proof}[Proof of Corollary~\ref{thm:the-reduced-anchored-detached-current-object-already-carries-the-canonical-reduced-observed-system}]
\label{proof:thm:the-reduced-anchored-detached-current-object-already-carries-the-canonical-reduced-observed-system}
Definition~\ref{def:reduced-anchored-detached-current-object} is formalized by
\lean{ClosureCurrent.ReducedAnchoredCurrentObject}. It keeps exactly two
pieces of data: the anchored detached current bridge and one reduced hidden
state assembly on that bridge.

Every explicit assembled detached current object determines such a smaller
object by
\lean{ClosureCurrent.AssembledUnifiedAnchoredCurrentObject.toReducedAnchoredCurrentObject}.
On that smaller object, the canonical reduced observed system is already the
definition
\lean{ClosureCurrent.ReducedAnchoredCurrentObject.reducedObservedSystem}. Its
iterated reduced hidden transport and selected observed readout are
\lean{ClosureCurrent.ReducedAnchoredCurrentObject.currentCoordinateIterate}
and
\lean{ClosureCurrent.ReducedAnchoredCurrentObject.currentObservedState}. The
exact finite-step hidden trajectory law on the reduced hidden carrier is
\lean{ClosureCurrent.ReducedAnchoredCurrentObject.currentCoordinateTrajectory},
and the hidden and observed current trajectories as trajectories of the same
canonical reduced observed system are
\lean{ClosureCurrent.ReducedAnchoredCurrentObject.systemHiddenTrajectory} and
\lean{ClosureCurrent.ReducedAnchoredCurrentObject.systemObservedTrajectory}.

So after passage to the reduced anchored object, the detached dynamics are
already completely carried by one anchored current bridge together with the
reduced hidden-state realization; the larger explicit assembled state-field
package is not needed to state those dynamics any longer.
\end{proof}

\begin{proof}[Proof of Corollary~\ref{thm:at-the-reduced-hidden-state-level-the-initial-reduced-coordinate-is-the-full-detached-state}]
\label{proof:thm:at-the-reduced-hidden-state-level-the-initial-reduced-coordinate-is-the-full-detached-state}
Fix one anchored detached current bridge and one reduced hidden-state assembly
on it. By the reduced trajectory factorization theorems,
\lean{ClosureCurrent.ReducedCurrentAssembly.systemHiddenTrajectory} and
\lean{ClosureCurrent.ReducedCurrentAssembly.systemObservedTrajectory}, every
hidden and observed future is obtained by applying the canonical reduced
observed system of that anchored bridge to the initial reduced hidden
coordinate.

So if two initial currents have the same reduced hidden coordinate, then the
same hidden iterate is being applied to the same starting point, and the same
observed readout is then being taken. This is formalized by
\lean{ClosureCurrent.ReducedCurrentAssembly.sameCoordinates_sameHiddenTrajectory}
and
\lean{ClosureCurrent.ReducedCurrentAssembly.sameCoordinates_sameObservedTrajectory}.

The same reasoning works across two reduced hidden-state assemblies on the same
anchored bridge, because both still factor through that same canonical reduced
observed system. If the initial reduced coordinates agree, then the future
hidden and observed trajectories agree as well. This is formalized by
\lean{ClosureCurrent.ReducedCurrentAssembly.sameCoordinates_acrossAssemblies_sameHiddenTrajectory}
and
\lean{ClosureCurrent.ReducedCurrentAssembly.sameCoordinates_acrossAssemblies_sameObservedTrajectory}.
Thus on the reduced detached surface, the reduced initial hidden coordinate is
already the entire dynamical state variable.
\end{proof}

\begin{proof}[Proof of Corollary~\ref{thm:current-free-state-space-form-of-the-detached-dynamics}]
\label{proof:thm:current-free-state-space-form-of-the-detached-dynamics}
On a fixed anchored detached current bridge, the canonical reduced observed
system already determines a current-free hidden trajectory from each initial
reduced hidden coordinate and a current-free observed trajectory from that same
coordinate. These are formalized by
\lean{ClosureCurrent.AnchoredCurrentObject.hiddenTrajectoryFrom}
and
\lean{ClosureCurrent.AnchoredCurrentObject.observedTrajectoryFrom}.

Now fix a reduced hidden-state assembly and an initial detached current. The
hidden and observed trajectories of that current are exactly the hidden and
observed trajectories of the canonical reduced observed system started from the
initial reduced hidden coordinate of that current. This is formalized by
\lean{ClosureCurrent.ReducedCurrentAssembly.hiddenTrajectory_from_initialCoordinate}
and
\lean{ClosureCurrent.ReducedCurrentAssembly.observedTrajectory_from_initialCoordinate}.

The same statement holds in packaged form on the smaller reduced anchored
current object, where the current-free hidden and observed state-space orbits
are
\lean{ClosureCurrent.ReducedAnchoredCurrentObject.hiddenTrajectoryFrom}
and
\lean{ClosureCurrent.ReducedAnchoredCurrentObject.observedTrajectoryFrom},
and the exact factorization of detached current trajectories through those
orbits is
\lean{ClosureCurrent.ReducedAnchoredCurrentObject.hiddenTrajectory_from_initialCoordinate}
and
\lean{ClosureCurrent.ReducedAnchoredCurrentObject.observedTrajectory_from_initialCoordinate}.
So after passage to the reduced hidden coordinate, the detached law is already
an honest discrete state-space system, with the hidden current entering only
through the initial state.
\end{proof}

\begin{proof}[Proof of Corollary~\ref{thm:current-free-recurrence-form-on-the-reduced-hidden-carrier}]
\label{proof:thm:current-free-recurrence-form-on-the-reduced-hidden-carrier}
The hidden current-free state-space orbit on one anchored bridge is generated
stepwise by the canonical reduced hidden update map. This is formalized by
\lean{ClosureCurrent.AnchoredCurrentObject.hiddenTrajectoryFrom_step}. The
selected observed orbit is then read from that hidden orbit by the canonical
selected projection, formalized by
\lean{ClosureCurrent.AnchoredCurrentObject.observedTrajectoryFrom_readout}.

The same recurrence and readout laws hold on the smaller reduced anchored
current object, where they are formalized by
\lean{ClosureCurrent.ReducedAnchoredCurrentObject.hiddenTrajectoryFrom_step}
and
\lean{ClosureCurrent.ReducedAnchoredCurrentObject.observedTrajectoryFrom_readout}.
So the detached law on the reduced hidden carrier is already a current-free
discrete recurrence together with a current-free selected readout.
\end{proof}

\begin{proof}[Proof of Corollary~\ref{thm:the-detached-law-is-a-reduced-initial-value-problem}]
\label{proof:thm:the-detached-law-is-a-reduced-initial-value-problem}
Definition~\ref{def:reduced-detached-initial-state} is formalized by
\lean{ClosureCurrent.ReducedInitialState}. It is exactly the usual initial-data
package for the reduced detached state-space system: one reduced anchored
current object together with one initial reduced hidden coordinate.

Its hidden and observed orbits are
\lean{ClosureCurrent.ReducedInitialState.hiddenTrajectory}
and
\lean{ClosureCurrent.ReducedInitialState.observedTrajectory}. The initial
condition is
\lean{ClosureCurrent.ReducedInitialState.hiddenTrajectory_zero}; the recurrence
law is
\lean{ClosureCurrent.ReducedInitialState.hiddenTrajectory_step}; and the
selected readout law is
\lean{ClosureCurrent.ReducedInitialState.observedTrajectory_readout}.

Every detached current on the reduced anchored system determines such initial
data by
\lean{ClosureCurrent.ReducedAnchoredCurrentObject.initialStateOf}. The exact
identification of the hidden and observed detached current futures with the
corresponding initial-state orbits is then
\lean{ClosureCurrent.ReducedAnchoredCurrentObject.hiddenTrajectory_initialStateOf}
and
\lean{ClosureCurrent.ReducedAnchoredCurrentObject.observedTrajectory_initialStateOf}.
So after passing to the reduced hidden coordinate, the detached microscopic law
is already an honest discrete initial-value problem.
\end{proof}

\begin{proof}[Proof of Corollary~\ref{thm:the-detached-law-already-factors-through-the-canonical-reduced-observed-system-plus-initial-coordinate}]
\label{proof:thm:the-detached-law-already-factors-through-the-canonical-reduced-observed-system-plus-initial-coordinate}
Definition~\ref{def:current-free-reduced-system-state} is formalized by
\lean{ClosureCurrent.ReducedObservedDynamicsState}. This is the smaller pure
current-free initial-data package: it keeps only the canonical reduced hidden
update map, the canonical selected observed readout, and the initial reduced
hidden coordinate.

Its hidden and observed orbits are
\lean{ClosureCurrent.ReducedObservedDynamicsState.hiddenTrajectory}
and
\lean{ClosureCurrent.ReducedObservedDynamicsState.observedTrajectory}. The
initial condition, recurrence law, and readout law are
\lean{ClosureCurrent.ReducedObservedDynamicsState.hiddenTrajectory_zero},
\lean{ClosureCurrent.ReducedObservedDynamicsState.hiddenTrajectory_step},
and
\lean{ClosureCurrent.ReducedObservedDynamicsState.observedTrajectory_readout},
all built from the underlying pure dynamics laws
\lean{ClosureCurrent.ReducedObservedDynamics.hiddenTrajectoryFrom_step}
and
\lean{ClosureCurrent.ReducedObservedDynamics.observedTrajectoryFrom_readout}.

For an actual detached current on a reduced anchored system, the corresponding
smaller initial data are given directly by
\lean{ClosureCurrent.ReducedAnchoredCurrentObject.dynamicsStateOf}. The exact
hidden and observed detached futures are then the hidden and observed orbits of
that smaller current-free initial state, formalized by
\lean{ClosureCurrent.ReducedAnchoredCurrentObject.hiddenTrajectory_dynamicsStateOf}
and
\lean{ClosureCurrent.ReducedAnchoredCurrentObject.observedTrajectory_dynamicsStateOf}.
So even the reduced anchored current object is more structure than the detached
dynamics themselves need once the canonical reduced observed system is fixed.
\end{proof}

\begin{proof}[Proof of Corollary~\ref{thm:once-rho3-is-fixed-the-detached-future-is-generated-by-the-selected-bridge}]
\label{proof:thm:once-rho3-is-fixed-the-detached-future-is-generated-by-the-selected-bridge}
Definition~\ref{def:selected-bridge-reduced-state} is formalized by
\lean{ClosureCurrent.SelectedReducedState}. This is the smaller bridge-driven
package: it retains only the selected bridge generator, the selected
projection, and the initial reduced hidden coordinate.

Its hidden and observed orbits are
\lean{ClosureCurrent.SelectedReducedState.hiddenTrajectory}
and
\lean{ClosureCurrent.SelectedReducedState.observedTrajectory}. The initial
condition, recurrence law, and readout law are
\lean{ClosureCurrent.SelectedReducedState.hiddenTrajectory_zero},
\lean{ClosureCurrent.SelectedReducedState.hiddenTrajectory_step},
and
\lean{ClosureCurrent.SelectedReducedState.observedTrajectory_readout}.

For an actual detached current on a reduced anchored system, the corresponding
smaller bridge-driven state is
\lean{ClosureCurrent.ReducedAnchoredCurrentObject.selectedReducedStateOf}. The
exact hidden and observed detached futures are then the hidden and observed
orbits of that selected-bridge reduced state, formalized by
\lean{ClosureCurrent.ReducedAnchoredCurrentObject.hiddenTrajectory_selectedReducedStateOf}
and
\lean{ClosureCurrent.ReducedAnchoredCurrentObject.observedTrajectory_selectedReducedStateOf}.
So once \(\rho_3(j_e)\) is fixed, the detached future is generated entirely by
the selected bridge generator and the selected projection.
\end{proof}

\begin{proof}[Proof of Corollary~\ref{thm:the-detached-future-already-depends-only-on-the-coordinate-assembly-and-the-selected-bridge}]
\label{proof:thm:the-detached-future-already-depends-only-on-the-coordinate-assembly-and-the-selected-bridge}
Definition~\ref{def:current-coordinate-assembly} is formalized by
\lean{ClosureCurrent.CurrentCoordinateAssembly} and
\lean{ClosureCurrent.CoordinateAnchoredCurrentObject}. This is the smaller
explicit package that retains only the signed readout, the genuine four-leg
frame, and the reduced transport law for the concrete coordinate assembly
\(\rho_3\).

For every detached current on that package, the corresponding selected-bridge
reduced state is
\lean{ClosureCurrent.CoordinateAnchoredCurrentObject.selectedReducedStateOf}.
The exact hidden and observed detached futures are then the hidden and observed
orbits of that selected-bridge reduced state, formalized by
\lean{ClosureCurrent.CoordinateAnchoredCurrentObject.hiddenTrajectory_selectedReducedStateOf}
and
\lean{ClosureCurrent.CoordinateAnchoredCurrentObject.observedTrajectory_selectedReducedStateOf}.
So after passing from the full explicit current-state assembly to the concrete
reduced-coordinate assembly, the detached future already depends only on that
coordinate assembly and the selected bridge.
\end{proof}

\begin{proof}[Proof of Corollary~\ref{cor:the-local-event-current-object-needs-only-the-coordinate-bridge-to-generate-the-detached-future}]
\label{proof:cor:the-local-event-current-object-needs-only-the-coordinate-bridge-to-generate-the-detached-future}
Definition~\ref{def:local-event-coordinate-bridge} is formalized by
\lean{ClosureCurrent.LocalEventCoordinateObject}. This is the smallest present
upstream microscopic package above the detached local-event current object: it
adds only the signed readout, the genuine four-leg frame, and the reduced
transport identity for the resulting concrete coordinate \(\rho_3\).

The bridge to the smaller anchored coordinate package is
\lean{ClosureCurrent.LocalEventCoordinateObject.toCoordinateAnchoredCurrentObject},
and the corresponding selected-bridge reduced state is
\lean{ClosureCurrent.LocalEventCoordinateObject.selectedReducedStateOf}. The
exact hidden and observed detached futures are then the hidden and observed
orbits of that selected-bridge reduced state, formalized by
\lean{ClosureCurrent.LocalEventCoordinateObject.hiddenTrajectory_selectedReducedStateOf}
and
\lean{ClosureCurrent.LocalEventCoordinateObject.observedTrajectory_selectedReducedStateOf}.
So once the detached local-event current object is fixed, the only remaining
upstream microscopic datum needed to generate the detached future is the
coordinate bridge for \(\rho_3\).
\end{proof}

\begin{proof}[Proof of Theorem~\ref{thm:the-coordinate-bridge-is-derived-from-the-local-event-state-bridge}]
\label{proof:thm:the-coordinate-bridge-is-derived-from-the-local-event-state-bridge}
Definition~\ref{def:local-event-state-bridge} is formalized by
\lean{ClosureCurrent.LocalEventStateBridge}. This strengthens the local-event
current object by keeping the same signed readout and genuine four-leg frame,
but requiring the exactifier/generator intertwining law already on the full
assembled four-state.

On that stronger package, the reduced transport identity for \(\rho_3\) is
derived by \lean{ClosureCurrent.LocalEventStateBridge.coordinateTransport}, so
the coordinate bridge is no longer a separate assumption. The resulting
local-event coordinate bridge is
\lean{ClosureCurrent.LocalEventStateBridge.toLocalEventCoordinateObject}, and
the packaged theorem \lean{ClosureCurrent.LocalEventStateBridge.surface}
recovers the same reduced hidden and observed initial-value dynamics as on the
coordinate-bridge surface.
\end{proof}

\begin{proof}[Proof of Corollary~\ref{cor:the-sharpened-detached-microscopic-law-is-already-a-pure-reduced-initial-value-system}]
\label{proof:cor:the-sharpened-detached-microscopic-law-is-already-a-pure-reduced-initial-value-system}
Once the local-event state bridge is included in the assumed detached
microscopic law, Theorem~\ref{thm:the-coordinate-bridge-is-derived-from-the-local-event-state-bridge}
already yields the local-event coordinate bridge through
\lean{ClosureCurrent.LocalEventStateBridge.toLocalEventCoordinateObject}.

The direct reduced observed dynamics is
\lean{ClosureCurrent.LocalEventStateBridge.reducedDynamics}, and the
corresponding reduced initial-value state for a detached current \(j_e\) is
\lean{ClosureCurrent.LocalEventStateBridge.initialStateOf}. The exact hidden
and observed trajectory factorization through that current-free initial-value
problem is given by
\lean{ClosureCurrent.LocalEventStateBridge.hiddenTrajectory_initialStateOf} and
\lean{ClosureCurrent.LocalEventStateBridge.observedTrajectory_initialStateOf}.
The packaged theorem \lean{ClosureCurrent.LocalEventStateBridge.surface}
compresses the same conclusion. So after sharpening the assumed detached
microscopic law by the local-event state bridge, no further microscopic datum
remains inside the detached lane: the law already determines a pure reduced
initial-value system.
\end{proof}

\begin{proof}[Proof of Theorem~\ref{thm:the-local-event-state-bridge-already-determines-a-pure-four-state-law}]
\label{proof:thm:the-local-event-state-bridge-already-determines-a-pure-four-state-law}
The current-free four-state dynamics itself is formalized by
\lean{ClosureCurrent.StateDynamics} and its initial-value package
\lean{ClosureCurrent.StateDynamicsState}. On the stronger microscopic package
\lean{ClosureCurrent.LocalEventStateBridge}, the corresponding update law and
initial assembled state are exported directly as
\lean{ClosureCurrent.LocalEventStateBridge.stateDynamics} and
\lean{ClosureCurrent.LocalEventStateBridge.stateInitialStateOf}.

The exact factorization of the detached four-state current future through that
current-free initial-value law is
\lean{ClosureCurrent.LocalEventStateBridge.stateTrajectory_initialStateOf}. The
fact that every such four-state orbit remains in the balanced zero-sum sector
is
\lean{ClosureCurrent.LocalEventStateBridge.stateTrajectory_balanced}. The
packaged theorem \lean{ClosureCurrent.LocalEventStateBridge.stateSurface}
combines these statements with the same nonempty current, direct/return
compatibility, and autonomous minimal visible quotient facts already carried by
the local-event state bridge.
\end{proof}

\begin{proof}[Proof of Corollary~\ref{cor:the-reduced-rho3-law-is-the-balanced-quotient-of-the-four-state-law}]
\label{proof:cor:the-reduced-rho3-law-is-the-balanced-quotient-of-the-four-state-law}
On the stronger microscopic package
\lean{ClosureCurrent.LocalEventStateBridge}, the reduced hidden trajectory is
formalized as the exact projection statement
\lean{ClosureCurrent.LocalEventStateBridge.reducedTrajectory_projectState}: the
reduced hidden orbit equals the first-three-coordinate quotient of the pure
four-state orbit.

Likewise the reduced observed future is formalized by
\lean{ClosureCurrent.LocalEventStateBridge.observedTrajectory_stateProjection}:
the observed orbit is exactly the selected projection of that same four-state
trajectory. So the sharpened \(\rho_3\)-law is already derived from one
current-free four-state law and is not an independent dynamical stage.
\end{proof}

\begin{proof}[Proof of Theorem~\ref{thm:the-no-stage-detached-primitive-is-already-a-four-state-law}]
\label{proof:thm:the-no-stage-detached-primitive-is-already-a-four-state-law}
The no-stage package is formalized by
\lean{ClosureCurrent.LocalEventFourStateLaw}. By definition it consists of a
detached local-event object, a signed readout, a genuine four-leg frame, and
the explicit four-state transport identity. Its conversion to the stronger
state-bridge package is
\lean{ClosureCurrent.LocalEventFourStateLaw.toLocalEventStateBridge}.

The packaged theorem \lean{ClosureCurrent.LocalEventFourStateLaw.surface} then
recovers the whole four-state and balanced-quotient surface at once: the pure
four-state initial-value law, the zero-sum balance of every four-state orbit,
the reduced hidden trajectory as the balanced quotient of that orbit, and the
reduced observed trajectory as the selected projection of the same orbit. So if
one wishes to stop before the continuum limit with no further hidden
microscopic stage, the local-event four-state law is already the clean
primitive form.
\end{proof}

\begin{proof}[Proof of Theorem~\ref{thm:the-no-stage-four-state-law-already-forces-a-four-state-continuum-law}]
\label{proof:thm:the-no-stage-four-state-law-already-forces-a-four-state-continuum-law}
The generic continuum package for a pure four-state update map is formalized by
\lean{ClosureCurrent.StateDynamics.refinementCompatibleFamily} and
\lean{ClosureCurrent.StateDynamics.reconstructionDatum}. The resulting family is
constant in the refinement index: every stage carries the same four-state map,
the comparison maps are identities, and the reconstruction maps are identities.

Accordingly the reconstructed observed law at stage \(n\) is already the same
four-state update map \(X\), so asymptotic consistency is exact and immediate.
This is formalized by
\lean{ClosureCurrent.StateDynamics.asymptoticConsistency}. The corresponding
continuum-bridge uniqueness theorem is
\lean{ClosureCurrent.StateDynamics.continuumUniqueness}, and the resulting
singleton forcing statement is
\lean{ClosureCurrent.StateDynamics.continuumLaw_forced}.

Applying that generic four-state continuum construction to the no-stage
microscopic four-state dynamics
\lean{ClosureCurrent.LocalEventFourStateLaw.stateDynamics} yields the detached
continuum forcing theorem
\lean{ClosureCurrent.LocalEventFourStateLaw.stateContinuumLaw_forced}. So the
continuum law is already forced by the discrete no-stage four-state law, not
chosen afterward as an endpoint representative.
\end{proof}

\begin{proof}[Proof of Corollary~\ref{cor:the-reduced-hidden-continuum-law-is-the-balanced-quotient-of-the-forced-four-state-continuum-law}]
\label{proof:cor:the-reduced-hidden-continuum-law-is-the-balanced-quotient-of-the-forced-four-state-continuum-law}
The reduced hidden continuum package is formalized generically by
\lean{ClosureCurrent.ReducedObservedDynamics.hiddenContinuumLaw_forced}. On the
no-stage microscopic four-state law it becomes
\lean{ClosureCurrent.LocalEventFourStateLaw.reducedContinuumLaw_forced}.

The relation between the four-state continuum law and the reduced hidden
continuum law is not extra data. It is exactly the balanced-quotient formula
\lean{ClosureCurrent.LocalEventFourStateLaw.reducedContinuumLaw_quotient},
which states that the reduced hidden continuum map is obtained by applying the
forced four-state continuum law to the balanced four-state reconstruction and
then projecting back to the first three coordinates.

Likewise the visible readout is not a second law. The selected observed
continuum readout is exactly
\lean{ClosureCurrent.LocalEventFourStateLaw.observedContinuumReadout}, i.e. the
selected projection of the same balanced hidden carrier. The packaged theorem
\lean{ClosureCurrent.LocalEventFourStateLaw.continuumSurface} combines these
four-state and reduced-hidden continuum facts with the already established
nonempty current, direct/return compatibility, autonomy, and minimal visible
quotient surfaces.
\end{proof}

\begin{proof}[Proof of Theorem~\ref{thm:the-no-stage-four-state-law-already-has-a-canonical-first-order-continuum-representative}]
\label{proof:thm:the-no-stage-four-state-law-already-has-a-canonical-first-order-continuum-representative}
The package
\lean{ClosureCurrent.ContinuumRepresentative} is the first honest continuum
representative class extracted from the detached no-stage law. It stores only:
\begin{enumerate}
\item the four-state continuum law;
\item the reduced hidden continuum law;
\item the selected observed readout;
\item the forcing statements for the four-state and reduced-hidden continuum
      laws;
\item the three structural identities saying that the four-state law is the
      matched selected generator, the reduced hidden law is its balanced
      transport quotient, and the observed law is its selected projection.
\end{enumerate}

On the no-stage four-state law this package is constructed by
\lean{ClosureCurrent.LocalEventFourStateLaw.toContinuumRepresentative}. The
point is that no extra analytic choice is made there. The four-state continuum
law is literally the matched selected generator on the explicit four-state
carrier, and its forcing statement is obtained directly from
\lean{ClosureCurrent.LocalEventFourStateLaw.stateContinuumLaw_forced}. The
reduced hidden law is literally the balanced transport
\lean{BalancedCoordinates.transport} of that same generator, and its forcing
statement comes directly from
\lean{ClosureCurrent.LocalEventFourStateLaw.reducedContinuumLaw_forced}. The
selected observed law is literally the selected projection of the same
balanced carrier.

All of those identities are packaged together by
\lean{ClosureCurrent.LocalEventFourStateLaw.representativeSurface}. So the
current formal surface already forces one first-order continuum representative
class; any later PDE or variational realization would be a representation of
this class, not an additional law.
\end{proof}

\begin{proof}[Proof of Theorem~\ref{thm:the-canonical-continuum-representative-is-already-additive}]
\label{proof:thm:the-canonical-continuum-representative-is-already-additive}
The extra content here is algebraic, not analytic. The reduced hidden carrier
is first equipped with its own additive operator interface
\lean{ClosureCurrent.CoordinateMap}. This is not a new continuum law; it is
just the balanced hidden carrier together with the two structural equations
that define an additive endomap. Any additive ambient state operator then
induces such a reduced hidden operator by reconstruction and projection via
\lean{ClosureCurrent.CoordinateMap.ofAddMap}. The key input for that
construction is that the balanced hidden carrier now has explicit coordinate
addition and zero, and the transport of an additive ambient operator respects
both of them.

The operator-level representative package itself is
\lean{ClosureCurrent.OperatorRepresentative}. Relative to the earlier
first-order representative, the only genuinely new fields are:
\begin{enumerate}
\item the four-state representative is stored as an additive map on
      \(\mathrm{State}\), not merely as a function;
\item the reduced hidden quotient is stored as an additive map on the balanced
      hidden carrier;
\item the forcing surface is still recorded on the underlying functions, so no
      extra continuum-identification axiom is introduced.
\end{enumerate}

On the no-stage four-state law this stronger package is constructed by
\lean{ClosureCurrent.LocalEventFourStateLaw.toOperatorRepresentative}. The
four-state operator is exactly the selected observed operator carried by the
active selection datum, which already lands in the certified additive-map
surface. The reduced hidden
operator is exactly the induced balanced transport of that same operator. The
packaged theorem
\lean{ClosureCurrent.LocalEventFourStateLaw.operatorRepresentativeSurface}
then states:
\begin{enumerate}
\item the four-state operator law is forced;
\item the reduced hidden operator law is forced;
\item both operators preserve addition and zero;
\item the four-state operator is still the matched selected generator;
\item the reduced hidden operator is still the balanced quotient of that same
      four-state operator;
\item the observed readout is still the selected projection.
\end{enumerate}

So the detached lane has now reached the strongest honest representative class
currently forced by the formal surface: not yet a PDE and not yet a
variational model, but already an additive four-state operator together with
its additive balanced hidden quotient.
\end{proof}

\begin{proof}[Proof of Theorem~\ref{thm:the-canonical-continuum-representative-is-already-quadratic-form-driven}]
\label{proof:thm:the-canonical-continuum-representative-is-already-quadratic-form-driven}
This step strengthens the representative in the variational direction, but only
as far as the current formal surface actually goes. The relevant point is that
the selected observed operator is not introduced externally at this stage. It
already comes from the compiler datum carried by the selected observer, namely
as the operator of the selected admissible quadratic form.

The package formalizing this stronger representative is
\lean{ClosureCurrent.QuadraticRepresentative}. It stores:
\begin{enumerate}
\item an admissible quadratic form on the explicit four-state carrier;
\item the forcing statement for the induced four-state operator;
\item the forcing statement for the balanced hidden quotient operator;
\item the identity of the four-state form with the selected observed law;
\item the identity of the induced operator with the matched selected
      generator;
\item the quotient and observed-projection identities already established
      earlier.
\end{enumerate}

On the no-stage microscopic law this package is built by
\lean{ClosureCurrent.LocalEventFourStateLaw.toQuadraticRepresentative}. The
selected quadratic form is literally
\lean{selected_observed_law} on the active selection datum, so its
admissibility is inherited directly from the derivation step carried in that
datum. No new optimization problem, positivity postulate, or external
variational principle is introduced here.

The theorem
\lean{ClosureCurrent.LocalEventFourStateLaw.quadraticRepresentativeSurface}
packages the result: the forced four-state continuum operator is already the
operator of an admissible quadratic form, the balanced hidden law is still its
quotient, and the observed law is still the selected projection. So the
current surface already reaches a genuine quadratic representative of the
detached law, while honestly stopping short of a full Lagrangian or
Hamiltonian continuum theory.
\end{proof}

\begin{proof}[Proof of Theorem~\ref{thm:the-canonical-quadratic-representative-is-already-psd-symmetry-compatible-and-minimal}]
\label{proof:thm:the-canonical-quadratic-representative-is-already-psd-symmetry-compatible-and-minimal}
The stronger package is formalized by
\lean{ClosureCurrent.MinimalQuadraticRepresentative}. It extends the earlier
quadratic representative by adding the selected finite symmetry datum together
with the exact compiler-side canonicality statements already present in the
selection surface.

On the no-stage four-state law it is built by
\lean{ClosureCurrent.LocalEventFourStateLaw.toMinimalQuadraticRepresentative}.
Nothing new is assumed. The selected quadratic form is still the selected
observed law on the same compiler derivation step, so positive semidefiniteness
and absolute minimality among positive semidefinite competitors come directly
from the admissibility certificate already stored on that derivation step.

The selected finite symmetry datum and its representation are likewise already
part of the same selection datum. Equivariance of the selected quadratic form
is therefore inherited directly from the certified compiler field
\lean{derivation_equivariant}. Once that is fixed, the equivariant minimality
statement is the direct compiler theorem
\lean{Compiler.symmetry_reduction}.

The packaged theorem
\lean{ClosureCurrent.LocalEventFourStateLaw.minimalQuadraticRepresentativeSurface}
collects these four statements into one surface. So the detached lane now has
not only an admissible quadratic representative, but a PSD, symmetry-compatible,
minimal selected quadratic law in the exact compiler sense.
\end{proof}

\begin{proof}[Proof of Theorem~\ref{thm:the-canonical-quadratic-representative-already-lives-on-the-least-motion-stationary-cut}]
\label{proof:thm:the-canonical-quadratic-representative-already-lives-on-the-least-motion-stationary-cut}
The stationary selected-cut package is formalized by
\lean{ClosureCurrent.StationaryQuadraticRepresentativeSurface}. On the
detached law it is produced directly by
\lean{ClosureCurrent.LocalEventFourStateLaw.stationaryQuadraticRepresentativeSurface}.

Nothing new is assumed. The same selected datum already contains:
\begin{enumerate}
\item the finite candidate family and minimizer certificate;
\item the selected projection;
\item the selected objective gradient;
\item the stationarity witness for that selected projection.
\end{enumerate}
So the least-motion part of the theorem is immediate from the explicit
selection surface: \lean{observerMotionValue_minimal} gives the minimizing
objective inequality, while the selected Schur bridge identifies the selected
projection both with the chosen finite minimizer and with the realized horizon
cut.

The stationarity clause then comes from the projector calculus alone.
\lean{stationarity_condition} converts the stored stationarity witness into the
vanishing projection commutator of the selected gradient, and
\lean{stationary_implies_offdiag_zero} immediately turns that into vanishing
\(PQ\) and \(QP\) gradient blocks.

The law-side clause is separate but equally direct. Exact visibility of the
selected quadratic law on the realized horizon cut is already supplied by the
selected Schur bridge. Rewriting along the equality between the selected
projection and the realized horizon cut gives exact visibility on the selected
projector itself. The target calculus then gives:
\begin{enumerate}
\item \lean{exact_visible_implies_commutes} for commutation with the cut;
\item \lean{commutes_implies_zeroPQ} for vanishing \(PQ\) block;
\item \lean{commutes_implies_zeroQP} for vanishing \(QP\) block.
\end{enumerate}
So the selected quadratic representative is already realized on one exact
least-motion stationary cut. This is stronger than mere admissibility or
minimality, but it still stops well short of a full PDE or action-principle
realization.
\end{proof}

\begin{proof}[Proof of Theorem~\ref{thm:the-canonical-quadratic-representative-is-already-the-exact-effective-law-on-the-least-motion-cut}]
\label{proof:thm:the-canonical-quadratic-representative-is-already-the-exact-effective-law-on-the-least-motion-cut}
This is where the rebuilt Schur/Feshbach calculus closes the next seam. The
previous theorem already placed the selected quadratic law on one exact
least-motion stationary cut and, crucially, made it exact-visible there. Once
that is in hand, the effective-law comparison is no longer an extra ansatz.

The formal package is
\lean{ClosureCurrent.EffectiveQuadraticRepresentativeSurface}, proved on the
detached no-stage law by
\lean{ClosureCurrent.LocalEventFourStateLaw.effectiveQuadraticRepresentativeSurface}.
Its proof is short because it only uses the target calculus already rebuilt on
the active spine.

First, exact visibility of the selected quadratic law on the chosen cut feeds
directly into \lean{exact_visible_implies_effective}. So for the canonical
selected shadow propagator, the selected quadratic operator already agrees
pointwise with its Schur/Feshbach effective observable law.

Second, the stationary-cut theorem already supplied vanishing \(PQ\)
transport on that same cut. By \lean{zeroPQ_implies_perfect}, this gives the
perfect-coherence condition needed to collapse the Schur correction itself.
Hence the same proof package identifies the Schur complement with the zero
operator on that cut.

Finally, \lean{returnedResidual_eq_schur} identifies the returned residual of
the canonical residual-return field with that same Schur correction. Since the
Schur correction is zero, the returned residual vanishes identically as well.
So the current exact-selected surface has already eliminated any further
effective-law correction layer: the selected quadratic law is already its own
exact observable effective law on the least-motion cut.
\end{proof}

\begin{proof}[Proof of Theorem~\ref{thm:the-exact-effective-selected-law-already-is-the-matched-selected-generator}]
\label{proof:thm:the-exact-effective-selected-law-already-is-the-matched-selected-generator}
This theorem packages the point at which the detached analytic lane already
has one common selected governing operator. The formal statement is
\lean{ClosureCurrent.LocalEventFourStateLaw.selectedGoverningSurface}, and
the governing residual itself is
\lean{ClosureCurrent.governingScalarResidual}.

Two equalities are fused there. First, the previous effective-law theorem
already gives pointwise equality between the selected quadratic law and the
exact effective selected law on the least-motion cut. Second, the earlier
quadratic representative surface already gives pointwise equality between the
selected quadratic law and the matched selected generator. So the detached
lane already satisfies
\[
A_{\mathrm{sel}} = A_{\mathrm{eff},\mathrm{sel}} = G_{\mathrm{sel}}
\]
on the same four-state carrier.

The residual
\[
r^{\mathrm{gov}}_{\mathrm{sel}}(x)=
\operatorname{SignedCount.ofNat}\!\bigl(\langle x,G_{\mathrm{sel}}x\rangle\bigr)
\]
is then just the rebuilt state pairing read against that matched selected
generator. Since the same generator already agrees pointwise with both the
selected quadratic law and the exact effective selected law, the same scalar
residual can be read from all three operators. So the theorem does not add a
new operator. It only makes explicit that the strongest current selected-cut
surface already carries one common governing operator before any PDE or
action-principle realization is attempted.
\end{proof}

\begin{proof}[Proof of Theorem~\ref{thm:the-selected-governing-operator-already-carries-an-exact-continuous-effective-flow-package}]
\label{proof:thm:the-selected-governing-operator-already-carries-an-exact-continuous-effective-flow-package}
This theorem is the first honest field/flow packaging step after the common
selected governing operator has been identified. The formal package is
\lean{ClosureCurrent.LocalEventFourStateLaw.selectedGoverningFlowSurface}.

The point is not to claim a PDE. The point is to show that the existing
continuum-bridge flow interfaces already apply without any extra analytic
hypothesis.
The selected projection and shadow are read as a constant interpolation datum
on the fixed least-motion cut. The block-derivative datum is then the literal
\(PP/PQ/QP/QQ\) block decomposition of the selected quadratic law on that same
cut.

The earlier selected-cut results already do the real work. Exact visibility
and commutation force the \(PQ\) and \(QP\) blocks to vanish. Exact visible
support also kills the \(QQ\) block, because the selected quadratic law
annihilates pure shadow input. Meanwhile the \(PP\) block is already the
matched selected generator by the governing-operator theorem.

From there the exact effective selected law is packaged as a
\lean{ContinuousEffectiveFlowData} object. Its flow is the exact effective
selected law itself, and the previous theorem already identifies that flow
pointwise with the matched selected generator. So the theorem adds no new
operator and no new analytic assumption. It simply makes explicit that the
selected governing operator already carries an exact continuum-bridge-style
field/flow package on the least-motion cut.
\end{proof}

\begin{proof}[Proof of Theorem~\ref{thm:the-selected-field-flow-package-admits-no-exogenous-constitutive-completion}]
\label{proof:thm:the-selected-field-flow-package-admits-no-exogenous-constitutive-completion}
This theorem is the constitutive counterpart of the selected field/flow
package. The formal statement is
\lean{ClosureCurrent.LocalEventFourStateLaw.selectedPromotionSurface}.

The selected promotion context is built in the most honest way available. Its
projector is the selected cut itself, its carried transport is the exact
effective selected law, its reconstruction is the identity on the four-state
carrier, and its only apparent constitutive term is the returned residual of
that same exact effective selected law.

But the exact effective-law theorem has already proved that returned residual
to vanish identically on the least-motion cut. So the promoted constitutive
term is zero. The promoted observed law therefore collapses to the same forced
selected continuum law, because only the carried transport term survives.

Once that is in hand, the generic continuum-bridge promotion wrappers add no further
freedom. The characteristic-scale promotion fields are witnessed directly by
equality with the same selected continuum law, the boundary identity of the
carried transport, and exact vanishing of the transported constitutive term.
The finite admissible promotion chain then becomes constant, and admissibility
forces every stage to select the same law. Thus the minimum-energy branch is
unique at every finite stage.

So the theorem does not create a new constitutive layer. It proves the
opposite: on the exact selected cut, apparent constitutive content is already
nothing but returned residual from hidden-sector elimination, and here even
that residual vanishes.
\end{proof}

\begin{proof}[Proof of Theorem~\ref{thm:the-scalarized-variational-residual-is-already-the-residual-of-the-exact-effective-law}]
\label{proof:thm:the-scalarized-variational-residual-is-already-the-residual-of-the-exact-effective-law}
This theorem closes the gap between the exact effective-law surface and the
later scalarized variational surface. The key point is that no new analytic
datum is inserted between them.

The formal statement is
\lean{ClosureCurrent.LocalEventFourStateLaw.effectiveVariationalRepresentativeSurface},
and the effective residual itself is
\lean{ClosureCurrent.effectiveScalarResidual}. By definition, that residual is
the rebuilt state pairing read directly against the exact effective observable
law on the least-motion cut.

But the previous theorem already proved that the selected quadratic law is
exactly that effective observable law on the same cut. So the new effective
residual agrees pointwise with the earlier scalar residual
\lean{ClosureCurrent.scalarResidual}. The same variational theorem used later,
\lean{ClosureCurrent.LocalEventFourStateLaw.variationalRepresentativeSurface},
already identifies that scalar residual with the matched selected-generator
residual and already supplies the dynamics triad, observer triad, and
Hamilton--Jacobi projection bridge.

Therefore the detached lane has not produced two different scalarizations. The
scalarized variational residual is already the residual of the exact effective
law itself, and the same regular/observer/Hamilton--Jacobi triad belongs to
that exact effective law on the selected cut.
\end{proof}

\begin{proof}[Proof of Theorem~\ref{thm:the-canonical-continuum-representative-already-carries-an-exact-scalarized-variational-triad}]
\label{proof:thm:the-canonical-continuum-representative-already-carries-an-exact-scalarized-variational-triad}
This is the first point at which the detached continuum representative reaches
an honestly variational presentation. The crucial boundary is that nothing new
is postulated beyond the already forced quadratic representative. Instead, one
simply reads a signed-count scalar residual from the selected quadratic
operator by the rebuilt state pairing.

That scalar residual is formalized as
\lean{ClosureCurrent.scalarResidual}. By definition it sends a four-state
carrier value \(x\) to the signed-count scalar
\[
\operatorname{SignedCount.ofNat}\!\bigl(\langle x,Q_{\mathrm{sel}}^{\mathrm{op}}x\rangle\bigr),
\]
where \(Q_{\mathrm{sel}}\) is the selected admissible quadratic form already
forced in the previous theorem. The first important point is that this is not
a new operator or a new action. The theorem
\lean{ClosureCurrent.LocalEventFourStateLaw.variationalRepresentativeSurface}
proves that the same scalar residual is equally the quadratic-form residual
and the matched-generator residual, because the quadratic operator is already
the selected generator on the detached continuum lane.

From that same scalar residual the file builds three explicit data:
\begin{enumerate}
\item the regular variational system
      \lean{ClosureCurrent.regularSystem};
\item the observer-side conservative transport datum
      \lean{ClosureCurrent.observerTransport};
\item the Hamilton--Jacobi projection bridge
      \lean{ClosureCurrent.projectionDatum}.
\end{enumerate}
Each one is built without any additional analytic structure. The Lagrangian,
Hamiltonian, and characteristic residuals in the first package are literally
the same scalar residual function. The Eulerian, Lagrangian, and
characteristic residuals in the second package are again literally the same
scalar residual function. The projection bridge is then immediate because the
observer-side characteristic residual and the dynamics-side characteristic
residual are definitionally equal.

The theorem
\lean{ClosureCurrent.LocalEventFourStateLaw.variationalRepresentativeSurface}
packages exactly the three foundation theorems that now apply:
\begin{enumerate}
\item \lean{dynamics_triad} for the regular variational datum;
\item \lean{observer_triad} for the conservative transport datum;
\item \lean{hj_projection} for the Hamilton--Jacobi bridge datum.
\end{enumerate}
So the detached no-stage law already carries a precise scalarized
variational/transport triad. But the theorem is careful not to say more than
that: there is still no full analytic action functional, no Hamiltonian flow
machinery, and no imported continuum mechanics layer. The gain is structural
rather than analytic.
\end{proof}

\begin{proof}[Proof of Theorem~\ref{thm:the-detached-analytic-lane-already-collapses-to-one-selected-operator-residual-law}]
\label{proof:thm:the-detached-analytic-lane-already-collapses-to-one-selected-operator-residual-law}
This theorem records the cleanest current law-shaped collapse of the detached
analytic lane. The selected-cut law object is formalized by
\lean{ClosureCurrent.SelectedFieldLaw}, and its detached instance is
\lean{ClosureCurrent.LocalEventFourStateLaw.selectedFieldLaw}. The actual
surface theorem is
\lean{ClosureCurrent.LocalEventFourStateLaw.selectedFieldLawSurface}.

The operator component of this selected law is simply the matched selected
generator \(G_{\mathrm{sel}}\). The first clause of the theorem says that the
forced four-state continuum law on the least-motion cut is already exactly
that operator. So the detached continuum dynamics has already collapsed to one
selected operator before any later flagship specialization is attempted.

The residual component is the governing signed-count scalar
\[
r^{\mathrm{gov}}_{\mathrm{sel}}(x)=
\operatorname{SignedCount.ofNat}\!\bigl(\langle x,G_{\mathrm{sel}}x\rangle\bigr).
\]
The same surface theorem identifies both earlier scalar readouts with this one
residual: the scalar residual previously read from the canonical quadratic law
and the scalar residual read from the exact effective selected law are both
already equal to \(r^{\mathrm{gov}}_{\mathrm{sel}}\).

Once that is in place, the remaining analytic layers become generic wrappers.
The selected law object canonically builds:
\begin{enumerate}
\item a regular variational system from the same residual;
\item an observer-side conservative transport datum from the same residual;
\item a Hamilton--Jacobi projection bridge from those two data;
\item a zero-constitutive promotion context from the selected projector and
      the same operator.
\end{enumerate}
The theorem then applies the rebuilt continuum-bridge equivalence theorems to these
generic constructions. So the variational, observer, and projection surfaces
are not extra analytic assumptions. They are simply presentations of the same
selected law object.

Finally, because the canonical promotion context has zero constitutive term by
construction, its promoted observed law is exactly the operator term again.
Thus even the promotion surface has already collapsed back to the same
selected law. This is the precise sense in which the detached analytic lane
has already become one selected operator-residual law together with multiple
equivalent presentations, rather than a pile of separate analytic layers.
\end{proof}

\begin{proof}[Proof of Theorem~\ref{thm:the-selected-operator-residual-law-is-uniquely-forced}]
\label{proof:thm:the-selected-operator-residual-law-is-uniquely-forced}
This theorem turns the previous law-shaped packaging result into a genuine
forcing statement. The admissibility surface is formalized by
\lean{ClosureCurrent.LocalEventFourStateLaw.selectedFieldLawClass}, and the
singleton forcing package itself is
\lean{ClosureCurrent.SelectedFieldLawForcingSurface}. The theorem witness is
\lean{ClosureCurrent.LocalEventFourStateLaw.selectedFieldLawForcingSurface}.

An admissible competitor \(\widetilde{\mathcal{L}}_{\mathrm{sel}}\) is allowed
to vary only in the one place where genuine freedom might still remain: its
operator. The selected projector is required to be the same
\(P_{\mathrm{sel}}\), and the residual is required to be the same governing
scalar residual \(r^{\mathrm{gov}}_{\mathrm{sel}}\). So the only substantive
question is whether a different admissible operator could survive.

But admissibility for this law class means admissibility for the same forced
four-state continuum limit class. The earlier selected-field theorem already
showed that the detached continuum law is exactly the operator term of
\(\mathcal{L}_{\mathrm{sel}}\). Therefore the existing continuum forcing
theorem forces any admissible competitor operator to agree pointwise with that
same selected operator.

Once the projector, operator, and residual all agree, the entire selected law
object agrees. The remaining fields in the formal structure are only proof
data attached to those same three mathematical components. So the detached
analytic lane now carries not merely one elegant selected operator-residual
law, but one uniquely forced selected operator-residual law.
\end{proof}

\begin{proof}[Proof of Theorem~\ref{thm:all-current-selected-cut-analytic-presentations-collapse-to-one-common-selected-equation}]
\label{proof:thm:all-current-selected-cut-analytic-presentations-collapse-to-one-common-selected-equation}
This theorem is formalized by
\lean{ClosureCurrent.LocalEventFourStateLaw.selectedFieldEquation},
\lean{ClosureCurrent.SelectedFieldEquationSurface}, and the theorem
\lean{ClosureCurrent.LocalEventFourStateLaw.selectedFieldEquationSurface}.

Its first block is a package of operator identities. The four-state continuum
law, the minimal quadratic law, the exact effective selected law, and the
promoted observed law all carry the same operator \(G_{\mathrm{sel}}\) on the
least-motion cut. So the detached analytic lane no longer contains several
competing governing operators; it already contains one common selected
operator.

Its second block is a package of residual identities. The quadratic scalar
residual and the effective scalar residual are both exactly the same governing
residual \(r_{\mathrm{sel}}\). So there are not several scalar equations left
either; there is already one common scalar equation
\(r_{\mathrm{sel}}(x)=0\).

The last block identifies the zero sets. The regular variational residuals,
the observer-side Eulerian/Lagrangian/characteristic residuals, and the
Hamilton--Jacobi projected characteristic residuals are all canonically built
from that same residual. Therefore their vanishing conditions are equivalent
to the single equation \(r_{\mathrm{sel}}(x)=0\).

So at the strongest current honest surface, the detached selected-cut analytic
lane has already collapsed to one uniquely forced selected equation with
multiple equivalent analytic presentations. What is still missing is not a
common equation surface, but only a stronger realization of that equation as a
PDE or action principle.
\end{proof}

\begin{proof}[Proof of Theorem~\ref{thm:any-later-pde-or-action-principle-flagship-form-must-realize-the-same-selected-equation}]
\label{proof:thm:any-later-pde-or-action-principle-flagship-form-must-realize-the-same-selected-equation}
The theorem deliberately does not assert that a PDE or action principle has
already been constructed. Instead it isolates the exact rigidity such a later
construction would have to satisfy. A PDE presentation is represented only by a
residual field and its zero set; an action presentation is represented only by
an Euler--Lagrange residual and its stationary set.

Once the detached lane has already forced one common selected equation
\(r_{\mathrm{sel}}(x)=0\), an honest PDE realization is exactly one whose
solution set equals that selected equation set. Likewise, an honest
action-style realization is exactly one whose stationary set equals that same
selected equation set.

Therefore any later PDE flagship form and any later action-principle flagship
form already have the same solution set. They would not introduce new
law-level freedom; they would only present the same uniquely forced selected
equation in a different analytic language.
\end{proof}

\begin{proof}[Proof of Theorem~\ref{thm:canonical-scalar-pde-form-and-action-form-realizations-already-exist}]
\label{proof:thm:canonical-scalar-pde-form-and-action-form-realizations-already-exist}
The proof is intentionally direct. The already forced selected residual
\(r_{\mathrm{sel}}\) is itself used as the PDE residual field, so the PDE-form
equation is simply \(\mathcal{E}_{\mathrm{sel}}(x)=r_{\mathrm{sel}}(x)=0\).
The same residual is also used as the Euler--Lagrange residual field, so the
action-form stationary equation is simply
\(\delta S_{\mathrm{sel}}(x)=r_{\mathrm{sel}}(x)=0\).

Thus the theorem really does provide current realization existence, not merely
future rigidity. What it does \emph{not} provide is a stronger analytic
interpretation of these residual fields as a spatial differential PDE or a
fully constructed action functional. Those remain later realization problems.
\end{proof}

\begin{proof}[Proof of Theorem~\ref{thm:the-current-pde-form-and-action-form-realizations-are-the-same-scalar-residual-field}]
\label{proof:thm:the-current-pde-form-and-action-form-realizations-are-the-same-scalar-residual-field}
The proof is immediate from the canonical constructions. The PDE-form
realization defines its residual field to be the selected residual
\(r_{\mathrm{sel}}\). The action-form realization defines its Euler--Lagrange
field to be that same selected residual \(r_{\mathrm{sel}}\). Therefore the
two fields agree pointwise by construction.

This is exactly the right reading of the present PDE language. The detached
lane already has a genuine equation object, namely one scalar residual field
whose vanishing defines the selected law. But it has not yet forced a spatial
differential operator or a full action functional beyond that scalar residual
surface.
\end{proof}

\begin{proof}[Proof of Theorem~\ref{thm:the-current-scalar-realizations-come-from-one-selected-quadratic-density}]
\label{proof:thm:the-current-scalar-realizations-come-from-one-selected-quadratic-density}
The proof is again direct, but it says something genuinely stronger than the
previous theorem. The selected density is defined pointwise by
\[
q_{\mathrm{sel}}(x)=\langle x,G_{\mathrm{sel}}x\rangle,
\]
where \(G_{\mathrm{sel}}\) is the same governing selected operator already
forced on the least-motion cut. By construction, the selected residual is its
signed-count scalarization:
\[
r_{\mathrm{sel}}(x)=\operatorname{SignedCount.ofNat}\!\bigl(q_{\mathrm{sel}}(x)\bigr).
\]

Now the canonical PDE-form realization uses exactly \(r_{\mathrm{sel}}\) as
its equation field, and the canonical action-form realization uses exactly
\(r_{\mathrm{sel}}\) as its Euler--Lagrange field. Therefore both are
pointwise equal to the same signed-count scalarization of the same selected
quadratic density.

So the present flagship realization is already more concrete than a bare
residual equation object. It is still not a spatial differential PDE or a full
action functional, but it is already the scalarization of one explicit
quadratic density on the selected state space.
\end{proof}

\begin{proof}[Proof of Theorem~\ref{thm:the-selected-operator-law-already-has-a-canonical-no-hidden-crime-realization-ledger}]
\label{proof:thm:the-selected-operator-law-already-has-a-canonical-no-hidden-crime-realization-ledger}
The construction is deliberately minimal. The realization law is chosen so
that its continuum residual is exactly the selected operator
\(\mathcal R_{\mathrm{cont}}(x)=G_{\mathrm{sel}}x\). The realization chain
then sets every realized residual representative equal to that same selected
operator. So the only remaining issue is whether any hidden realization ledger
crime survives.

The theorem shows that none does. Projection, geometry, quadrature, aliasing,
split, linearization, algebraic, boundary, and interface crime all vanish
pointwise. Therefore the realized residual is exactly the continuum residual.
So the selected operator law already has an exact host-side
realization-chain presentation.

Finally, the same theorem reads the selected density and scalar residual
directly from that exact realized residual:
\[
q_{\mathrm{sel}}(x)=\langle x,\mathcal R_{\mathrm{sel}}(x)\rangle,
\qquad
r_{\mathrm{sel}}(x)=\operatorname{SignedCount.ofNat}\!\bigl(\langle x,\mathcal R_{\mathrm{sel}}(x)\rangle\bigr).
\]
Thus the present flagship representative is exact in a second sense as well:
not only on the selected-cut law side, but also on the explicit realization
ledger.
\end{proof}

\begin{proof}[Proof of Theorem~\ref{thm:the-selected-quadratic-density-is-already-an-exact-adjoint-goal-ledger}]
\label{proof:thm:the-selected-quadratic-density-is-already-an-exact-adjoint-goal-ledger}
The certificate uses the identity operator, so its exact solved state is just
its right-hand side \(x+G_{\mathrm{sel}}x\), while its algebraic residual is
exactly \(G_{\mathrm{sel}}x\). The adjoint-goal identity then applies without
any extra analytic input: the goal error equals the adjoint-weighted algebraic
residual, and also equals the weighted residual sum of the listed ledger
terms.

Because the certificate carries the goal vector \(x\), adjoint weight \(x\),
and a singleton residual ledger \([G_{\mathrm{sel}}x]\), those general
identities collapse immediately to
\[
\eta_{\mathrm{sel}}(x)
  =\langle x,G_{\mathrm{sel}}x\rangle
  =q_{\mathrm{sel}}(x),
\]
and the weighted residual ledger has exactly that one contribution.
The final scalar residual statement follows by the previously proved density
scalarization surface.
\end{proof}

\begin{proof}[Proof of Theorem~\ref{thm:the-selected-cut-already-carries-one-canonical-scalar-law-object}]
\label{proof:thm:the-selected-cut-already-carries-one-canonical-scalar-law-object}
This theorem identifies the current scalar surfaces already proved separately
as one canonical scalar law on the selected cut.

The theorem bundles six previously proved ingredients on the same selected
cut:
\begin{enumerate}
\item the exact selected-equation boundary for PDE/action realizations;
\item the existence of the canonical scalar PDE-form and scalar action-form
      realizations;
\item the pointwise equality of those two scalar residual fields;
\item the explicit selected quadratic density whose signed-count scalarization
      gives that common residual;
\item the exact no-hidden-crime realization ledger;
\item the exact adjoint-goal ledger.
\end{enumerate}

The first three ingredients say that the forced selected equation,
the canonical scalar PDE-form realization, and the canonical scalar
action-form realization already have the same zero set. The density theorem
adds that all of those residual fields are pointwise the same
signed-count scalarization of
\[
q_{\mathrm{sel}}(x)=\langle x,G_{\mathrm{sel}}x\rangle.
\]
The realization-ledger theorem then rewrites that same density as
\[
q_{\mathrm{sel}}(x)=\langle x,\mathcal R_{\mathrm{sel}}(x)\rangle,
\]
and the adjoint-goal theorem rewrites it again as the exact goal error
\[
q_{\mathrm{sel}}(x)=\eta_{\mathrm{sel}}(x).
\]
So the same selected cut already carries one canonical scalar law object
before any further differential realization is attempted.
\end{proof}

\begin{proof}[Proof of Theorem~\ref{thm:the-selected-cut-already-carries-one-canonical-first-order-law-object}]
\label{proof:thm:the-selected-cut-already-carries-one-canonical-first-order-law-object}
This theorem identifies the exact selected flow package and the scalar-law
package as the same first-order law object on the selected cut.

The theorem bundles three previously proved ingredients on the same selected
cut:
\begin{enumerate}
\item the uniqueness of the selected operator-residual law, via
      \lean{ClosureCurrent.LocalEventFourStateLaw.selectedFieldLawForcingSurface};
\item the canonical scalar law object, via
      \lean{ClosureCurrent.LocalEventFourStateLaw.selectedScalarLawSurface};
\item the exact selected effective-flow package, via
      \lean{ClosureCurrent.LocalEventFourStateLaw.selectedGoverningFlowSurface}.
\end{enumerate}

The selected-flow theorem gives that every stage of the exact effective flow
already equals the same selected operator \(G_{\mathrm{sel}}\). The scalar-law
theorem identifies the governing residual with the signed-count scalarization
of the selected quadratic density, and the new first-order package rewrites
that same density through the exact selected flow value itself. Therefore:
\begin{enumerate}
\item the selected equation is exactly the vanishing of the scalar readout of
      the exact selected flow;
\item the canonical scalar PDE-form realization has that same zero set;
\item the canonical scalar action-form realization has that same zero set.
\end{enumerate}
So the detached flagship law is already first-order before any stronger
differential realization is attempted.
\end{proof}

\begin{proof}[Proof of Theorem~\ref{thm:the-selected-cut-already-carries-one-uniquely-forced-autonomous-evolution-law}]
\label{proof:thm:the-selected-cut-already-carries-one-uniquely-forced-autonomous-evolution-law}
This theorem packages the already established exact selected flow, common
selected generator, and governing residual into one autonomous evolution law
object.

The first-order theorem already proves that every stage of the exact selected
flow is the same selected operator and that the scalar selected equation is the
same zero set as the scalar readout of that flow. The new theorem promotes
that surface to the concrete evolution law
\[
(\Phi_{\mathrm{sel}}^{(n)},G_{\mathrm{sel}},r_{\mathrm{sel}})
\]
with stagewise identity
\[
\Phi_{\mathrm{sel}}^{(n)}(x)=G_{\mathrm{sel}}x.
\]

The forcing part is then immediate from the previously established forcing of
the selected operator itself. The admissible-law class now asks only for three
things:
\begin{enumerate}
\item the generator lies in the same forced four-state continuum class;
\item the residual is the same governing selected residual;
\item the stagewise flow is exactly that generator.
\end{enumerate}
Once these are imposed, the generator is already forced by the earlier
selected-operator theorem, and the stagewise flow is forced with it. So the
autonomous evolution law is unique.

Finally, because the residual of this evolution law is exactly
\(r_{\mathrm{sel}}\), its solution set is exactly the selected equation
\[
r_{\mathrm{sel}}(x)=0,
\]
and therefore exactly the current scalar PDE-form solution set and the current
scalar action-form stationary set.
\end{proof}

\begin{proof}[Proof of Theorem~\ref{thm:the-selected-cut-already-carries-one-uniquely-forced-autonomous-scalar-law}]
\label{proof:thm:the-selected-cut-already-carries-one-uniquely-forced-autonomous-scalar-law}
This theorem collapses the already-established evolution-law and scalar-law
surfaces to one autonomous scalar law object.

The structural gain is that the autonomous scalar law already carries its own
stagewise autonomy, density identity, and scalarization identity as part of
its definition. So the admissible-law class now asks only for one thing: the
generator lies in the same forced four-state continuum class. Once that
generator is forced, the stagewise flow is forced with it, the density is
forced to be
\[
q_{\mathrm{sel}}(x)=\langle x,G_{\mathrm{sel}}x\rangle,
\]
and the residual is forced to be
\[
r_{\mathrm{sel}}(x)=\operatorname{sc}\!\bigl(q_{\mathrm{sel}}(x)\bigr).
\]

The scalar-law theorem then adds the two exact ledger identifications
\[
q_{\mathrm{sel}}(x)
 =\langle x,\mathcal R_{\mathrm{sel}}(x)\rangle
 =\eta_{\mathrm{sel}}(x),
\]
while the evolution theorem identifies the zero set of the new law with the
selected equation and hence with the canonical scalar PDE/action zero set and
with the autonomous evolution-law solution set. So the detached selected cut
already carries one uniquely forced autonomous scalar law before any stronger
differential realization is attempted.
\end{proof}

\begin{proof}[Proof of Theorem~\ref{thm:the-selected-cut-already-fixes-the-public-contract-for-any-later-stronger-flagship-realization}]
\label{proof:thm:the-selected-cut-already-fixes-the-public-contract-for-any-later-stronger-flagship-realization}
This theorem is formalized by
\lean{ClosureCurrent.FlagshipRealization},
\lean{ClosureCurrent.SelectedFlagshipRealizationSurface}, and the witness
\lean{ClosureCurrent.LocalEventFourStateLaw.selectedFlagshipRealizationSurface}.
It records the exact public contract any later stronger flagship form must
satisfy.

The key point is that the realization interface only leaves one genuinely free
public datum: the velocity field. The theorem requires that velocity to lie in
the same forced four-state continuum class and then requires the remaining
public quantities to be read honestly from it:
\begin{enumerate}
\item the density is the state--velocity pairing;
\item the residual is the signed-count scalarization of that density;
\item the solution set is the zero set of that residual and is exact-selected.
\end{enumerate}
Because the velocity lies in the same forced continuum class, it is already
forced to equal \(G_{\mathrm{sel}}\). Once that happens, the density is
forced to be \(q_{\mathrm{sel}}\), the residual is forced to be
\(r_{\mathrm{sel}}\), and the solution set is forced to be the same selected
equation carried by the autonomous scalar law.

The last two clauses then sharpen that contract further by comparison with the
already established ledger surfaces: the same density is also
\[
\langle x,\mathcal R_{\mathrm{sel}}(x)\rangle
\quad\text{and}\quad
\eta_{\mathrm{sel}}(x),
\]
and the same solution set is also the canonical scalar PDE solution set, the
canonical scalar action stationary set, and the autonomous evolution-law
solution set. So any later stronger flagship language is forced to present
this same public law data rather than to introduce a different law.
\end{proof}

\begin{proof}[Proof of Theorem~\ref{thm:the-selected-cut-already-fixes-the-exact-realization-ledger-adjoint-goal-contract-for-any-later-stronger-flagship-realization}]
\label{proof:thm:the-selected-cut-already-fixes-the-exact-realization-ledger-adjoint-goal-contract-for-any-later-stronger-flagship-realization}
This theorem is formalized by
\lean{ClosureCurrent.LedgerGoalFlagshipRealization},
\lean{ClosureCurrent.SelectedLedgerGoalFlagshipRealizationSurface}, and the
witness
\lean{ClosureCurrent.LocalEventFourStateLaw.selectedLedgerGoalFlagshipRealizationSurface}.
It sharpens the flagship contract on the exact host-realization side.

The realization interface extends the earlier public flagship contract by
requiring:
\begin{enumerate}
\item one exact no-hidden-crime realization chain for the canonical selected
      realization law;
\item one exact adjoint-goal certificate;
\item exact identification of the continuum residual and the algebraic
      residual with the same realized flagship velocity field;
\item exact identification of the goal error with the same realized flagship
      density.
\end{enumerate}
Because the flagship velocity is already forced to be \(G_{\mathrm{sel}}\),
the continuum residual is forced to be the same selected operator. Exact
no-hidden-crime then forces the terminal realized residual to be that same
selected operator. The public flagship theorem had already forced the density
to be \(q_{\mathrm{sel}}\), and therefore also to equal the realized density
pairing and the exact goal error. So the solved state, algebraic residual,
and exact goal error of any later stronger flagship realization are already
fixed by the selected cut. This sharpens the contract from the public law data
to the exact realization-ledger / adjoint-goal ledger themselves.
\end{proof}

\begin{proof}[Proof of Theorem~\ref{thm:the-selected-cut-already-fixes-the-differential-flow-contract-for-any-later-stronger-flagship-realization}]
\label{proof:thm:the-selected-cut-already-fixes-the-differential-flow-contract-for-any-later-stronger-flagship-realization}
This theorem is formalized by
\lean{ClosureCurrent.DifferentialFlagshipRealization},
\lean{ClosureCurrent.SelectedDifferentialFlagshipRealizationSurface}, and the
witness
\lean{ClosureCurrent.LocalEventFourStateLaw.selectedDifferentialFlagshipRealizationSurface}.
It sharpens the flagship contract from the public law data to the exact
continuum-bridge differential-flow package.

The realization interface extends the earlier flagship contract by requiring:
\begin{enumerate}
\item one concrete continuous block-derivative datum;
\item one concrete continuous effective-flow datum;
\item exact identification of the block operator, the \(PP\) block, and the
      effective flow with the same realized flagship velocity field;
\item exact vanishing of the \(PQ\), \(QP\), and \(QQ\) blocks.
\end{enumerate}
Because the flagship velocity is already forced to be \(G_{\mathrm{sel}}\),
the block operator is forced with it. The selected stationary-cut theorem then
forces the \(PP\) block to be the same selected generator and the
\(PQ/QP/QQ\) blocks to vanish exactly. The selected governing-flow theorem
forces the effective flow to be that same generator stagewise. So any later
differential flagship language is constrained not only at the common zero-set
level, but also at the exact selected block-derivative/effective-flow level.
\end{proof}

\begin{proof}[Proof of Theorem~\ref{thm:the-selected-cut-already-fixes-the-derivative-subsampling-contract-for-any-later-stronger-flagship-realization}]
\label{proof:thm:the-selected-cut-already-fixes-the-derivative-subsampling-contract-for-any-later-stronger-flagship-realization}
This theorem is formalized by
\lean{ClosureCurrent.DerivativeSubsamplingFlagshipRealization},
\lean{ClosureCurrent.SelectedDerivativeSubsamplingFlagshipSurface}, and the
witness
\lean{ClosureCurrent.LocalEventFourStateLaw.selectedDerivativeSubsamplingFlagshipSurface}.
It sharpens the flagship contract again, now in the continuum-bridge classical
derivative-subsampling language.

The realization interface extends the earlier flagship contract by requiring:
\begin{enumerate}
\item one stagewise projection family;
\item one derivative field;
\item one observed restriction family, identified with the observed stagewise
      readout of that derivative through the selected projection;
\item one defect family, required to vanish exactly on the selected cut.
\end{enumerate}
Because the flagship velocity is already forced to be \(G_{\mathrm{sel}}\),
the derivative is forced with it. Once the selected cut is fixed, the
projection family is forced to be the constant selected projection, and the
observed restriction is forced to be the selected projection of the same
selected generator. The defect then collapses identically to zero. So any
later derivative-like flagship language is constrained not only at the
zero-set level and not only at the block-flow level, but also at this first
classical derivative-subsampling level.
\end{proof}

\begin{proof}[Proof of Theorem~\ref{thm:the-selected-cut-already-fixes-the-variational-contract-for-any-later-stronger-flagship-realization}]
\label{proof:thm:the-selected-cut-already-fixes-the-variational-contract-for-any-later-stronger-flagship-realization}
This theorem is formalized by
\lean{ClosureCurrent.VariationalFlagshipRealization},
\lean{ClosureCurrent.SelectedVariationalFlagshipRealizationSurface}, and the
witness
\lean{ClosureCurrent.LocalEventFourStateLaw.selectedVariationalFlagshipRealizationSurface}.
It sharpens the flagship contract on the action side.

The realization interface extends the earlier flagship contract by requiring:
\begin{enumerate}
\item one regular variational system;
\item one observer-side conservative transport datum;
\item one Hamilton--Jacobi projection bridge;
\item exact identification of every residual in those three packages with the
      same realized flagship residual.
\end{enumerate}
Because the flagship residual is already forced to be \(r_{\mathrm{sel}}\),
every residual in the dynamics-side variational triad and every residual in
the observer-side transport/Hamilton--Jacobi bridge is forced with it. So any
later action-like flagship language is constrained not only at the common
zero-set level, but also at the exact variational / observer /
Hamilton--Jacobi residual level.
\end{proof}

\begin{proof}[Proof of Theorem~\ref{thm:the-selected-cut-already-fixes-the-exact-combined-pre-pde-contract-for-any-later-stronger-flagship-realization}]
\label{proof:thm:the-selected-cut-already-fixes-the-exact-combined-pre-pde-contract-for-any-later-stronger-flagship-realization}
This theorem is formalized by
\lean{ClosureCurrent.CompositeFlagshipRealization},
\lean{ClosureCurrent.SelectedCompositeFlagshipRealizationSurface}, and the
witness
\lean{ClosureCurrent.LocalEventFourStateLaw.selectedCompositeFlagshipRealizationSurface}.
It is the combined realization-side boundary theorem.

The input is one later stronger flagship language that carries, on one same
selected realization, all of the exact contract layers already established
separately:
\begin{enumerate}
\item the public flagship contract;
\item the exact realization-ledger / adjoint-goal contract;
\item the continuum-bridge differential-flow contract;
\item the continuum-bridge derivative-subsampling contract;
\item the Chapter~7 variational / observer / Hamilton--Jacobi contract.
\end{enumerate}
Each earlier theorem had already forced its own visible surfaces to agree with
the selected cut. The new theorem then proves the joint collapse: once all
those languages are required to live on one same selected realization, they do
not carry independent law data anymore. They collapse to one common operator,
one common density, one common scalar residual, and one common zero set. So
the remaining boundary before a final spatial differential PDE or full action
functional is now a pure realization problem rather than a remaining law-choice
problem.
\end{proof}

\begin{proof}[Proof of Theorem~\ref{thm:the-selected-cut-already-fixes-the-eventual-realized-form-up-to-differential-language}]
\label{proof:thm:the-selected-cut-already-fixes-the-eventual-realized-form-up-to-differential-language}
The theorem is formalized by
\lean{ClosureCurrent.SelectedRealizedFormSurface} and
\lean{ClosureCurrent.LocalEventFourStateLaw.selectedRealizedFormSurface}. It is
the strongest law-shaped restatement currently available of the combined
realization-side contract.

Its input is still one later stronger flagship language carrying, on one same
selected realization, all five exact contract layers already established
earlier:
\begin{enumerate}
\item the public flagship contract;
\item the exact realization-ledger / adjoint-goal contract;
\item the continuum-bridge differential-flow contract;
\item the continuum-bridge derivative-subsampling contract;
\item the Chapter~7 variational / observer / Hamilton--Jacobi contract.
\end{enumerate}
The previous combined theorem had already shown that those five layers collapse
jointly to one operator, one density, one scalar residual, and one zero set.
The new theorem sharpens that collapse into explicit realized form. On every
stage \(n\) and state \(x\):
\begin{enumerate}
\item the public velocity, block operator, derivative field, and effective
      flow all coincide with the same selected generator \(G_{\mathrm{sel}}\);
\item the mixed and hidden differential blocks vanish, so the realized
      differential presentation is already block-diagonal on the selected cut;
\item the derivative defect vanishes, so the observed restriction is already
      exact on that same cut;
\item the common density is exactly the selected quadratic pairing
      \(\langle x,G_{\mathrm{sel}}x\rangle\), and the common scalar residual is
      exactly its signed-count scalarization;
\item the flagship solution set is exactly the selected equation and hence
      exactly the same scalar PDE/action and evolution zero sets forced
      earlier.
\end{enumerate}
So before any final spatial differential PDE or full action functional is
constructed, the detached selected cut already fixes the eventual realized form
up to differential language: first-order, autonomous, quadratic, and
selected-cut exact.
\end{proof}

\begin{proof}[Proof of Theorem~\ref{thm:the-detached-selected-law-already-has-one-underlying-state-object}]
\label{proof:thm:the-detached-selected-law-already-has-one-underlying-state-object}
The theorem is formalized by
\lean{ClosureCurrent.SelectedStateObject},
\lean{ClosureCurrent.SelectedStateObjectSurface}, and
\lean{ClosureCurrent.LocalEventFourStateLaw.selectedStateObjectSurface}. It is
the cleanest current answer to the question of what the detached selected law
is actually about.

The key move is to stop reading the current flagship surfaces as parallel
packages and to identify the common carrier they already share. That carrier is
the explicit four-coordinate state space \(\mathrm{State}\) from the rebuilt
foundation. On that carrier, the theorem packages one selected state object
consisting of:
\begin{enumerate}
\item the selected projection \(P_{\mathrm{sel}}\);
\item the selected generator \(G_{\mathrm{sel}}\);
\item the quadratic density
      \(q_{\mathrm{sel}}(x)=\langle x,G_{\mathrm{sel}}x\rangle\);
\item the scalar residual
      \(r_{\mathrm{sel}}(x)=\operatorname{sc}(q_{\mathrm{sel}}(x))\);
\item the zero set \(r_{\mathrm{sel}}(x)=0\), which is exactly the selected
      equation.
\end{enumerate}
So the current law is not merely a bundle of a velocity field, density,
residual, and equation. Those are already read directly from one common state
object.

The theorem then adds the realization-side rigidity. Any later stronger
flagship realization already agrees with this same state object on its public
law faces: same generator, same density, same scalar residual, and same zero
set. Thus the present unification target is not one scalar but one underlying
state object with several derived readouts.
\end{proof}

\begin{proof}[Proof of Theorem~\ref{thm:split-completion-bridge-recovers-the-no-stage-completion-law}]
\label{proof:thm:split-completion-bridge-recovers-the-no-stage-completion-law}
Fix a constructive sector law. The completion data above it split naturally
into two independent pieces.

The first piece is the current-side four-state assembly. It consists of a
signed readout, a genuine four-leg frame, and the exactifier/generator
intertwining identity on the assembled hidden state. The readout and frame
determine the four-state hidden realization, and adjoining the intertwining
identity upgrades that realization to the full four-state assembly. Hence the
current-side completion input may be read either as one assembled package or as
the smaller pair ``readout/frame + exactifier realization''.

The second piece is the Step~4 completion datum over the same intrinsic sector
generation. It now consists of two independent clauses: endpoint rigidity, and
a reduced representative package carrying the four endpoint-sensitive canonical
representatives together with the already classified kinetic face. Carrier
classification is recovered from that same reduced representative package.
Since the underlying constructive sector law already fixes the selected bundle
and the intrinsic sector generation, no additional choice of selected cut or
physical law is inserted at this stage; those identifications are inherited
from the same law.

Combining the current-side four-state assembly with the aligned Step~4 datum
recovers the split completion bridge. Once this bridge has been reconstructed,
the no-stage completion law follows immediately, because the bridge is exactly
the remaining ingredient needed to pass from the constructive sector law to the
completion witness.

The same decomposition shows why the earlier stronger no-stage completion law
already provides a concrete witness. That stronger law contains the same
readout/frame data, the same exactifier identity, and the same Step~4
rigidity and reduced representative clauses, from which the same weaker
classification target is recovered internally. Therefore it realizes the split
completion target directly, not merely the older assembled completion package.

Thus the completion seam has been reduced to a genuinely smaller statement.
One does not first assume the no-stage completion law as a monolithic object.
It is recovered from two visibly distinct ingredients over the constructive
sector law: the four-state assembly on the current side and the Step~4
alignment on the completion side.
\end{proof}

\begin{proof}[Proof of Theorem~\ref{thm:matched-no-stage-continuum-law-closes-on-the-active-selected-datum}]
\label{proof:thm:matched-no-stage-continuum-law-closes-on-the-active-selected-datum}
The selected Schur / HFT-2 exact-closure surface is packaged here as the
proposition
\lean{ClosureCurrent.SelectedSchurHFT2Surface}. The fact that the active
law-side selected datum on the Lean PartIV surface already satisfies that
packaged surface is
\lean{ClosureCurrent.selectedSchurHFT2Surface}; it is just a manuscript-facing
alias of the previously established selected Schur / HFT-2 closure theorem on
the active Lean surface.

The new law-side completion datum is
\lean{ClosureCurrent.LocalEventFourStateCompletionLaw}. It is deliberately the
no-stage package: the no-stage four-state law, the intrinsic sector-generation
clauses, and the Step~4 endpoint-rigidity / carrier-classification clauses all
live on one and the same selected bridge. The active natural-operator
completion is then reconstructed internally by
\lean{ClosureCurrent.LocalEventFourStateCompletionLaw.toNaturalOperatorCompletion}.

The actual closure theorem is
\lean{ClosureCurrent.LocalEventFourStateCompletionLaw.attachedSurface}. It
combines four ingredients already established separately in Lean:
\begin{enumerate}
\item the detached intrinsic selected-bundle and intrinsic-sector forcing
      surface, recovered internally from the same no-stage law;
\item the forced four-state continuum law and forced reduced hidden continuum
      law, carried by the no-stage four-state theorem
      \lean{ClosureCurrent.LocalEventFourStateLaw.continuumSurface};
\item the selected Schur / HFT-2 exact-closure surface on that same selected
      bridge, via \lean{ClosureCurrent.selectedSchurHFT2Surface};
\item the full active Lean PartIV law-completion statement on the internally
      reconstructed completion object.
\end{enumerate}
Thus the detached continuum extension is no longer merely matched to an
external active completion package. It closes back into the active selected-cut
completion surface from one law-side datum alone, with no extra completion
hypothesis and no separate bridge-identification step.
\end{proof}

\begin{proof}[Proof of Theorem~\ref{thm:no-stage-completion-law-reduces-the-remaining-flagship-seam-to-the-four-endpoint-frontiers}]
\label{proof:thm:no-stage-completion-law-reduces-the-remaining-flagship-seam-to-the-four-endpoint-frontiers}
The theorem is formalized by
\lean{ClosureCurrent.LocalEventFourStateCompletionLaw.frontierBoundarySurface}.
Its first component is exactly the law-native completion theorem already
proved earlier: the no-stage completion law reconstructs the active
\lean{PartIVLawCompletionStatement} on its own selected datum.

The second component is the sharper point. It is universal over
\lean{ClosureCurrent.FlagshipFrontiers}: once a phase frontier, a wave
frontier, a gauge frontier, and a geometric frontier are all read on the one
intrinsically selected bundle generated by the no-stage completion law, the
theorem returns the four recognizable identity witnesses directly. Internally,
it applies the already established lane-wise theorems
\lean{ClosureCurrent.LocalEventFourStateCompletionLaw.phaseSurface},
\lean{ClosureCurrent.LocalEventFourStateCompletionLaw.waveSurface},
\lean{ClosureCurrent.LocalEventFourStateCompletionLaw.gaugeSurface}, and
\lean{ClosureCurrent.LocalEventFourStateCompletionLaw.geometricSurface}, and
then repackages their nonempty recognizable witnesses into one aggregate
surface.

So the theorem does not claim that the common completion law already produces
the endpoint frontiers. It says something more precise: once the law-native
completion surface is reconstructed, the remaining flagship burden is exactly
those four frontier witness bundles and no extra completion or datum-matching
input survives.
\end{proof}

\begin{proof}[Proof of Theorem~\ref{thm:the-remaining-flagship-seam-has-an-exact-endpoint-witness-form}]
\label{proof:thm:the-remaining-flagship-seam-has-an-exact-endpoint-witness-form}
The sharper theorem is formalized by
\lean{ClosureCurrent.LocalEventFourStateCompletionLaw.endpointWitnessBoundarySurface}.
Its inputs are not the coarser frontier bundles anymore, but the exact
lane-specific endpoint witness packages:
\lean{ClosureCurrent.PhaseEndpointWitness},
\lean{ClosureCurrent.WaveEndpointWitness},
\lean{ClosureCurrent.GaugeEndpointWitness},
and \lean{ClosureCurrent.GeometricEndpointWitness}, grouped by
\lean{ClosureCurrent.FlagshipEndpointWitnesses}.

These witness packages do not add new content. They unpack the previous
frontier layer into the exact endpoint-collapse datum plus the minimal
auxiliary corollary witnesses still needed in each lane:
\begin{enumerate}
\item phase: phase-compatible Schur-tower data and observable self-adjoint
      scalarization;
\item wave: selected-channel transport generation, scalar recursion, local
      internal datum, and admissible transport;
\item gauge: selected-channel gauge covariance;
\item geometry: orthogonal symmetric-tensor scalarization.
\end{enumerate}
The conversion maps
\lean{ClosureCurrent.PhaseEndpointWitness.toFrontier},
\lean{ClosureCurrent.WaveEndpointWitness.toFrontier},
\lean{ClosureCurrent.GaugeEndpointWitness.toFrontier},
\lean{ClosureCurrent.GeometricEndpointWitness.toFrontier}, and
\lean{ClosureCurrent.FlagshipEndpointWitnesses.toFrontiers}
repackage these exact witnesses into the coarser frontier bundles, while
\lean{ClosureCurrent.PhaseFrontier.toEndpointWitness},
\lean{ClosureCurrent.WaveFrontier.toEndpointWitness},
\lean{ClosureCurrent.GaugeFrontier.toEndpointWitness},
\lean{ClosureCurrent.GeometricFrontier.toEndpointWitness}, and
\lean{ClosureCurrent.FlagshipFrontiers.toEndpointWitnesses}
run the same passage in the reverse direction.

Above one fixed no-stage completion witness, these exact witness packages are
now also named as four separate targets:
\lean{ClosureCurrent.PhaseEndpointWitnessTarget},
\lean{ClosureCurrent.WaveEndpointWitnessTarget},
\lean{ClosureCurrent.GaugeEndpointWitnessTarget}, and
\lean{ClosureCurrent.GeometricEndpointWitnessTarget}. The theorem
\lean{ClosureCurrent.NoStageCompletionWitness.flagshipExactBoundaryTargetOfLaneTargets}
packages those four lane targets into the exact fixed-witness flagship
boundary target \lean{ClosureCurrent.FlagshipExactBoundaryTarget}. The theorem
\lean{ClosureCurrent.NoStageCompletionWitness.endpointWitnessTargetOfLaneTargets}
shows that these four lane targets reassemble the aggregate
\lean{ClosureCurrent.EndpointWitnessTarget}.

The equivalence with the frontier presentation is now formal at the target
level, not just at the datum level. The theorems
\lean{ClosureCurrent.NoStageCompletionWitness.phaseEndpointWitnessTargetOfFrontierTarget},
\lean{ClosureCurrent.NoStageCompletionWitness.waveEndpointWitnessTargetOfFrontierTarget},
\lean{ClosureCurrent.NoStageCompletionWitness.gaugeEndpointWitnessTargetOfFrontierTarget},
and
\lean{ClosureCurrent.NoStageCompletionWitness.geometricEndpointWitnessTargetOfFrontierTarget}
show that the four law-native frontier targets already realize the four exact
endpoint-witness lane targets on the same fixed completion witness. The
aggregate comparison theorems
\lean{ClosureCurrent.NoStageCompletionWitness.flagshipBoundaryTargetOfExactBoundaryTarget}
and
\lean{ClosureCurrent.NoStageCompletionWitness.flagshipExactBoundaryTargetOfBoundaryTarget}
show that the frontier-side and exact-side fixed-witness flagship boundary
targets are interchangeable. The aggregate theorem is
\lean{ClosureCurrent.NoStageCompletionWitness.endpointWitnessTargetOfFrontierTargets},
and the corresponding endpoint-complete flagship surface is
\lean{ClosureCurrent.NoStageCompletionWitness.flagshipSurfaceOfFrontierTargets}.
On the exact side, the matching surface theorem is
\lean{ClosureCurrent.NoStageCompletionWitness.flagshipSurfaceOfExactBoundaryTarget}.

The theorem then simply applies the already established frontier-boundary
theorem after that repackaging. So the gain is not a stronger flagship
identity theorem yet. The gain is a sharper honest description of what still
remains explicit: the exact lane endpoint-collapse packages and their minimal
auxiliary witness data, rather than an opaque frontier placeholder.

This exact endpoint-witness layer now also connects directly to the direct
normal-form surface. The theorem
\lean{ClosureCurrent.NoStageCompletionWitness.normalFormTargetOfEndpointWitnessTarget}
shows that, above one fixed no-stage completion witness, the exact
endpoint-witness target already implies the direct target
\lean{ClosureCurrent.NormalFormTarget}; the corresponding surface theorem is
\lean{ClosureCurrent.NoStageCompletionWitness.normalFormTargetSurfaceOfEndpointWitnessTarget}.
At the existence level, this yields
\lean{ClosureCurrent.MicroscopicClosureLaw.normalFormTargetOfFlagshipTarget}
and
\lean{ClosureCurrent.MicroscopicClosureLaw.normalFormTargetSurfaceOfFlagshipTarget}.
So the endpoint-witness theorem is no longer only a sharper boundary
description. It is now also a proved bridge from the stronger flagship target
surface to the direct recognizable normal-form target.

That bridge is now split one step more finely than before. The theorem
\lean{ClosureCurrent.NoStageCompletionWitness.normalFormTargetOfEndpointLaneTargets}
shows directly that the four exact lane-specific endpoint-witness targets
already imply \lean{ClosureCurrent.NormalFormTarget}; the corresponding
surface theorem is
\lean{ClosureCurrent.NoStageCompletionWitness.normalFormTargetSurfaceOfEndpointLaneTargets}.
Because the four frontier targets already imply those four lane-specific
endpoint targets, the frontier and exact-endpoint presentations are now
interchangeable on a fixed no-stage completion witness.
The exact fixed-witness flagship boundary target makes that equivalence
explicit: \lean{ClosureCurrent.NoStageCompletionWitness.normalFormTargetOfExactBoundaryTarget}
and
\lean{ClosureCurrent.NoStageCompletionWitness.normalFormTargetSurfaceOfExactBoundaryTarget}
show that the exact-side boundary target already yields the same direct
normal-form target and surface as the frontier-side boundary target.
The older aggregate theorem
\lean{ClosureCurrent.NoStageCompletionWitness.normalFormTargetOfEndpointWitnessTarget}
is now just the wrapper obtained by unpacking one aggregate endpoint witness
bundle into those four lane targets.

This exact-side form is now also promoted to the law level. The target
\lean{ClosureCurrent.FlagshipExactTarget} records one microscopic law together
with one no-stage completion witness carrying that exact fixed-witness
boundary target. The theorems
\lean{ClosureCurrent.MicroscopicClosureLaw.flagshipExactTargetOfExactBoundaryTarget}
and
\lean{ClosureCurrent.MicroscopicClosureLaw.flagshipExactTargetOfBoundaryTarget}
package the fixed-witness exact-side data into that existence-level exact
target. Then
\lean{ClosureCurrent.MicroscopicClosureLaw.flagshipTargetOfExactTarget}
and
\lean{ClosureCurrent.MicroscopicClosureLaw.flagshipSurfaceOfExactTarget}
show that the exact-side law target already yields the existence-level
flagship target and its endpoint-complete surface, while
\lean{ClosureCurrent.MicroscopicClosureLaw.normalFormTargetOfExactTarget}
and
\lean{ClosureCurrent.MicroscopicClosureLaw.normalFormTargetSurfaceOfExactTarget}
show that the same exact-side law target already yields the direct normal-form
target and its surface.

The stronger endpoint-complete law already inhabits this sharper exact-side
law target by
\lean{ClosureCurrent.LocalEventFourStateFlagshipLaw.flagshipExactTarget}.
\end{proof}

\begin{proof}[Proof of Theorem~\ref{thm:split-normal-form-bridge-yields-the-flagship-identities}]
\label{proof:thm:split-normal-form-bridge-yields-the-flagship-identities}
The bridge is formalized by \lean{ClosureCurrent.NormalFormBridge}. It extends
the earlier split completion bridge by four lane-specific sector-rigidity
packages together with explicit equalities saying that all four live on the
same aligned active completion datum. These are the fields
\lean{ClosureCurrent.NormalFormBridge.phaseUsesCompletion},
\lean{ClosureCurrent.NormalFormBridge.waveUsesCompletion},
\lean{ClosureCurrent.NormalFormBridge.gaugeUsesCompletion}, and
\lean{ClosureCurrent.NormalFormBridge.geometricUsesCompletion}.

Before reassembling that bridge, the four lane-specific packages are now named
separately as
\lean{ClosureCurrent.PhaseRigidityDatum},
\lean{ClosureCurrent.WaveRigidityDatum},
\lean{ClosureCurrent.GaugeRigidityDatum}, and
\lean{ClosureCurrent.GeometricRigidityDatum}. Their exact existence targets
above one fixed completion bridge are
\lean{ClosureCurrent.PhaseRigidityTarget},
\lean{ClosureCurrent.WaveRigidityTarget},
\lean{ClosureCurrent.GaugeRigidityTarget}, and
\lean{ClosureCurrent.GeometricRigidityTarget}. Any actual lane datum realizes
its target by the corresponding theorems
\lean{ClosureCurrent.PhaseRigidityDatum.target},
\lean{ClosureCurrent.WaveRigidityDatum.target},
\lean{ClosureCurrent.GaugeRigidityDatum.target}, and
\lean{ClosureCurrent.GeometricRigidityDatum.target}. That older bridge still
exists, but it is no longer the smallest public seam.

The sharper theorem now works one step lower, directly above a fixed
no-stage completion witness. The four smaller lane packages are
\lean{ClosureCurrent.PhaseFrontierTarget},
\lean{ClosureCurrent.WaveFrontierTarget},
\lean{ClosureCurrent.GaugeFrontierTarget}, and
\lean{ClosureCurrent.GeometricFrontierTarget}. Each one simply asks for the
corresponding exact endpoint frontier on the witness's own selected bundle.
These are the law-native inputs already consumed by the completion-law
surface theorems.

One step below those frontier targets sit the smaller recognizable lane
targets
\lean{ClosureCurrent.PhaseRecognizableTarget},
\lean{ClosureCurrent.WaveRecognizableTarget},
\lean{ClosureCurrent.GaugeRecognizableTarget}, and
\lean{ClosureCurrent.GeometricRecognizableTarget}. The reduction theorems
\lean{ClosureCurrent.NoStageCompletionWitness.phaseRecognizableTargetOfFrontierTarget},
\lean{ClosureCurrent.NoStageCompletionWitness.waveRecognizableTargetOfFrontierTarget},
\lean{ClosureCurrent.NoStageCompletionWitness.gaugeRecognizableTargetOfFrontierTarget},
and
\lean{ClosureCurrent.NoStageCompletionWitness.geometricRecognizableTargetOfFrontierTarget}
show that each law-native frontier target already yields the corresponding
recognizable target by applying the completion law's own selected-bundle
closure theorem in that lane. So the recognizable data are no longer primitive
at the sharpest public seam. They are forced one step above the completion
law-native frontier input.

If one unpacks those recognizable targets, the underlying lane data are
\lean{ClosureCurrent.PhaseRecognizableDatum},
\lean{ClosureCurrent.WaveRecognizableDatum},
\lean{ClosureCurrent.GaugeRecognizableDatum}, and
\lean{ClosureCurrent.GeometricRecognizableDatum}. Each one keeps only:
\begin{enumerate}
\item the lane-specific endpoint collapse;
\item the proof that this collapse is read on the witness's own selected
      bundle;
\item the resulting recognizable identity in that lane.
\end{enumerate}

The reduction theorem is then
\lean{ClosureCurrent.NoStageCompletionWitness.normalFormTargetOfFrontierTargets}.
It first converts the four frontier targets into those four recognizable lane
targets, and then repackages the resulting recognizable witnesses into
\lean{ClosureCurrent.FlagshipNormalForms}, so the direct target
\lean{ClosureCurrent.NormalFormTarget} follows without taking endpoint witness
bundles or sector-rigidity packages as primitive input. The corresponding
surface theorem is
\lean{ClosureCurrent.NoStageCompletionWitness.normalFormTargetSurfaceOfFrontierTargets}.

This witness-level frontier theorem is now also promoted to an exact
existence-level target above one microscopic law:
\lean{ClosureCurrent.FlagshipFrontierTarget}. The theorems
\lean{ClosureCurrent.NoStageCompletionWitness.flagshipBoundaryTargetOfFrontierTargets},
\lean{ClosureCurrent.NoStageCompletionWitness.flagshipExactBoundaryTargetOfBoundaryTarget},
\lean{ClosureCurrent.NoStageCompletionWitness.flagshipBoundaryTargetOfExactBoundaryTarget},
\lean{ClosureCurrent.NoStageCompletionWitness.endpointWitnessTargetOfBoundaryTarget},
\lean{ClosureCurrent.NoStageCompletionWitness.flagshipSurfaceOfBoundaryTarget},
\lean{ClosureCurrent.NoStageCompletionWitness.flagshipSurfaceOfExactBoundaryTarget},
\lean{ClosureCurrent.NoStageCompletionWitness.normalFormTargetOfBoundaryTarget},
\lean{ClosureCurrent.NoStageCompletionWitness.normalFormTargetSurfaceOfBoundaryTarget},
\lean{ClosureCurrent.NoStageCompletionWitness.normalFormTargetOfExactBoundaryTarget},
and
\lean{ClosureCurrent.NoStageCompletionWitness.normalFormTargetSurfaceOfExactBoundaryTarget}
first package the four law-native frontier targets above one fixed completion
witness into the exact fixed-witness flagship boundary target
\lean{ClosureCurrent.FlagshipBoundaryTarget}. The exact-side form
\lean{ClosureCurrent.FlagshipExactBoundaryTarget} is interchangeable with it
and yields the same flagship and normal-form surfaces. For the direct
normal-form branch there is now also a smaller fixed-witness target:
\lean{ClosureCurrent.FlagshipRecognizableBoundaryTarget}. The theorems
\lean{ClosureCurrent.NoStageCompletionWitness.flagshipRecognizableBoundaryTargetOfRecognizableTargets},
\lean{ClosureCurrent.NoStageCompletionWitness.flagshipRecognizableBoundaryTargetOfFrontierTargets},
\lean{ClosureCurrent.NoStageCompletionWitness.flagshipRecognizableBoundaryTargetOfBoundaryTarget},
\lean{ClosureCurrent.NoStageCompletionWitness.flagshipRecognizableBoundaryTargetOfExactBoundaryTarget},
\lean{ClosureCurrent.NoStageCompletionWitness.normalFormTargetOfRecognizableBoundaryTarget},
and
\lean{ClosureCurrent.NoStageCompletionWitness.normalFormTargetSurfaceOfRecognizableBoundaryTarget}
show that the four smaller recognizable lane targets already reassemble one
smaller recognizable boundary target on that same witness, and that this
smaller target already suffices for the direct normal-form target and its
surface. The theorem
\lean{ClosureCurrent.MicroscopicClosureLaw.flagshipFrontierTargetOfBoundaryTarget}
then promotes one such fixed-witness boundary target to the existence-level
frontier target. The theorems
\lean{ClosureCurrent.MicroscopicClosureLaw.flagshipRecognizableTargetOfRecognizableBoundaryTarget},
\lean{ClosureCurrent.MicroscopicClosureLaw.flagshipRecognizableTargetOfBoundaryTarget},
\lean{ClosureCurrent.MicroscopicClosureLaw.flagshipRecognizableTargetOfExactBoundaryTarget},
\lean{ClosureCurrent.MicroscopicClosureLaw.flagshipRecognizableTargetOfFrontierTarget},
\lean{ClosureCurrent.MicroscopicClosureLaw.normalFormTargetOfRecognizableTarget},
and
\lean{ClosureCurrent.MicroscopicClosureLaw.normalFormTargetSurfaceOfRecognizableTarget}
package the same witness-level recognizable datum into the smaller
existence-level recognizable target
\lean{ClosureCurrent.FlagshipRecognizableTarget}, which already suffices for
the direct normal-form target. The theorems
\lean{ClosureCurrent.MicroscopicClosureLaw.flagshipTargetOfFrontierTarget}
and
\lean{ClosureCurrent.MicroscopicClosureLaw.flagshipSurfaceOfFrontierTarget}
show that one such frontier witness already yields the existence-level
endpoint-witness flagship target and its endpoint-complete surface. Likewise,
\lean{ClosureCurrent.MicroscopicClosureLaw.normalFormTargetOfFrontierTarget}
and
\lean{ClosureCurrent.MicroscopicClosureLaw.normalFormTargetSurfaceOfFrontierTarget}
show that the same frontier witness already yields the direct recognizable
normal-form target and its surface.

This sharper target language is already concretely inhabited by the older
stronger endpoint-complete law. The theorems
\lean{ClosureCurrent.LocalEventFourStateFlagshipLaw.phaseFrontierTarget},
\lean{ClosureCurrent.LocalEventFourStateFlagshipLaw.waveFrontierTarget},
\lean{ClosureCurrent.LocalEventFourStateFlagshipLaw.gaugeFrontierTarget},
and
\lean{ClosureCurrent.LocalEventFourStateFlagshipLaw.geometricFrontierTarget}
read the four smaller frontier targets directly from the law's endpoint
witness bundles, and
\lean{ClosureCurrent.LocalEventFourStateFlagshipLaw.flagshipRecognizableBoundaryTarget}
and
\lean{ClosureCurrent.LocalEventFourStateFlagshipLaw.flagshipRecognizableTarget}
package the resulting recognizable data into the smaller fixed-witness and
existence-level recognizable targets. Also,
\lean{ClosureCurrent.LocalEventFourStateFlagshipLaw.flagshipExactBoundaryTarget}
packages the same stored endpoint witnesses into the exact-side fixed-witness
boundary target. Likewise,
\lean{ClosureCurrent.LocalEventFourStateFlagshipLaw.flagshipBoundaryTarget}
first packages them into the fixed-witness boundary target on the law's
canonical completion witness. Then
\lean{ClosureCurrent.LocalEventFourStateFlagshipLaw.flagshipFrontierTarget}
packages that boundary datum into one witness for
\lean{ClosureCurrent.FlagshipFrontierTarget}. Then
\lean{ClosureCurrent.LocalEventFourStateFlagshipLaw.normalFormTargetOfFrontierTargets}
then recovers the abstract normal-form target from them, while
\lean{ClosureCurrent.LocalEventFourStateFlagshipLaw.normalFormTargetOfRecognizableTargets}
records the same reduction through the smaller recognizable target. The older theorems
\lean{ClosureCurrent.LocalEventFourStateFlagshipLaw.phaseRecognizableTarget},
\lean{ClosureCurrent.LocalEventFourStateFlagshipLaw.waveRecognizableTarget},
\lean{ClosureCurrent.LocalEventFourStateFlagshipLaw.gaugeRecognizableTarget},
and
\lean{ClosureCurrent.LocalEventFourStateFlagshipLaw.geometricRecognizableTarget}
now appear as derived corollaries of that frontier-level reduction.

So the theorem is now sharper about the remaining gap. It does not yet say
that the common no-stage completion law itself forces those four law-native
frontier targets. It says something narrower and rigorous: once those
completion-native frontier targets are supplied on one no-stage completion
witness, they package into one exact fixed-witness flagship boundary target,
and that sharper target already yields the aggregate endpoint-witness target,
the recognizable targets, and hence the flagship identities.
\end{proof}

\begin{proof}[Proof of Corollary~\ref{cor:frontier-target-closes-the-full-detached-capstone}]
\label{proof:cor:frontier-target-closes-the-full-detached-capstone}
This is now a direct composition and is formalized by
\lean{ClosureCurrent.MicroscopicClosureLaw.completionTargetOfFrontierTarget},
\lean{ClosureCurrent.MicroscopicClosureLaw.originClosureTargetOfFrontierTarget},
\lean{ClosureCurrent.MicroscopicClosureLaw.sectorClosureTargetOfFrontierTarget},
\lean{ClosureCurrent.MicroscopicClosureLaw.flagshipExactTargetOfFrontierTarget},
\lean{ClosureCurrent.MicroscopicClosureLaw.frontierTargetFullClosureWitnessSurface},
and
\lean{ClosureCurrent.MicroscopicClosureLaw.frontierTargetFullNormalFormSurface}.
The supplied frontier witness already carries a no-stage completion witness on
the same microscopic law. That same completion witness already forces origin
closure and sector closure there. It also already yields the existence-level
exact endpoint-witness target, so the sharper frontier-target closure and
normal-form capstone theorems then apply directly. So no new law-side datum is
being added here: the existence-level frontier target is itself the single
remaining detached-law seam before both capstone surfaces.

The older stronger endpoint-complete law already inhabits these sharper
capstones. The concrete witnesses are
\lean{ClosureCurrent.LocalEventFourStateFlagshipLaw.fullClosureWitnessSurface}
and
\lean{ClosureCurrent.LocalEventFourStateFlagshipLaw.fullNormalFormSurface}.
\end{proof}

\begin{proof}[Proof of Corollary~\ref{cor:exact-target-closes-the-full-detached-closure-capstone}]
\label{proof:cor:exact-target-closes-the-full-detached-closure-capstone}
This is the exact-side version of the same reduction and is formalized by
\lean{ClosureCurrent.MicroscopicClosureLaw.completionTargetOfExactTarget},
\lean{ClosureCurrent.MicroscopicClosureLaw.originClosureTargetOfExactTarget},
\lean{ClosureCurrent.MicroscopicClosureLaw.sectorClosureTargetOfExactTarget},
\lean{ClosureCurrent.MicroscopicClosureLaw.flagshipRecognizableTargetOfExactTarget},
\lean{ClosureCurrent.MicroscopicClosureLaw.exactTargetFullClosureWitnessSurface},
and
\lean{ClosureCurrent.MicroscopicClosureLaw.exactTargetFullNormalFormSurface}.
The supplied exact witness already carries a no-stage completion witness on
the same microscopic law. That same completion witness already forces origin
closure and sector closure there. It also already yields the smaller
recognizable target. The sharper exact-side closure and
normal-form capstone theorems then apply directly. So for the detached
capstones the exact-side law target already suffices: the stronger frontier
presentation is not needed at this final reduction step.

The older stronger endpoint-complete law realizes that sharper exact-side
input by \lean{ClosureCurrent.LocalEventFourStateFlagshipLaw.flagshipExactTarget},
and its resulting concrete capstones remain
\lean{ClosureCurrent.LocalEventFourStateFlagshipLaw.fullClosureWitnessSurface}
and
\lean{ClosureCurrent.LocalEventFourStateFlagshipLaw.fullNormalFormSurface}.
\end{proof}

\begin{proof}[Proof of Corollary~\ref{cor:recognizable-target-closes-the-full-detached-normal-form-capstone}]
\label{proof:cor:recognizable-target-closes-the-full-detached-normal-form-capstone}
This is the sharper normal-form-only reduction and is formalized by
\lean{ClosureCurrent.MicroscopicClosureLaw.completionTargetOfRecognizableTarget},
\lean{ClosureCurrent.MicroscopicClosureLaw.originClosureTargetOfRecognizableTarget},
\lean{ClosureCurrent.MicroscopicClosureLaw.sectorClosureTargetOfRecognizableTarget},
and
\lean{ClosureCurrent.MicroscopicClosureLaw.recognizableTargetFullNormalFormSurface}.
The supplied recognizable witness already carries a no-stage completion witness
on the same microscopic law. That same completion witness already forces
origin closure and sector closure there, and the sharper recognizable-target
capstone theorem then yields the full detached normal-form surface directly.
So the frontier target is not the minimal remaining input for the normal-form
capstone: the smaller existence-level recognizable target already suffices.

The older stronger endpoint-complete law realizes that sharper input by
\lean{ClosureCurrent.LocalEventFourStateFlagshipLaw.flagshipRecognizableTarget},
and its resulting concrete capstone remains
\lean{ClosureCurrent.LocalEventFourStateFlagshipLaw.fullNormalFormSurface}.
\end{proof}

\begin{proof}[Proof of Theorem~\ref{thm:endpoint-complete-no-stage-law-closes-the-remaining-flagship-seam}]
\label{proof:thm:endpoint-complete-no-stage-law-closes-the-remaining-flagship-seam}
The detached flagship package now carries three linked kinds of datum on one
common selected bundle: the realized same-bundle temporal law on the faithful
cut class, a detached no-stage completion law on that same realized datum, and
the exact phase, wave, gauge, and geometric carrier-face data on that detached
completion law. In Lean this is recorded by
\lean{ClosureCurrent.DetachedFlagshipPackage.toCoherentRealizationNormalForm},
while the grouped closure statement is
\lean{ClosureCurrent.DetachedFlagshipPackage.targetSurface} together
with \lean{ClosureCurrent.DetachedFlagshipPackage.fullClosureWitnessSurface}
and \lean{ClosureCurrent.DetachedFlagshipPackage.fullNormalFormSurface}.

The detached package therefore already supplies both the public selected-cut
surface and its own completion-side normal form on the same datum.

The completion part of the law closes the selected-cut completion surface on
that same datum. The exact boundary identifications then force the four
recognizable flagship identities on the corresponding carrier faces. Hence no
separate frontier package, boundary package, or additional completion datum is
left outside the law statement.

At the same time, the canonical-origin part of the package shows that the selected
anchor, the visible origin, and the coherent selected-cut closure transport are
also internal to that same law-side datum. They are not auxiliary witnesses
supplied from outside. So the detached package closes both seams at once: the origin seam and the
flagship boundary seam.

Therefore the mathematical gain is not a different family of consequences, but
a closed detached law surface. The public temporal realization law yields the
selected-cut realization surface on one common datum, and the detached
flagship package then yields the canonical-origin surface, the completion
surface, the exact flagship surface, and the recognizable flagship identities
on that same datum. The completion-side normal form is simply the detached
reading of those same package-internal data.
\end{proof}

\begin{proof}[Proof of Theorem~\ref{thm:no-stage-completion-law-forces-the-recognizable-flagship-identities}]
\label{proof:thm:no-stage-completion-law-forces-the-recognizable-flagship-identities}
The previous theorems have already isolated the exact remaining flagship seam,
first at the frontier level, then at the sharper endpoint-witness level, and
finally by absorbing those witnesses into the endpoint-complete no-stage law.
The theorem proved here is the direct
application of that boundary statement to one supplied
\lean{ClosureCurrent.FlagshipFrontiers} package.

Formally, this aggregate application is
\lean{ClosureCurrent.LocalEventFourStateCompletionLaw.flagshipSurface}. It
combines the internal completion surface from
\lean{ClosureCurrent.LocalEventFourStateCompletionLaw.attachedSurface} with
the four recognizable witnesses obtained from the supplied phase, wave,
gauge, and geometric frontiers. Hence the detached no-stage completion law
already forces the recognizable Schr\"odinger, wave/Klein--Gordon, Maxwell,
and Einstein identities on its own selected datum once those exact endpoint
frontiers are read there.
\end{proof}

\begin{proof}[Proof of Theorem~\ref{thm:endpoint-complete-no-stage-law-already-carries-the-strongest-current-analytic-flagship-surface}]
\label{proof:thm:endpoint-complete-no-stage-law-already-carries-the-strongest-current-analytic-flagship-surface}
The proposition
\lean{ClosureCurrent.AnalyticFlagshipSurface} is a manuscript-facing package
for the strongest current analytic surface already available from the
detached flagship package. It is not a stronger primitive and it does not add
any extra analytic structure. It merely places the current detached analytic
representative and the recognizable flagship identities into one auditable
statement.

The theorem itself is
\lean{ClosureCurrent.DetachedFlagshipPackage.analyticSurface}. Its
proof is immediate once the flagship package is read through its
completion-side normal form and the previously established package-native
surfaces are placed side by side.

First, the four-state law carried inside the endpoint-complete package still
satisfies
\lean{ClosureCurrent.LocalEventFourStateLaw.minimalQuadraticRepresentativeSurface}.
That theorem gives the compiler-side canonicality of the selected quadratic
law: positivity, equivariance, absolute minimality among PSD competitors, and
equivariant minimality within the selected symmetry class.

Second, the same four-state law still
satisfies
\lean{ClosureCurrent.LocalEventFourStateLaw.stationaryQuadraticRepresentativeSurface}.
That theorem adds the least-motion stationary cut on which the same quadratic
law is actually realized: the selected projection is the finite minimizer and
the realized horizon cut, the selected gradient is stationary there, and the
selected quadratic law is exact-visible and commuting on that same cut.

Third, the same four-state law still
satisfies
\lean{ClosureCurrent.LocalEventFourStateLaw.effectiveQuadraticRepresentativeSurface}.
That theorem sharpens the same selected cut further: the quadratic law already
agrees there with its Schur/Feshbach effective observable law for the
canonical selected shadow propagator, the Schur correction vanishes, and so
does the associated returned residual field.

Fourth, the same four-state law still
satisfies
\lean{ClosureCurrent.LocalEventFourStateLaw.selectedGoverningSurface}. That
theorem packages the next operator identity explicitly: on the same
least-motion cut, the exact effective selected law already coincides
pointwise with the matched selected generator. So the detached analytic lane
already has one common selected governing operator before any further
analytic realization.

Fifth, the same four-state law still
satisfies
\lean{ClosureCurrent.LocalEventFourStateLaw.selectedGoverningFlowSurface}. That
theorem upgrades the common governing operator to a genuine selected
field/flow package in the generic continuum-bridge sense: constant selected
interpolation, exact block decomposition, vanishing off-diagonal/shadow
blocks, and an exact effective flow already identified with the matched
selected generator.

Sixth, the same four-state law still
satisfies
\lean{ClosureCurrent.LocalEventFourStateLaw.selectedFieldLawSurface}. That
theorem compresses the detached selected-cut analytic lane to one selected
operator-residual law: the forced continuum law, the quadratic/effective
scalar residuals, the regular variational datum, the observer-side transport
datum, the Hamilton--Jacobi bridge, and the zero-constitutive promoted law
are all presentations of that same selected law object.

Seventh, the same four-state law still
satisfies
\lean{ClosureCurrent.LocalEventFourStateLaw.selectedFieldLawForcingSurface}.
That theorem sharpens the previous reduction to uniqueness: once the selected
projector and governing residual are fixed, any operator admissible for the
same forced four-state continuum class already agrees with that same selected
law. So the selected operator-residual law is not merely convenient; it is
already uniquely forced on the detached selected cut.

Eighth, the same four-state law still
satisfies
\lean{ClosureCurrent.LocalEventFourStateLaw.selectedPromotionSurface}. That
theorem turns the field/flow package into the strongest current constitutive
surface: the apparent constitutive term is the returned residual of the exact
effective selected law, that term vanishes identically on the exact selected
cut, and the finite promotion chain therefore collapses to one unique branch.

Ninth, the same four-state law still
satisfies
\lean{ClosureCurrent.LocalEventFourStateLaw.effectiveVariationalRepresentativeSurface}.
That theorem identifies the scalarized variational residual itself with the
residual of the exact effective observable law on the same least-motion cut.
So the later variational presentation is already tied to the exact selected
effective law and is not a second postulated layer.

Tenth, the same four-state law still
satisfies
\lean{ClosureCurrent.LocalEventFourStateLaw.variationalRepresentativeSurface}.
That theorem gives admissibility of the selected quadratic form, equality of
the scalar residual read from the quadratic operator and from the matched
selected generator, the regular variational triad, the conservative transport
triad, the Hamilton--Jacobi projection bridge, and the exact identification of
the projection datum with those transported variational data.

The later scalar-law-side unification is now packaged by
\lean{ClosureCurrent.LocalEventFourStateLaw.selectedScalarLawSurface}. That
theorem combines the selected equation boundary, the current scalar PDE/action
realizations, the explicit selected quadratic density, the exact no-hidden-crime
realization ledger, and the exact adjoint-goal ledger into one statement: the
forced selected equation, the two scalar flagship residuals, the density
scalarization, the realized-residual scalarization, and the goal-error
scalarization are already the same scalar law object on the selected cut.

The next law-shaped strengthening is
\lean{ClosureCurrent.LocalEventFourStateLaw.selectedFirstOrderLawSurface}. That
theorem combines the uniquely forced selected operator law, the exact
continuum-bridge
continuous effective flow, and the scalar-law package above into one
first-order surface: every stage of the exact selected flow is already the
same selected operator, and the selected equation together with its current
scalar PDE/action realizations is exactly the zero set of the scalar readout
of that same first-order flow.

The next strengthening after that is
\lean{ClosureCurrent.LocalEventFourStateLaw.selectedEvolutionLawSurface}. That
theorem promotes the exact selected flow, the common selected generator, and
the governing residual to one autonomous evolution law object and proves that
this law is uniquely forced among all autonomous first-order laws compatible
with the same forced continuum class and the same governing residual. Its
solution set is exactly the selected equation and therefore exactly the same
current scalar PDE/action zero set.

The next collapse after that is
\lean{ClosureCurrent.LocalEventFourStateLaw.selectedAutonomousScalarLawSurface}.
That theorem fuses the autonomous evolution law, the canonical scalar law
object, the exact realization ledger, and the exact adjoint-goal ledger into
one uniquely forced autonomous scalar law. Its flow is the same selected
evolution flow, its density is the same selected quadratic density, its
residual is the same scalar residual, and its zero set is the same selected
equation.

The next boundary sharpening after that is
\lean{ClosureCurrent.LocalEventFourStateLaw.selectedFlagshipRealizationSurface}.
That theorem fixes the public contract for any later stronger flagship
realization: once a later language exposes a velocity field, a scalar density,
a scalar residual, and a solution set compatible with the same forced
continuum class and the same selected equation, all four of those public
surfaces are already forced to agree with the autonomous scalar law.

The next sharpening after that is
\lean{ClosureCurrent.LocalEventFourStateLaw.selectedLedgerGoalFlagshipRealizationSurface}.
That theorem tightens the exact host-realization side of the same contract:
any later stronger flagship realization that claims an exact no-hidden-crime
realization chain and an exact adjoint-goal certificate is already forced to
use the same continuum residual, the same terminal realized residual, the same
realized density pairing, the same exact solved state, the same algebraic
residual, and the same exact goal error.

The next sharpening after that is
\lean{ClosureCurrent.LocalEventFourStateLaw.selectedDifferentialFlagshipRealizationSurface}.
That theorem refines the contract from public law data to the exact
continuum-bridge
differential-flow language: any later stronger flagship realization that
claims a continuous block-derivative/effective-flow presentation is already
forced to use the same selected block operator, the same \(PP/PQ/QP/QQ\)
blocks, the same exact effective flow, and the same common zero set.

The next sharpening after that is
\lean{ClosureCurrent.LocalEventFourStateLaw.selectedDerivativeSubsamplingFlagshipSurface}.
That theorem refines the contract again to the continuum-bridge classical
derivative-subsampling language: any later stronger flagship realization that
claims a projection/derivative/observed-restriction presentation is already
forced to use the same constant selected projection family, the same selected
generator as derivative field, the same selected observed restriction, the
same vanishing defect, and the same common zero set.

The next sharpening after that is
\lean{ClosureCurrent.LocalEventFourStateLaw.selectedVariationalFlagshipRealizationSurface}.
That theorem tightens the action side in the Chapter~7 language: any later
stronger flagship realization that claims a regular variational system, an
observer-side transport datum, and a Hamilton--Jacobi projection bridge is
already forced to use the same Lagrangian, Hamiltonian, characteristic,
Eulerian, observer-characteristic, and projected characteristic residuals,
and therefore the same common zero set.

The next sharpening after that is
\lean{ClosureCurrent.LocalEventFourStateLaw.selectedCompositeFlagshipRealizationSurface}.
That theorem is the combined realization-side boundary: once the public,
ledger, differential, derivative, and variational languages are all carried on
one same selected realization, they already collapse jointly to one operator,
one density, one scalar residual, and one zero set.

The next sharpening after that is
\lean{ClosureCurrent.LocalEventFourStateLaw.selectedRealizedFormSurface}. That
theorem rewrites the same combined contract in the strongest differential
language now honestly forced: the common realized form is already first-order
and autonomous, already carried by \(G_{\mathrm{sel}}\), already block-diagonal
with vanishing mixed and hidden blocks, already defect-free on the selected
restriction, and already governed by the selected quadratic density and common
selected equation.

The next sharpening after that is
\lean{ClosureCurrent.LocalEventFourStateLaw.selectedStateObjectSurface}. That
theorem identifies the common carrier of all these readouts: one selected state
object on the explicit four-coordinate carrier \(\mathrm{State}\). The
generator, density, residual, and equation are then different readouts of that
one object rather than parallel primitive data.

The next sharpening after that is
\lean{ClosureCurrent.LocalEventFourStateLaw.selectedStateFieldLawSurface}. That
theorem fuses the selected state object and the autonomous scalar law into one
canonical selected state-field law on the detached selected cut: one selected
projection, one exact autonomous flow, one common generator, one density, one
residual, and one zero set on one same carrier.

The next sharpening after that is
\lean{ClosureCurrent.LocalEventFourStateLaw.selectedStateFieldRealizationSurface}.
That theorem then says the current stronger realization languages are not
parallel formalisms aimed at different hidden laws. The differential-flow,
derivative-subsampling, and variational languages already present that same
selected state-field law exactly.

Thirteenth, the endpoint-complete law already satisfies
\lean{ClosureCurrent.LocalEventFourStateFlagshipLaw.surface}. That theorem
provides the internally reconstructed \lean{PartIVLawCompletionStatement}
together with the recognizable Schr\"odinger, wave/Klein--Gordon, Maxwell,
and Einstein identities on the same intrinsically selected datum.

The new theorem simply concatenates these law-native surfaces. So the gain is
not a new analytic derivation beyond the existing Lean surface. The gain is
that the draft now exposes one final law-native culmination theorem: the
strongest current analytic representative, its least-motion stationary cut,
its already-collapsed exact effective law, its one common selected governing
operator, its one uniquely forced selected operator-residual law, its exact
selected field/flow package, its exact no-hidden-crime realization ledger, its
exact adjoint-goal ledger for that same density, its one canonical scalar law
object, its one canonical first-order law object, its exact constitutive
collapse, the explicit selected quadratic density whose scalarization gives
the same exact effective residual, its one uniquely forced autonomous
evolution law, its one uniquely forced autonomous scalar law, its public
flagship-realization contract, its exact realization-ledger / adjoint-goal
flagship contract, its differential-flow flagship contract, its
derivative-subsampling flagship contract, its variational flagship contract,
its exact combined pre-PDE flagship contract, its exact realized-form boundary
up to differential language, its one underlying selected state object on the
concrete carrier \(\mathrm{State}\), its one canonical selected state-field
law, its state-field realization target surface across the current stronger
languages, and the recognizable textbook equations already coexist on the same
selected datum of the endpoint-complete no-stage law.
\end{proof}

\begin{proof}[Proof of Corollary~\ref{cor:endpoint-complete-no-stage-law-puts-the-flagship-equations-on-one-shared-selected-state}]
\label{proof:cor:endpoint-complete-no-stage-law-puts-the-flagship-equations-on-one-shared-selected-state}
This is the manuscript-facing restatement of
\lean{ClosureCurrent.LocalEventFourStateFlagshipLaw.sharedStateSurface}, whose
packaging proposition is
\lean{ClosureCurrent.FlagshipSharedStateSurface}. The point is not to add new
analytic structure. The point is to isolate one consequence that is easy to
miss when reading only the larger culmination theorem.

By Theorem~\ref{thm:endpoint-complete-no-stage-law-already-carries-the-strongest-current-analytic-flagship-surface},
the endpoint-complete no-stage law already provides the one selected state
object of
\lean{ClosureCurrent.LocalEventFourStateLaw.selectedStateObjectSurface}, the
exact realized-form collapse of
\lean{ClosureCurrent.LocalEventFourStateLaw.selectedRealizedFormSurface}, and
the recognizable Schr\"odinger, wave/Klein--Gordon, Maxwell, and Einstein
identities on one same intrinsically selected datum.

The new corollary simply reads these pieces together. The selected state
object identifies the common carrier: one concrete four-coordinate
\(\mathrm{State}\)-valued selected object with one selected projection, one
selected generator, one density, one residual, and one zero set. The
realized-form theorem then says that every current public realization package
already reads its block operator, derivative field, observed restriction,
density, residual, and solution set from that same object. So within the
present flagship surface the visible textbook equations do not carry separate
primitive state variables; they are all read on one shared selected state.
\end{proof}

\begin{proof}[Proof of Corollary~\ref{cor:detached-selected-cut-already-carries-one-canonical-selected-state-field-law}]
\label{proof:cor:detached-selected-cut-already-carries-one-canonical-selected-state-field-law}
This is the manuscript-facing restatement of
\lean{ClosureCurrent.LocalEventFourStateLaw.selectedStateFieldLawSurface}. The
key point is that two previously separate surfaces have now fused into one
law-shaped object on the detached selected cut.

The first surface is
\lean{ClosureCurrent.LocalEventFourStateLaw.selectedStateObjectSurface}, which
already identifies one selected state object on the concrete carrier
\(\mathrm{State}\). The second is
\lean{ClosureCurrent.LocalEventFourStateLaw.selectedAutonomousScalarLawSurface},
which already identifies one uniquely forced autonomous flow, one common
generator, one density, one residual, and one zero set on that same carrier.

The new theorem packages these together into one selected state-field law:
the selected projection comes from the state-object surface, the exact
autonomous flow comes from the autonomous-law surface, and the generator,
density, residual, and zero set are identified across both. The additional
pointwise agreement with the current public realization packages is supplied by
\lean{ClosureCurrent.LocalEventFourStateLaw.selectedRealizedFormSurface}.

So the detached selected cut now already carries one explicit
\(\mathrm{State}\)-valued law object. This is still weaker than a final
spatial differential PDE or a full action functional, but it makes the
remaining target precise: any stronger realization must present this same
selected state-field law rather than a different underlying law.
\end{proof}

\begin{proof}[Proof of Corollary~\ref{cor:current-stronger-realization-languages-already-present-one-same-selected-state-field-law}]
\label{proof:cor:current-stronger-realization-languages-already-present-one-same-selected-state-field-law}
This is the manuscript-facing restatement of
\lean{ClosureCurrent.LocalEventFourStateLaw.selectedStateFieldRealizationSurface},
whose packaging proposition is
\lean{ClosureCurrent.SelectedStateFieldRealizationSurface}. The theorem does
not add a stronger law. It says the current stronger realization languages are
already exact presentations of the same law object.

By Corollary~\ref{cor:detached-selected-cut-already-carries-one-canonical-selected-state-field-law},
the detached selected cut already carries one canonical selected state-field
law on \(\mathrm{State}\). The differential-flow contract of
\lean{ClosureCurrent.LocalEventFourStateLaw.selectedDifferentialFlagshipRealizationSurface}
already fixes the same operator, the same \(PP/PQ/QP/QQ\) blocks, the same
exact effective flow, and the same common zero set. The derivative-subsampling
contract of
\lean{ClosureCurrent.LocalEventFourStateLaw.selectedDerivativeSubsamplingFlagshipSurface}
already fixes the same selected projection, the same derivative field, the
same observed restriction, the same vanishing defect, and the same common zero
set. And the variational contract of
\lean{ClosureCurrent.LocalEventFourStateLaw.selectedVariationalFlagshipRealizationSurface}
already fixes the same Lagrangian, Hamiltonian, characteristic, Eulerian,
observer-characteristic, and projected-characteristic residuals, hence the
same common zero set.

The new corollary simply reads those already-proved contracts through the
state-field law. The generator in all three languages is the generator of that
same state-field law, the selected projection in the derivative language is
the projection of that same law, the residuals in the variational language are
the residual of that same law, and the solution sets in all three languages
are the same zero set of that same law. So the current stronger realization
languages already present one same selected state-field law rather than
different candidate laws.
\end{proof}

\begin{proof}[Proof of Corollary~\ref{cor:endpoint-complete-no-stage-law-puts-the-flagship-equations-on-one-shared-selected-state-field-law}]
\label{proof:cor:endpoint-complete-no-stage-law-puts-the-flagship-equations-on-one-shared-selected-state-field-law}
This is the manuscript-facing restatement of
\lean{ClosureCurrent.LocalEventFourStateFlagshipLaw.sharedStateFieldSurface},
whose packaging proposition is
\lean{ClosureCurrent.FlagshipSharedStateFieldSurface}. It is the flagship-side
counterpart of the previous two detached selected-cut corollaries.

By Theorem~\ref{thm:endpoint-complete-no-stage-law-already-carries-the-strongest-current-analytic-flagship-surface},
the endpoint-complete no-stage law already provides the recognizable
Schr\"odinger, wave/Klein--Gordon, Maxwell, and Einstein identities on one same
intrinsically selected datum. By
Corollary~\ref{cor:detached-selected-cut-already-carries-one-canonical-selected-state-field-law},
that same selected datum already carries one canonical selected state-field law
on \(\mathrm{State}\). And by
Corollary~\ref{cor:current-stronger-realization-languages-already-present-one-same-selected-state-field-law},
the current stronger realization languages on that selected datum already all
present that same selected state-field law exactly.

So the flagship statement is stronger than mere shared hidden state. Within the
present formal surface, the recognizable textbook equations already live on one
same selected state-field law.
\end{proof}

\begin{proof}[Proof of Corollary~\ref{cor:endpoint-complete-no-stage-law-already-yields-one-unified-field-package}]
\label{proof:cor:endpoint-complete-no-stage-law-already-yields-one-unified-field-package}
This is the manuscript-facing restatement of
\lean{ClosureCurrent.LocalEventFourStateFlagshipLaw.unifiedFieldSurface},
whose packaging proposition is \lean{ClosureCurrent.UnifiedFieldSurface}. The
point is not to add a stronger theorem than the previous corollary, but to
name the actual object now present on the formal surface.

By Corollary~\ref{cor:endpoint-complete-no-stage-law-puts-the-flagship-equations-on-one-shared-selected-state-field-law},
the endpoint-complete no-stage law already puts the recognizable
Schr\"odinger, wave/Klein--Gordon, Maxwell, and Einstein identities on one same
selected datum under one same canonical selected state-field law on
\(\mathrm{State}\), and the current stronger realization languages already all
present that same law exactly. The new Lean structure
\lean{ClosureCurrent.UnifiedField} simply packages those already-proved pieces
into one explicit object:
\begin{enumerate}[label=(\roman*), leftmargin=2em, itemsep=0.2em]
\item the canonical selected state-field law itself;
\item its exact realization-target surface across the current stronger
      languages;
\item the phase, wave, gauge, and geometric recognizable identities on that
      same selected datum.
\end{enumerate}
So the unified-field claim is now explicit at the object level. What is still
absent is not the common law package, but only its final spatial differential
PDE or full action-functional realization.
\end{proof}

\begin{proof}[Proof of Corollary~\ref{cor:endpoint-complete-no-stage-law-already-yields-one-unified-differential-field-package}]
\label{proof:cor:endpoint-complete-no-stage-law-already-yields-one-unified-differential-field-package}
This is the manuscript-facing restatement of
\lean{ClosureCurrent.LocalEventFourStateFlagshipLaw.unifiedDifferentialFieldSurface},
whose packaging proposition is
\lean{ClosureCurrent.UnifiedDifferentialFieldSurface}. The new point is that
the present unified field package already carries the exact differential data
that any final stronger carrier presentation must respect.

By Corollary~\ref{cor:endpoint-complete-no-stage-law-already-yields-one-unified-field-package},
the endpoint-complete no-stage law already yields one explicit unified field
package: one canonical selected state-field law together with its four
recognizable flagship faces. On the detached selected cut, the law also
already carries one exact continuous block-derivative package and one exact
effective-flow package on the same carrier \(\mathrm{State}\), and the current
derivative-subsampling and variational languages already present those same
data exactly; this is the content of
\lean{ClosureCurrent.LocalEventFourStateLaw.selectedDifferentialFlagshipRealizationSurface},
\lean{ClosureCurrent.LocalEventFourStateLaw.selectedDerivativeSubsamplingFlagshipSurface},
and
\lean{ClosureCurrent.LocalEventFourStateLaw.selectedVariationalFlagshipRealizationSurface}.

The new Lean structure \lean{ClosureCurrent.UnifiedDifferentialField} packages
those already-proved pieces into one explicit differential object:
\begin{enumerate}[label=(\roman*), leftmargin=2em, itemsep=0.2em]
\item the unified field package itself;
\item the exact continuous block-derivative datum on \(\mathrm{State}\);
\item the exact effective-flow datum on \(\mathrm{State}\);
\item the exact differential, derivative, and variational realization
      surfaces of that same object.
\end{enumerate}
So the best current rigorous closure is now explicit: the common object is not
only a shared selected state-field law, but a shared state-valued differential
field. What still remains is only the final choice of spatial differential or
action-functional carrier language for that same object.
\end{proof}

\begin{proof}[Proof of Corollary~\ref{cor:endpoint-complete-no-stage-law-leaves-no-new-primitive-carrier-on-the-detached-flagship-surface}]
\label{proof:cor:endpoint-complete-no-stage-law-leaves-no-new-primitive-carrier-on-the-detached-flagship-surface}
This is the manuscript-facing restatement of
\lean{ClosureCurrent.LocalEventFourStateFlagshipLaw.stateCarrierRigiditySurface}.
The claim is not that later work may not choose a different external
presentation space. The claim is that on the current detached flagship surface
there is no remaining primitive carrier freedom.

By Corollary~\ref{cor:endpoint-complete-no-stage-law-already-yields-one-unified-differential-field-package},
the endpoint-complete no-stage law already yields one unified differential
field package on the concrete carrier \(\mathrm{State}\). This is already the
continuum-bridge realization language on that carrier: continuous block derivatives,
exact effective flow, derivative subsampling, and the
variational/observer/Hamilton--Jacobi package.

The second input is
\lean{ClosureCurrent.LocalEventFourStateLaw.selectedCompositeFlagshipRealizationSurface}.
That theorem proves that once the current public, host-ledger, differential,
derivative-subsampling, and variational languages are carried on one same
selected realization, they already collapse pointwise to one selected
generator, one quadratic density, one scalar residual, and one common zero
set. The theorem
\lean{ClosureCurrent.LocalEventFourStateFlagshipLaw.noNewPrimitiveCarrier}
reads off exactly the flagship-facing part of that collapse: the public
velocity, density, residual, and solution surfaces of any admitted stronger
realization already agree pointwise with the same selected law on
\(\mathrm{State}\).

So there is no remaining primitive carrier choice on the detached flagship
surface. What remains is only a presentation problem: write the already-fixed
\(\mathrm{State}\)-carried law in a final spatial differential PDE language or
as a full action functional.
\end{proof}

\begin{proof}[Proof of Corollary~\ref{cor:endpoint-complete-no-stage-law-already-yields-one-unified-transport-field}]
\label{proof:cor:endpoint-complete-no-stage-law-already-yields-one-unified-transport-field}
This is the manuscript-facing restatement of
\lean{ClosureCurrent.LocalEventFourStateFlagshipLaw.unifiedTransportFieldSurface}.
It is the cleanest place where the rebuilt foundation and the continuum bridge
machinery re-enter the detached flagship lane without changing the carrier or
the law.

By Corollary~\ref{cor:endpoint-complete-no-stage-law-already-yields-one-unified-differential-field-package},
the endpoint-complete no-stage law already yields one unified differential
field package on the concrete carrier \(\mathrm{State}\). What is still
missing at that stage is an explicit statement that this same field already
lives inside the active boundary/transport grammar carried by the selected
datum itself.

That extra datum is already present. The selected datum in the detached law is
an explicit \lean{SelectionDatum}, hence it already carries one concrete
\lean{CharacteristicHorizonRealization} and one fixed selected horizon. By
\lean{transfer_part_one} and \lean{characteristic_realization}, that
realization inherits the active Part~I boundary calculus and the transport of
observed boundary and defect across the selected horizon. On the same carrier,
the selected state-field law already carries the exact observer-side
conservative transport triad and the exact Hamilton--Jacobi projection bridge
by \lean{ClosureCurrent.SelectedFieldLaw.observerSurface} and
\lean{ClosureCurrent.SelectedFieldLaw.projectionSurface}.

The new Lean structure \lean{ClosureCurrent.UnifiedTransportField} simply
packages those already-proved pieces together:
\begin{enumerate}[label=(\roman*), leftmargin=2em, itemsep=0.2em]
\item the unified differential field on \(\mathrm{State}\);
\item the selected datum's characteristic horizon realization and selected
      horizon;
\item the observer-side conservative transport datum on \(\mathrm{State}\);
\item the Hamilton--Jacobi projection datum on \(\mathrm{State}\).
\end{enumerate}
So the detached flagship surface is now explicitly attached to the active
boundary/transport grammar on one same carrier. The remaining task is still
only the final spatial differential or action-functional presentation of that
same \(\mathrm{State}\)-valued transport field.
\end{proof}

\begin{proof}[Proof of Corollary~\ref{cor:endpoint-complete-no-stage-law-already-yields-one-unified-boundary-field}]
\label{proof:cor:endpoint-complete-no-stage-law-already-yields-one-unified-boundary-field}
Corollary~\ref{cor:endpoint-complete-no-stage-law-already-yields-one-unified-boundary-field}
uses exactly the next piece of the active Part~I boundary calculus on the same
selected datum.

By Corollary~\ref{cor:endpoint-complete-no-stage-law-already-yields-one-unified-transport-field},
the endpoint-complete no-stage law already yields one unified transport field
on the concrete carrier \(\mathrm{State}\). The selected datum therefore
already carries:
\begin{enumerate}[label=(\roman*), leftmargin=2em, itemsep=0.2em]
\item the realized selected horizon cut;
\item the observed boundary, leakage, return, and defect maps associated to
      that horizon;
\item transport of observed boundary and defect by the same selected
      realization.
\end{enumerate}

The only additional ingredient is the explicit Part~I through-horizon
decomposition \lean{through_horizon_decomposition}. Combined with the already
available identity
\lean{selectedProjection_eq_horizonCut}, it gives the exact split
\[
d u = d_{\mathrm{obs}}u + \ell u
\]
for every selected state \(u\) with \(\Pi_{\mathrm{sel}}u = u\), together with
\[
d_{\mathrm{obs}}^2 u = -\,\delta u.
\]
The equation \(\delta = r \circ \ell\) is just the defining factorization of
\lean{boundaryDefect}, while \(d^2=0\) comes from the active full-boundary
closure law. The transport identities for \(d_{\mathrm{obs}}\) and \(\delta\)
are exactly the two transport clauses already packaged by
\lean{characteristic_realization}.

The new Lean structure \lean{ClosureCurrent.UnifiedBoundaryField} packages
those already-forced pieces together on one same \(\mathrm{State}\)-valued
selected-horizon field. So the present detached flagship surface now carries
its own explicit boundary/leakage/defect grammar on the same carrier, with no
new primitive carrier and no new primitive law.
\end{proof}

\begin{proof}[Proof of Corollary~\ref{cor:endpoint-complete-no-stage-law-already-yields-one-classwise-unified-boundary-field}]
\label{proof:cor:endpoint-complete-no-stage-law-already-yields-one-classwise-unified-boundary-field}
Corollary~\ref{cor:endpoint-complete-no-stage-law-already-yields-one-classwise-unified-boundary-field}
is the classwise canonicity step.

By Corollary~\ref{cor:endpoint-complete-no-stage-law-already-yields-one-unified-boundary-field},
the endpoint-complete no-stage law already yields one explicit selected-horizon
boundary field on the least-motion selected datum. The same selected bridge is
not just one isolated datum, however. It already carries an admissible realized
class, one least-motion root, and classwise horizon-preserving equivalences.

Those ingredients are already part of the canonical selected Schur bridge. The
theorem \lean{canonical_selected_schur_realization} gives, for every datum in
the class, equality of horizons and transport of the selected observed law from
the least-motion root. The remaining transport statements for observed boundary
and defect are then exactly the two clauses already built into
\lean{HorizonPreservingEquivalence}. Rewriting the root datum as the selected
datum gives transport of the observed boundary and defect from the same root
boundary field across the whole class.

The new Lean structure \lean{ClosureCurrent.UnifiedBoundaryClass} simply
packages that content together: one selected-horizon boundary field at the
least-motion root, plus its classwise horizon-preserving transport through the
entire admissible realized class. So the remaining PDE/action step is no
longer even a choice of selected datum within that class.
\end{proof}

\begin{proof}[Proof of Corollary~\ref{cor:endpoint-complete-no-stage-law-already-yields-one-physical-field}]
\label{proof:cor:endpoint-complete-no-stage-law-already-yields-one-physical-field}
Corollary~\ref{cor:endpoint-complete-no-stage-law-already-yields-one-physical-field}
is now immediate from the two previous closures.

By Corollary~\ref{cor:endpoint-complete-no-stage-law-already-yields-one-classwise-unified-boundary-field},
the endpoint-complete no-stage law already determines one canonical
selected-horizon boundary field across the entire admissible realized class on
the concrete carrier \(\mathrm{State}\). By the flagship surface theorem for
the same endpoint-complete law, the recognizable phase, wave, gauge, and
geometric identities already hold on that same selected datum.

The new proposition \lean{ClosureCurrent.PhysicalFieldSurface} does not add new
law data. It records the stronger reading of what has already been proved:
those four recognizable identities are not merely parallel consequences of the
same abstract completion surface, but faces of one same classwise boundary
field. So once the endpoint-complete no-stage law is fixed, no additional
carrier-side state object remains to be chosen for the flagship equations.
\end{proof}

\begin{proof}[Proof of Corollary~\ref{cor:endpoint-complete-no-stage-law-already-yields-one-exact-physical-equation-field}]
\label{proof:cor:endpoint-complete-no-stage-law-already-yields-one-exact-physical-equation-field}
Corollary~\ref{cor:endpoint-complete-no-stage-law-already-yields-one-exact-physical-equation-field}
is the final abstract closure step.

By Corollary~\ref{cor:endpoint-complete-no-stage-law-already-yields-one-classwise-unified-boundary-field},
the endpoint-complete no-stage law already determines one canonical
selected-horizon boundary field across the entire admissible realized class on
the concrete carrier \(\mathrm{State}\). By
Theorem~\ref{thm:the-selected-operator-law-already-has-a-canonical-no-hidden-crime-realization-ledger},
the same selected operator law already has one exact no-hidden-crime
realization chain whose realized residual is exactly the selected operator. By
Theorem~\ref{thm:the-selected-quadratic-density-is-already-an-exact-adjoint-goal-ledger},
the same selected quadratic density is already the exact goal-error ledger of
the canonical adjoint-goal certificate. And by
Corollary~\ref{cor:endpoint-complete-no-stage-law-already-yields-one-physical-field},
the recognizable phase, wave, gauge, and geometric identities already live on
that same selected datum.

The new Lean structure \lean{ClosureCurrent.PhysicalEquationField} does not add
further law data. It simply fuses those already-proved surfaces together and
records the two remaining identifications explicitly: the exact realized
residual is the generator of the same selected state-field law, and the exact
goal error is the density of that same selected state-field law. So the
abstract detached closure is complete before any later external PDE or action
notation is chosen.
\end{proof}

\begin{proof}[Proof of Theorem~\ref{thm:observation-field-detached-current-object-recovers-the-observer-realized-surface}]
\label{proof:thm:observation-field-detached-current-object-recovers-the-observer-realized-surface}
Definition~\ref{def:observation-field-detached-current-object} is formalized by
\lean{ClosureCurrent.VisibleObservationField} and
\lean{ClosureCurrent.ObservationFieldObject}. The new ingredient is a
function-level observation field: one selected-observer family together with a
role-indexed map assigning a visible quotient value to each selected observer.
The constructor
\lean{ClosureCurrent.VisibleObservationField.toVisibleObserverRealization}
derives the earlier predicate-style observer realization from that observation
field. The constructor
\lean{ClosureCurrent.ObservationFieldObject.toObservedCurrentObject} then
recovers the previous observer-realized detached current object, and the
packaged theorem \lean{ClosureCurrent.ObservationFieldObject.surface} gives the
same downstream consequences: nonempty local pair current, direct/return
compatibility, autonomous minimal visible quotients, and nonempty intrinsic
selected-bundle existence.
\end{proof}

\begin{proof}[Proof of Corollary~\ref{cor:observation-field-detached-current-sector-object-closes-the-reduced-seam}]
\label{proof:cor:observation-field-detached-current-sector-object-closes-the-reduced-seam}
The corollary is formalized by
\lean{ClosureCurrent.ObservationFieldSectorObject}. This extends the
function-level observation-field object by the same intrinsic non-dictionary
sector-generation clauses as before. The constructor
\lean{ClosureCurrent.ObservationFieldSectorObject.toObservedCurrentSectorObject}
recovers the previous observer-realized sector object, and the packaged
theorem \lean{ClosureCurrent.ObservationFieldSectorObject.surface} then yields
nonempty intrinsic selected-bundle existence, nonempty intrinsic sector
generation, and the canonical logical profile \([3,2,1]\), together with the
same current-side direct/return and autonomous-quotient facts.
\end{proof}

\begin{proof}[Proof of Theorem~\ref{thm:observer-fiber-detached-current-object-recovers-the-observer-realized-surface}]
\label{proof:thm:observer-fiber-detached-current-object-recovers-the-observer-realized-surface}
Definition~\ref{def:observer-fiber-detached-current-object} is formalized by
\lean{ClosureCurrent.SelectedObserverFiber},
\lean{ClosureCurrent.FiberObservationField}, and
\lean{ClosureCurrent.FiberObservationObject}. The new input is smaller than the
function-level observation-field package: instead of a dependent observation
map on a predicate of selected observers, it uses one selected-observer fiber
and one role-indexed observation map on that fiber. The constructor
\lean{ClosureCurrent.FiberObservationField.toVisibleObserverRealization}
derives the earlier predicate-style observer realization directly from the
fiber data by reading visible states off fiber points. The constructor
\lean{ClosureCurrent.FiberObservationObject.toObservedCurrentObject} then
recovers the observer-realized detached current object, and the packaged
theorem \lean{ClosureCurrent.FiberObservationObject.surface} yields the same
reduced current consequences: nonempty local pair current, direct/return
compatibility, autonomous minimal visible quotients, and nonempty intrinsic
selected-bundle existence.
\end{proof}

\begin{proof}[Proof of Corollary~\ref{cor:observer-fiber-detached-current-sector-object-closes-the-reduced-seam}]
\label{proof:cor:observer-fiber-detached-current-sector-object-closes-the-reduced-seam}
The corollary is formalized by
\lean{ClosureCurrent.FiberObservationSectorObject}. This extends the
observer-fiber detached current object by the same intrinsic non-dictionary
sector-generation clauses used earlier. The constructor
\lean{ClosureCurrent.FiberObservationSectorObject.toObservedCurrentSectorObject}
recovers the observer-realized sector object, and the packaged theorem
\lean{ClosureCurrent.FiberObservationSectorObject.surface} then yields
nonempty intrinsic selected-bundle existence, nonempty intrinsic sector
generation, and the canonical logical profile \([3,2,1]\), together with the
same current-side direct/return and autonomous-quotient facts.
\end{proof}
